% Fonctions associées et étude de fonctions — Mathématiques 1ère A — Cours signé OnBuch+
% Programme officiel MINESEC. Compiler avec : tectonic NA04-etude-fonctions-a.tex
\newcommand{\DOCMATIERE}{Mathématiques}
\newcommand{\DOCNIVEAU}{1ère A}
\newcommand{\DOCMODULE}{Relations et opérations fondamentales dans l'ensemble des nombres réels}
\newcommand{\DOCLECON}{NA04}
\newcommand{\DOCTITRE}{Fonctions associées et étude de fonctions}
% OnBuch+ — préambule « cours signés » ENRICHI (compilé avec tectonic / XeLaTeX)
% Système visuel « Pop » d'OnBuch+ : fond crème (jamais blanc), bordures encre
% épaisses, ombres « dures » décalées, titres Archivo Black en capitales.
% Version MATHÉMATIQUES : graphes (pgfplots/tikz), boîtes pédagogiques
% (définition, propriété, méthode, attention, le savais-tu, exercice résolu,
% exercice, TP, bilan), page de garde et signature OnBuch+ ancrée en bas.
%
% Macros attendues dans main.tex AVANT \input{preamble} :
%   \DOCMATIERE  \DOCNIVEAU  \DOCMODULE  \DOCLECON (numéro)  \DOCTITRE
\documentclass[11pt]{article}

\usepackage{fontspec}
\usepackage[french]{babel}
\usepackage[a4paper,margin=2cm,top=3.4cm,bottom=2.8cm,headheight=1.4cm,footskip=1.3cm]{geometry}
\usepackage{amsmath,amssymb,amsfonts}
\usepackage{mathtools} % \xmapsto, \coloneqq... souvent utilisés par les modèles
\usepackage{cancel} % simplifications barrées
% \abs{z} (module, valeur absolue) et \norm{u} (norme), utilisés par les modèles
\providecommand{\abs}{}\let\abs\relax\DeclarePairedDelimiter{\abs}{\lvert}{\rvert}
\providecommand{\norm}{}\let\norm\relax\DeclarePairedDelimiter{\norm}{\lVert}{\rVert}
\usepackage[table]{xcolor} % [table] : fournit aussi \rowcolors (alternance de lignes)
\usepackage{pagecolor}
\usepackage{tikz}
\usetikzlibrary{arrows.meta,positioning,calc,shapes.geometric,shapes.misc,shapes.arrows,decorations.pathmorphing,decorations.pathreplacing,decorations.markings,patterns,shadows,mindmap,trees,fit,backgrounds,matrix,intersections,angles,quotes}
\usepackage{pgfplots}
\usepackage{pgf-pie} % diagrammes circulaires \pie (statistiques)
\pgfplotsset{compat=1.18}
\usepgfplotslibrary{fillbetween,groupplots} % aires entre courbes (calcul intégral)
\usepackage{enumitem}
\usepackage{fancyhdr}
% [most] charge toutes les bibliothèques tcolorbox (listings, minted, raster,
% documentation...) et gonfle inutilement la mémoire TeX sur un document long
% (~300 boîtes) : on ne charge que ce qui est réellement utilisé.
\usepackage[skins,breakable]{tcolorbox}
\usepackage{graphicx}
\usepackage{colortbl}
\usepackage{booktabs}
\usepackage{tabularx}
\usepackage{diagbox} % case d'en-tête diagonale des tableaux à double entrée
% Colonnes extensibles centrée / gauche / droite (C, L, R), souvent utilisées par les modèles
\newcolumntype{C}{>{\centering\arraybackslash}X}
\newcolumntype{L}{>{\raggedright\arraybackslash}X}
\newcolumntype{R}{>{\raggedleft\arraybackslash}X}
\usepackage{array}
\usepackage{multicol}
\usepackage{siunitx}
% parse-numbers=false : accepte aussi \SI{2,5 \times 10^{-3}}{...}
\sisetup{locale=FR,output-decimal-marker={,},group-minimum-digits=4,parse-numbers=false}
\DeclareSIUnit\ohm{\ensuremath{\symup{\Omega}}} % U+2126 absent de la police
\usepackage{qrcode}
% \degree : commande fréquemment utilisée par les modèles pour un angle ou une
% température en dehors de \SI{}{} (non fournie par les paquets chargés).
\providecommand{\degree}{^{\circ}}
% \qed (fin de démonstration), utilisé par les modèles sans amsthm
\providecommand{\qed}{\hfill$\square$}
% \barre : double barre des valeurs interdites dans un tableau de variations (tkz-tab)
\providecommand{\barre}{\ensuremath{\|}}
% environnement proof (amsthm non chargé) utilisé par les modèles
\ifdefined\proof\else\newenvironment{proof}[1][Démonstration]{\par\noindent\textit{#1.}\ }{\qed\par}\fi
\usepackage{lastpage}
\usepackage{hyperref}
\hypersetup{hidelinks}

% ============================================================
% Palette « Pop » OnBuch+ — mêmes hex que la classe P (plus_ui.dart)
% ============================================================
\definecolor{poporange}{HTML}{F59321}
\definecolor{popdark}{HTML}{E07A0C}
\definecolor{popcream}{HTML}{FFF6E8}
\definecolor{popink}{HTML}{1C1714}
\definecolor{poppurple}{HTML}{7A5AE0}
\definecolor{popgreen}{HTML}{2E9E6A}
\definecolor{poppink}{HTML}{E0457B}
\definecolor{popblue}{HTML}{2F7DE1}
\definecolor{popgold}{HTML}{C98A12}
\definecolor{popmuted}{HTML}{8A7F76}
\definecolor{popline}{HTML}{EADFD2}
\definecolor{popgray}{HTML}{8A7F76}
\colorlet{popgrey}{popgray}
% Alias tolérants (le modèle nomme parfois les couleurs différemment)
\colorlet{popwhite}{white}
\colorlet{popblack}{popink}
\colorlet{popred}{poppink}
\colorlet{popyellow}{popgold}
% Teintes claires pour fonds de tableaux / zones
\colorlet{popblueL}{popblue!10!white}
\colorlet{poporangeL}{poporange!14!white}
\colorlet{popgreenL}{popgreen!12!white}
\colorlet{poppurpleL}{poppurple!12!white}
\colorlet{poppinkL}{poppink!12!white}
\colorlet{popgoldL}{popgold!14!white}

\pagecolor{popcream}

% ============================================================
% Typographie
% ============================================================
\defaultfontfeatures{Ligatures=TeX}
\setmainfont{PlusJakartaSans-Medium.ttf}[Path=../../../fonts/,BoldFont=PlusJakartaSans-Bold.ttf]
% Maths en sans-serif (Fira Math) avec chiffres et lettres droites de Plus
% Jakarta, pour que les formules se fondent dans le texte.
\usepackage[math-style=ISO,bold-style=ISO]{unicode-math}
\setmathfont{FiraMath-Regular.otf}[Path=../../../fonts/]
\setmathfont{PlusJakartaSans-Medium.ttf}[Path=../../../fonts/,range={up/num,up/{Latin,latin},\mathup}]
\usepackage{icomma} % virgule décimale sans espace en mode math
% Caractères mathématiques Unicode absents des polices de texte (Plus Jakarta,
% Archivo Black) : rendus par leur équivalent mathématique, en texte comme en
% formule — sinon ils s'affichent en carré vide (ex. « Arithmétique dans ℤ »).
\usepackage{newunicodechar}
\newunicodechar{ℝ}{\ensuremath{\mathbb{R}}}
\newunicodechar{ℤ}{\ensuremath{\mathbb{Z}}}
\newunicodechar{ℂ}{\ensuremath{\mathbb{C}}}
\newunicodechar{ℕ}{\ensuremath{\mathbb{N}}}
\newunicodechar{ℚ}{\ensuremath{\mathbb{Q}}}
\newunicodechar{𝔻}{\ensuremath{\mathbb{D}}}
\newunicodechar{α}{\ensuremath{\alpha}}
\newunicodechar{β}{\ensuremath{\beta}}
\newunicodechar{γ}{\ensuremath{\gamma}}
\newunicodechar{δ}{\ensuremath{\delta}}
\newunicodechar{ε}{\ensuremath{\varepsilon}}
\newunicodechar{θ}{\ensuremath{\theta}}
\newunicodechar{λ}{\ensuremath{\lambda}}
\newunicodechar{μ}{\ensuremath{\mu}}
\newunicodechar{σ}{\ensuremath{\sigma}}
\newunicodechar{τ}{\ensuremath{\tau}}
\newunicodechar{φ}{\ensuremath{\varphi}}
\newunicodechar{ψ}{\ensuremath{\psi}}
\newunicodechar{ω}{\ensuremath{\omega}}
\newunicodechar{Σ}{\ensuremath{\Sigma}}
% majuscules produites par \MakeUppercase dans les titres (α-aminés -> Α)
\newunicodechar{Α}{\ensuremath{\alpha}}
\newunicodechar{Β}{\ensuremath{\beta}}
\newunicodechar{Γ}{\ensuremath{\Gamma}}
\newunicodechar{Δ}{\ensuremath{\Delta}}
\newunicodechar{Ε}{\ensuremath{\varepsilon}}
\newunicodechar{Θ}{\ensuremath{\Theta}}
\newunicodechar{Λ}{\ensuremath{\Lambda}}
\newunicodechar{Μ}{\ensuremath{\mu}}
\newunicodechar{Τ}{\ensuremath{\tau}}
\newunicodechar{Φ}{\ensuremath{\Phi}}
\newunicodechar{Ψ}{\ensuremath{\Psi}}
\newunicodechar{Ω}{\ensuremath{\Omega}}
\newunicodechar{↦}{\ensuremath{\mapsto}}
\newunicodechar{∈}{\ensuremath{\in}}
\newunicodechar{∉}{\ensuremath{\notin}}
\newunicodechar{∧}{\ensuremath{\wedge}}
\newunicodechar{⇔}{\ensuremath{\Leftrightarrow}}
\newunicodechar{⟺}{\ensuremath{\Longleftrightarrow}}
\newunicodechar{⇒}{\ensuremath{\Rightarrow}}
\newunicodechar{⟹}{\ensuremath{\Longrightarrow}}
\newunicodechar{∘}{\ensuremath{\circ}}
\newunicodechar{∪}{\ensuremath{\cup}}
\newunicodechar{∩}{\ensuremath{\cap}}
\newunicodechar{≡}{\ensuremath{\equiv}}
\newunicodechar{⊂}{\ensuremath{\subset}}
\newunicodechar{⊥}{\ensuremath{\perp}}
\newunicodechar{∃}{\ensuremath{\exists}}
\newunicodechar{∀}{\ensuremath{\forall}}
\newunicodechar{∛}{\ensuremath{\sqrt[3]{\,}}}
\newunicodechar{∜}{\ensuremath{\sqrt[4]{\,}}}
\newunicodechar{⃗}{\ensuremath{{}^{\to}}}
\newunicodechar{ⁿ}{\ifmmode^{n}\else\textsuperscript{\ensuremath{n}}\fi}
\newunicodechar{ˣ}{\ifmmode^{x}\else\textsuperscript{\ensuremath{x}}\fi}
\newunicodechar{ᵇ}{\ifmmode^{b}\else\textsuperscript{\ensuremath{b}}\fi}
\newunicodechar{ʳ}{\ifmmode^{r}\else\textsuperscript{\ensuremath{r}}\fi}
\newunicodechar{ᵘ}{\ifmmode^{u}\else\textsuperscript{\ensuremath{u}}\fi}
\newunicodechar{ʸ}{\ifmmode^{y}\else\textsuperscript{\ensuremath{y}}\fi}
\newunicodechar{ᵃ}{\ifmmode^{a}\else\textsuperscript{\ensuremath{a}}\fi}
\newunicodechar{ᵉ}{\ifmmode^{e}\else\textsuperscript{\ensuremath{e}}\fi}
\newunicodechar{ᵏ}{\ifmmode^{k}\else\textsuperscript{\ensuremath{k}}\fi}
\newunicodechar{ᵖ}{\ifmmode^{p}\else\textsuperscript{\ensuremath{p}}\fi}
\newunicodechar{ᴺ}{\ifmmode^{N}\else\textsuperscript{\ensuremath{N}}\fi}
\newunicodechar{ᶜ}{\ifmmode^{c}\else\textsuperscript{\ensuremath{c}}\fi}
\newunicodechar{ʰ}{\ifmmode^{h}\else\textsuperscript{\ensuremath{h}}\fi}
\newunicodechar{ᵠ}{\ifmmode^{q}\else\textsuperscript{\ensuremath{q}}\fi}
\newunicodechar{ₙ}{\ifmmode_{n}\else\textsubscript{\ensuremath{n}}\fi}
\newunicodechar{ₐ}{\ifmmode_{a}\else\textsubscript{\ensuremath{a}}\fi}
\newunicodechar{ᵢ}{\ifmmode_{i}\else\textsubscript{\ensuremath{i}}\fi}
\newunicodechar{ᵣ}{\ifmmode_{r}\else\textsubscript{\ensuremath{r}}\fi}
\newunicodechar{ₖ}{\ifmmode_{k}\else\textsubscript{\ensuremath{k}}\fi}
\newunicodechar{ᵦ}{\ifmmode_{\beta}\else\textsubscript{\ensuremath{\beta}}\fi}
\newunicodechar{ₓ}{\ifmmode_{x}\else\textsubscript{\ensuremath{x}}\fi}
\newfontfamily\popdisplay{ArchivoBlack-Regular.ttf}[Path=../../../fonts/]
\newfontfamily\poplogo{SpaceGrotesk-Bold.ttf}[Path=../../../fonts/]
\newfontfamily\popsemi{PlusJakartaSans-SemiBold.ttf}[Path=../../../fonts/]
\newfontfamily\popxbold{PlusJakartaSans-ExtraBold.ttf}[Path=../../../fonts/]
\newfontfamily\popxboldls{PlusJakartaSans-ExtraBold.ttf}[Path=../../../fonts/,LetterSpace=3.0]

\setlist[enumerate]{leftmargin=1.7em,itemsep=5pt,topsep=4pt}
\setlist[itemize]{leftmargin=1.7em,itemsep=4pt,topsep=4pt,label={\color{poporange}\textbullet}}
\setlength{\parindent}{0pt}
\setlength{\parskip}{5pt}
\renewcommand{\arraystretch}{1.25}

% pgfplots : style Pop par défaut
\pgfplotsset{
  every axis/.append style={axis line style={popink,line width=1pt},tick style={popink},
    grid style={popline,line width=0.5pt},label style={font=\small},tick label style={font=\footnotesize}},
  popcurve/.style={poporange,line width=1.8pt,no markers,smooth},
}
% Même style disponible dans un \draw TikZ simple (hors pgfplots)
\tikzset{popcurve/.style={poporange,line width=1.8pt}}

% ============================================================
% Pill / squiggle
% ============================================================
\newcommand{\poppill}[3]{%
  \tikz[baseline=-0.55ex]\node[draw=popink,line width=1.2pt,rounded corners=2.6pt,%
    fill=#2,inner xsep=6.5pt,inner ysep=3pt]{%
    \popxboldls\fontsize{7}{7}\selectfont\color{#3}\MakeUppercase{#1}};%
}
\newcommand{\popsquiggle}[1]{%
  \tikz[baseline=-0.4ex]{
    \draw[#1,line width=2.6pt,line cap=round]
      plot [smooth] coordinates {(0,0.09) (0.34,0.01) (0.68,0.21) (1.02,0.01) (1.36,0.10)};
  }%
}
% Mot-clé surligné (marqueur orange sous le texte)
\newcommand{\cle}[1]{{\popxbold\color{popdark}#1}}

% ============================================================
% En-tête / pied
% ============================================================
\pagestyle{fancy}
\fancyhf{}
\renewcommand{\headrulewidth}{0pt}
\renewcommand{\footrulewidth}{0pt}
\lhead{\raisebox{-0.15ex}{\poplogo\fontsize{13}{13}\selectfont\color{popink}OnBuch\color{poporange}+}}
\rhead{\poppill{\DOCMATIERE\ \textbullet\ \DOCNIVEAU\ \textbullet\ Leçon \DOCLECON}{white}{popink}}
\lfoot{\popxbold\fontsize{7.6}{7.6}\selectfont\color{popmuted}\MakeUppercase{Cours signé OnBuch+}\ \textbullet\ {\popsemi\DOCTITRE}}
\rfoot{\tikz[baseline=-0.6ex]\node[rounded corners=8.5pt,fill=popink,minimum height=17pt,minimum width=17pt,inner xsep=5pt,inner sep=0]{\popxbold\fontsize{8}{8}\selectfont\color{popcream}\thepage\,/\,\pageref*{LastPage}};}

% ============================================================
% Titres
% ============================================================
\newcommand{\chaptitre}[2]{%
  \begin{tcolorbox}[enhanced,colback=poporange,colframe=popink,boxrule=2.6pt,arc=11pt,
      drop shadow=popink,left=15pt,right=15pt,top=12pt,bottom=11pt,before skip=3pt,after skip=18pt]
    \poppill{#2}{popink}{popcream}\par\vspace{9pt}
    {\raggedright\hyphenpenalty=10000\exhyphenpenalty=10000\popdisplay\fontsize{22}{24}\selectfont\color{popcream}\MakeUppercase{#1}\par}
  \end{tcolorbox}
}
% Section numérotée : pastille orange avec le numéro + titre Archivo
\newcounter{popsec}
\newcommand{\coursec}[1]{%
  \stepcounter{popsec}\setcounter{popsub}{0}%
  \par\vspace{16pt}\noindent
  \begin{minipage}{\linewidth}
  \tikz[baseline=-0.7ex]\node[draw=popink,line width=1.6pt,fill=poporange,rounded corners=4pt,minimum size=22pt,inner sep=0]{\popdisplay\fontsize{12}{12}\selectfont\color{popcream}\thepopsec};%
  \hspace{8pt}{\popdisplay\fontsize{15}{17}\selectfont\color{popink}\MakeUppercase{#1}}\par
  \vspace{4pt}\noindent{\color{poporange}\rule{36pt}{3.6pt}}
  \end{minipage}\par\nopagebreak\vspace{8pt}
}
\newcounter{popsub}
\newcommand{\courssub}[1]{%
  \stepcounter{popsub}%
  \par\vspace{9pt}\noindent{\popxbold\fontsize{11.5}{13}\selectfont\color{popink}{\color{poporange}\thepopsec.\thepopsub}\hspace{6pt}#1}\par\nopagebreak\vspace{4pt}
}

% ============================================================
% Boîtes pédagogiques — même recette : fond blanc, bordure épaisse colorée,
% ombre dure encre, badge plein. Titre optionnel [#1].
% ============================================================
\tcbset{popbase/.style={breakable,enhanced,colback=white,boxrule=2.3pt,arc=9pt,
  shadow={3pt}{-3pt}{0pt}{popink},left=14pt,right=14pt,top=11pt,bottom=11pt,before skip=11pt,after skip=15pt,
  fontupper=\popsemi\fontsize{10.6}{14.5}\selectfont\color{popink},
  fontlower=\popsemi\fontsize{10.6}{14.5}\selectfont\color{popink}}}
% Coche dessinée (le glyphe ✓ n'existe pas dans les polices du document)
\newcommand{\popcheck}[1]{\tikz[baseline=-0.5ex]\draw[#1,line width=1.6pt,line cap=round,line join=round](-0.13,0.0)--(-0.04,-0.09)--(0.13,0.1);}
\newcommand{\popboxhead}[3]{\poppill{#1}{#2}{white}\ifx\relax#3\relax\else\hspace{6pt}{\popxbold\fontsize{10.6}{12}\selectfont\color{popink}#3}\fi\par\vspace{7pt}}

\newtcolorbox{aretenir}[1][]{popbase,colframe=popgreen,before upper={\popboxhead{À retenir}{popgreen}{#1}}}
\newtcolorbox{exemplebox}[1][]{popbase,colframe=poporange,before upper={\popboxhead{Exemple}{poporange}{#1}}}
\newtcolorbox{definition}[1][]{popbase,colframe=popblue,before upper={\popboxhead{Définition}{popblue}{#1}}}
\newtcolorbox{propriete}[1][]{popbase,colframe=poppurple,before upper={\popboxhead{Propriété}{poppurple}{#1}}}
\newtcolorbox{methode}[1][]{popbase,colframe=popgold,before upper={\popboxhead{Méthode}{popgold}{#1}}}
\newtcolorbox{attention}[1][]{popbase,colframe=poppink,before upper={\popboxhead{Attention}{poppink}{#1}}}
\newtcolorbox{experience}[1][]{popbase,colframe=popblue,colback=popblueL,before upper={\popboxhead{Expérience}{popblue}{#1}}}
% « Le savais-tu ? » : carte violette pleine (contexte, culture, Cameroun)
\newtcolorbox{savaistu}[1][]{popbase,colback=poppurple,colframe=popink,
  fontupper=\popsemi\fontsize{10.4}{14.2}\selectfont\color{white},
  before upper={\poppill{Le savais-tu\,?}{white}{poppurple}\ifx\relax#1\relax\else\hspace{6pt}{\popxbold\color{white}#1}\fi\par\vspace{7pt}}}
% Exercice résolu : énoncé en haut, \tcblower puis la solution sur fond crème
\newcounter{popexo}\newcounter{popexr}
\newtcolorbox{exoresolu}[1][]{popbase,colframe=poporange,colbacklower=popcream,
  segmentation style={popink,line width=1.2pt,dashed},
  before upper={\stepcounter{popexr}\popboxhead{Exercice résolu \thepopexr}{poporange}{#1}},
  before lower={\poppill{Solution}{popink}{popcream}\par\vspace{6pt}}}
% Exercice (non résolu en place) — la correction est dans la partie Corrigés
\newtcolorbox{exercice}[1][]{popbase,colframe=popink,
  before upper={\stepcounter{popexo}\popboxhead{Exercice \thepopexo}{popink}{#1}}}
% Corrigé
\newtcolorbox{corrige}[1][]{popbase,colframe=popgreen,colback=popgreenL,
  before upper={\popboxhead{Corrigé}{popgreen}{#1}}}
% Objectifs / prérequis / bilan
\newtcolorbox{objectifs}{popbase,colframe=popink,colback=poporangeL,
  before upper={\popboxhead{Objectifs de la leçon}{popink}{}}}
\newtcolorbox{prerequis}{popbase,colframe=popink,colback=popblueL,
  before upper={\popboxhead{Prérequis}{popblue}{}}}
\newtcolorbox{bilan}{popbase,colframe=popink,colback=white,boxrule=2.8pt,
  before upper={\popboxhead{Fiche bilan}{poporange}{L'essentiel à connaître pour l'examen}}}
% Figure encadrée avec légende
\newtcolorbox{popfigure}[1][]{enhanced,breakable=false,colback=white,colframe=popline,boxrule=1.4pt,arc=8pt,
  left=10pt,right=10pt,top=10pt,bottom=8pt,before skip=10pt,after skip=12pt,halign=center,
  lower separated=false,halign lower=center,
  fontlower=\popsemi\fontsize{9}{11.5}\selectfont\color{popmuted},
  #1}
\newcommand{\legende}[1]{\par\vspace{6pt}{\popsemi\fontsize{9}{11.5}\selectfont\color{popmuted}#1}}

% ============================================================
% Signature OnBuch+
% ============================================================
\newcommand{\poplogoinline}[1]{{\poplogo\fontsize{#1}{#1}\selectfont\color{popink}OnBuch\color{poporange}+}}
\newcommand{\signatureonbuch}{%
  \begin{tcolorbox}[enhanced,colback=popink,colframe=popink,boxrule=2.2pt,arc=10pt,
      drop shadow=poporange,left=14pt,right=14pt,top=10pt,bottom=10pt,before skip=0pt,after skip=0pt]
    \tikz[baseline=-0.7ex]\node[circle,fill=popgreen,draw=popcream,line width=1.3pt,minimum size=22pt,inner sep=0]{\popcheck{white}};%
    \hspace{9pt}\begin{minipage}[c]{0.62\linewidth}
      {\popxbold\fontsize{12}{13}\selectfont\color{popcream}Signé\ \textbullet\ {\poplogo OnBuch\color{poporange}+}}\par\vspace{2pt}
      {\raggedright\popsemi\fontsize{8}{10.5}\selectfont\color{popline}\MakeUppercase{Équipe pédagogique OnBuch+}\\\MakeUppercase{Conforme au programme officiel MINESEC}\par}
    \end{minipage}\hfill
    {\poplogo\fontsize{22}{22}\selectfont\color{popcream}OnBuch\color{poporange}+}
  \end{tcolorbox}%
}

% ============================================================
% Page de garde
% #1 titre · #2 matière · #3 niveau · #4 module · #5 n° leçon · #6 durée ·
% #7 description / objectifs (texte ou itemize)
% ============================================================
\newcommand{\pagegarde}[7]{%
  \thispagestyle{empty}
  \newgeometry{a4paper,margin=1.7cm,top=1.6cm,bottom=1.4cm}%
  \begin{tikzpicture}[remember picture,overlay]
    \fill[poporange!18] ([xshift=-3.2cm,yshift=-3.6cm]current page.north east) circle (4.6cm);
    \fill[poppurple!14] ([xshift=2.4cm,yshift=7.6cm]current page.south west) circle (3.8cm);
  \end{tikzpicture}%
  {\poplogo\fontsize{30}{30}\selectfont\color{popink}OnBuch\color{poporange}+}\hfill
  \raisebox{0.6ex}{\poppill{Cours signé \textbullet\ Premium}{popink}{popcream}}\par
  \vspace{1pt}\popsquiggle{poppurple}\par
  \vspace{8pt}
  \begin{tcolorbox}[enhanced,colback=poporange,colframe=popink,boxrule=2.8pt,arc=13pt,
      drop shadow=popink,left=18pt,right=18pt,top=12pt,bottom=13pt,after skip=10pt]
    \poppill{#2 \textbullet\ #3}{popink}{popcream}\hspace{5pt}\poppill{#4}{white}{popink}\par\vspace{8pt}
    {\popdisplay\fontsize{14}{14}\selectfont\color{popink}LEÇON #5}\par\vspace{4pt}
    {\raggedright\hyphenpenalty=10000\exhyphenpenalty=10000\popdisplay\fontsize{25}{27}\selectfont\color{popcream}\MakeUppercase{#1}\par}
    \vspace{8pt}
    {\popxbold\fontsize{9.5}{9.5}\selectfont\color{popink}\MakeUppercase{Durée conseillée : #6}}
  \end{tcolorbox}
  \begin{tcolorbox}[enhanced,colback=white,colframe=popink,boxrule=2.2pt,arc=10pt,
      drop shadow=popink,left=16pt,right=16pt,top=10pt,bottom=10pt,after skip=8pt]
    \poppill{Dans cette leçon}{poppurple}{white}\par\vspace{5pt}
    \popsemi\fontsize{9.3}{12.6}\selectfont\color{popink}#7
  \end{tcolorbox}
  \noindent\begin{minipage}[t]{0.487\linewidth}
    \begin{tcolorbox}[enhanced,colback=popgreen,colframe=popink,boxrule=2.2pt,arc=10pt,
        drop shadow=popink,left=11pt,right=11pt,top=10pt,bottom=10pt,height=3.1cm,valign=center]
      \tcbox[colback=white,colframe=popink,boxrule=1.3pt,arc=5pt,boxsep=0pt,nobeforeafter,
        left=3pt,right=3pt,top=3pt,bottom=3pt,valign=center]{\qrcode[height=1.9cm]{https://play.google.com/store/apps/details?id=com.luvvix.onbuch&hl=fr}}%
      \hspace{8pt}\begin{minipage}[c]{0.5\linewidth}
        {\popdisplay\fontsize{11}{12.5}\selectfont\color{white}\MakeUppercase{Télécharge OnBuch}}\par\vspace{4pt}
        {\raggedright\popsemi\fontsize{8.4}{11}\selectfont\color{white}Gratuit \textbullet\ Google Play\\Tuteur IA, annales, concours, résultats\par}
      \end{minipage}
    \end{tcolorbox}
  \end{minipage}\hfill
  \begin{minipage}[t]{0.487\linewidth}
    \begin{tcolorbox}[enhanced,colback=poppurple,colframe=popink,boxrule=2.2pt,arc=10pt,
        drop shadow=popink,left=11pt,right=11pt,top=10pt,bottom=10pt,height=3.1cm,valign=center]
      \tikz[baseline=-0.7ex]\node[circle,fill=white,draw=popink,line width=1.3pt,minimum size=1.6cm,inner sep=0]{\popdisplay\fontsize{24}{24}\selectfont\color{poppurple}+};%
      \hspace{8pt}\begin{minipage}[c]{0.6\linewidth}
        {\popdisplay\fontsize{11}{12.5}\selectfont\color{white}\MakeUppercase{Passe à OnBuch+}}\par\vspace{4pt}
        {\raggedright\popsemi\fontsize{8.4}{11}\selectfont\color{white}Tous les cours signés, Tuteur IA illimité, fiches PDF — onglet OnBuch+ de l'app\par}
      \end{minipage}
    \end{tcolorbox}
  \end{minipage}
  \vspace{8pt}
  \popcontenu
  \vfill
  \signatureonbuch
  \restoregeometry
  \newpage
}

% Rangée « Ce cours contient » (page de garde)
\newcommand{\poptile}[3]{% #1 couleur #2 gros texte #3 légende
  \begin{tcolorbox}[enhanced,colback=#1,colframe=popink,boxrule=2pt,arc=9pt,shadow={2.5pt}{-2.5pt}{0pt}{popink},
      left=6pt,right=6pt,top=9pt,bottom=9pt,halign=center,valign=center,height=1.75cm,width=\linewidth,nobeforeafter]
    {\popdisplay\fontsize{12.5}{13}\selectfont\color{white}\MakeUppercase{#2}}\par\vspace{3pt}
    {\centering\popsemi\fontsize{7.8}{9.5}\selectfont\color{white}#3\par}
  \end{tcolorbox}}
\newcommand{\popcontenu}{%
  \par\noindent{\popxboldls\fontsize{7.5}{7.5}\selectfont\color{popmuted}CE COURS SIGNÉ CONTIENT}\par\vspace{7pt}
  \noindent\begin{minipage}[t]{0.235\linewidth}\poptile{popblue}{Cours}{complet, illustré, pas à pas}\end{minipage}\hfill
  \begin{minipage}[t]{0.235\linewidth}\poptile{poporange}{Méthodes}{et exercices résolus}\end{minipage}\hfill
  \begin{minipage}[t]{0.235\linewidth}\poptile{poppink}{Exercices}{type examen + corrigés}\end{minipage}\hfill
  \begin{minipage}[t]{0.235\linewidth}\poptile{popgreen}{Fiche bilan}{l'essentiel à retenir}\end{minipage}\par}

% Sommaire
\newtcolorbox{sommaire}{enhanced,colback=white,colframe=popink,boxrule=2.2pt,arc=10pt,
  drop shadow=popink,left=18pt,right=18pt,top=13pt,bottom=13pt,before skip=6pt,after skip=20pt,
  before upper={\poppill{Sommaire}{poppurple}{white}\par\vspace{9pt}\popsemi\fontsize{10.6}{14.5}\selectfont\color{popink}}}

% Signature de fin de cours — poussée tout en bas de la dernière page
\newcommand{\findecours}{%
  \par\vfill\nopagebreak
  \begin{center}\popsquiggle{poporange}\end{center}\vspace{4pt}
  \signatureonbuch
}
\AtBeginDocument{\def\llbracket{\mathopen{[\mkern-3.5mu[}}\def\rrbracket{\mathclose{]\mkern-3.5mu]}}} % intervalles d'entiers (glyphe ⟦ absent de la police)
% Glyphes absents de Fira Math : redéfinis à partir de glyphes disponibles
\AtBeginDocument{%
  \def\setminus{\mathbin{\backslash}}\let\smallsetminus\setminus
  \def\vdots{\mathord{\vbox{\baselineskip3.2pt\lineskiplimit0pt\kern4pt\hbox{.}\hbox{.}\hbox{.}}}}%
  \def\ddots{\mathinner{\mkern1mu\raise6.5pt\hbox{.}\mkern2mu\raise3.6pt\hbox{.}\mkern2mu\raise0.7pt\hbox{.}\mkern1mu}}%
  \def\triangle{\mathord{\scalebox{0.9}{$\symup{\Delta}$}}}%
  \def\nabla{\mathord{\raisebox{\depth}{\scalebox{1}[-1]{$\symup{\Delta}$}}}}%
  \def\checkmark{\popcheck{popdark}}%
  \def\star{\mathbin{\ast}}\def\bigstar{\mathord{\ast}}%
  \def\diamond{\mathbin{\scalebox{0.6}{$\Diamond$}}}%
  \def\bigcup{\mathop{\mathchoice{\raisebox{-0.3em}{\scalebox{1.7}{$\cup$}}}{\scalebox{1.25}{$\cup$}}{\cup}{\cup}}\displaylimits}%
  \def\bigcap{\mathop{\mathchoice{\raisebox{-0.3em}{\scalebox{1.7}{$\cap$}}}{\scalebox{1.25}{$\cap$}}{\cap}{\cap}}\displaylimits}%
  \def\lesseqqgtr{\mathrel{\lessgtr}}\def\bigoplus{\mathop{\oplus}\displaylimits}%
}
\providecommand{\ssi}{si et seulement si} % abréviation fréquente
\AtBeginDocument{\providecommand{\wideparen}{\overparen}} % arc de cercle

%% FIGFIX-BEGIN v3 — correctifs globaux des figures (docs/onbuch-generation/tools/figfix)
\usepackage{adjustbox}\usepackage{etoolbox}
% toute figure TikZ/pgfplots plus large que la ligne est réduite à la largeur disponible
\BeforeBeginEnvironment{tikzpicture}{\begin{adjustbox}{max width=\linewidth}}
\AfterEndEnvironment{tikzpicture}{\end{adjustbox}}
% légendes pgfplots lisibles par-dessus les courbes
\pgfplotsset{every axis/.append style={legend style={fill=white,fill opacity=0.92,text opacity=1,draw=black!15,inner sep=3pt}}}
% étiquettes d'arêtes (syntaxe edge["..."]) avec fond blanc
\tikzset{every edge quotes/.append style={fill=white,inner sep=1.2pt}}
% bulles de cartes mentales : texte à la ligne dans la bulle, bulle un peu plus grande
\tikzset{every concept/.append style={text width=2.0cm,align=center,minimum size=2.3cm,inner sep=2pt}}
%% FIGFIX-END
\begin{document}

\pagegarde{Fonctions associées et étude de fonctions}{Mathématiques}{1ère A}{Relations et opérations fondamentales dans l'ensemble des nombres réels}{NA04}{14 h}{Cette leçon approfondit l'étude des fonctions polynomiales du second degré et homographiques sur des intervalles bornés : ensemble de définition, asymptotes, variations, extrema, tableaux et représentations graphiques. Elle maîtrise ensuite les transformations géométriques (translations, symétries, valeur absolue) pour construire les courbes des fonctions associées sans recalculer.
\vspace{4pt}
\begin{multicols}{2}\small
\begin{itemize}
\item À la fin de cette leçon, je sais déterminer l'ensemble de définition d'une fonction polynôme du second degré ou homographique sur un intervalle borné.
\item Je sais mettre une fonction du second degré sous forme canonique et une fonction homographique sous la forme x ↦ α + k/(x-β) pour en déduire les translations par rapport aux fonctions de référence x ↦ ax² et x ↦ k/x.
\item Je sais déterminer l'équation de l'asymptote verticale d'une fonction homographique et tracer son tableau de variations sur un intervalle borné ne contenant pas la valeur interdite.
\item Je sais dresser le tableau de variations complet d'une fonction polynôme du second degré ou homographique sur un intervalle borné en utilisant le signe de la fonction dérivée.
\item Je sais déterminer les extremums (relatifs et absolus) d'une fonction sur un intervalle borné et construire sa courbe représentative avec rigueur.
\item Je sais résoudre graphiquement une équation f(x) = 0 ou f(x) = m et une inéquation f(x) > 0 sur un intervalle borné en exploitant la courbe représentative.
\end{itemize}\end{multicols}}

\begin{sommaire}
\begin{multicols}{2}
\begin{enumerate}
\item 1. Outils préliminaires : Dérivée, forme canonique et ensemble de définition
\item 2. Étude complète d'une fonction polynôme du second degré sur un intervalle borné
\item 3. Étude complète d'une fonction homographique sur un intervalle borné
\item 4. Fonctions associées : Transformations géométriques des courbes
\item 5. Résolution graphique d'équations et d'inéquations ; Paramètre m
\item 6. Synthèse et stratégie globale d'étude de fonction
\item 7. Modélisation et problèmes concrets (Applications camerounaises)
\item Activité d'intégration
\item Méthodes et exercices résolus
\item Exercices
\item Corrigés des exercices
\item Fiche bilan
\end{enumerate}
\end{multicols}
\end{sommaire}

\begin{objectifs}
\begin{itemize}
\item À la fin de cette leçon, je sais déterminer l'ensemble de définition d'une fonction polynôme du second degré ou homographique sur un intervalle borné.
\item Je sais mettre une fonction du second degré sous forme canonique et une fonction homographique sous la forme x ↦ α + k/(x-β) pour en déduire les translations par rapport aux fonctions de référence x ↦ ax² et x ↦ k/x.
\item Je sais déterminer l'équation de l'asymptote verticale d'une fonction homographique et tracer son tableau de variations sur un intervalle borné ne contenant pas la valeur interdite.
\item Je sais dresser le tableau de variations complet d'une fonction polynôme du second degré ou homographique sur un intervalle borné en utilisant le signe de la fonction dérivée.
\item Je sais déterminer les extremums (relatifs et absolus) d'une fonction sur un intervalle borné et construire sa courbe représentative avec rigueur.
\item Je sais résoudre graphiquement une équation f(x) = 0 ou f(x) = m et une inéquation f(x) > 0 sur un intervalle borné en exploitant la courbe représentative.
\item Je sais construire la courbe des fonctions x ↦ -f(x), x ↦ |f(x)|, x ↦ f(x-a), x ↦ f(x)+b et x ↦ f(x-a)+b à partir de celle de f par transformations géométriques successives.
\end{itemize}
\end{objectifs}

\begin{prerequis}
\begin{itemize}
\item Notion de fonction, ensemble de définition, courbe représentative (Seconde).
\item Résolution d'équations et inéquations du second degré : discriminant, forme factorisée, signe d'un trinôme (Seconde).
\item Nombre dérivé et fonction dérivée des fonctions usuelles (x\^{}n, 1/x, sommes) ; lien signe dérivée / variations (Début 1ère).
\item Géométrie analytique : repère, vecteurs, translations, symétries axiales (Seconde).
\item Fonctions homographiques : définition, ensemble de définition, asymptotes (notions abordées en Seconde ou début 1ère).
\end{itemize}
\end{prerequis}

\newpage

\coursec{Outils préliminaires : Dérivée, forme canonique et ensemble de définition}

\noindent \textbf{Situation d'accroche.}\quad Au marché Central de Yaoundé, Maman Ngo vend des mangues. Elle constate que si elle fixe le prix trop haut, elle en vend peu~; trop bas, son bénéfice diminue. Son fils, élève en 1ère A, modélise son bénéfice journalier $B(x)$ en fonction du prix de vente $x$ (en FCFA) par une fonction du second degré sur l'intervalle $[200\,; 800]$. Comment déterminer le prix optimal pour maximiser le gain~? Et si une taxe fixe de $50$~FCFA par mangue est instaurée, comment la courbe du bénéfice se transforme-t-elle~? Pour répondre à ces questions, nous devons maîtriser les outils de base : la dérivée pour étudier les variations, l'ensemble de définition pour savoir où la fonction «~vit~», et les formes canoniques pour visualiser les transformations géométriques.

\courssub{Rappels sur la dérivation : nombre dérivé, fonction dérivée et règles de calcul}

\begin{definition}[Nombre dérivé en un point]
Soit $f$ une fonction définie sur un intervalle $I$ et $a \in I$. Si le taux d'accroissement $\dfrac{f(x)-f(a)}{x-a}$ admet une limite finie lorsque $x \to a$, on dit que $f$ est \cle{dérivable} en $a$. Cette limite est appelée \cle{nombre dérivé} de $f$ en $a$ et se note $f'(a)$~:
\[
f'(a) = \lim_{x \to a} \frac{f(x)-f(a)}{x-a} = \lim_{h \to 0} \frac{f(a+h)-f(a)}{h}.
\]
Géométriquement, $f'(a)$ est la pente de la tangente à la courbe $\mathcal{C}_f$ au point d'abscisse $a$.
\end{definition}

\begin{definition}[Fonction dérivée sur un intervalle]
Si $f$ est dérivable en tout point d'un intervalle $I$, on définit la \cle{fonction dérivée} $f' : I \to \mathbb{R}$ qui associe à tout $x \in I$ le nombre dérivé $f'(x)$.
\end{definition}

\begin{popfigure}
\begin{tikzpicture}[
    mindmap, grow cyclic, every node/.style=concept, concept color=popblue!60,
    level 1/.append style={level distance=4.5cm, sibling angle=90},
    level 2/.append style={level distance=3.2cm, sibling angle=45},
    level 3/.append style={level distance=2.8cm, sibling angle=30},
    text=white, font=\small\bfseries
]
\node {Règles de dérivation}
    child[concept color=poporange] { node {Linéarité}
        child { node {$ (u+v)' = u' + v' $ } }
        child { node {$ (\lambda u)' = \lambda u' $ } }
    }
    child[concept color=popgreen] { node {Puissances}
        child { node[align=center] {$ (x^n)' = n x^{n-1} $ \\ ($n \in \mathbb{N}$) } }
        child { node {$ (\sqrt{x})' = \frac{1}{2\sqrt{x}} $ } }
    }
    child[concept color=poppurple] { node {Inverse}
        child { node {$ \left(\frac{1}{u}\right)' = -\frac{u'}{u^2} $ } }
    }
    child[concept color=poppink] { node {Composée}
        child { node {$ (u \circ v)' = v' \cdot (u' \circ v) $ } }
    };
\end{tikzpicture}
\legende{Arbre des principales règles de dérivation (niveau 1ère A : on retient surtout la linéarité, les puissances entières naturelles et l'inverse).}
\end{popfigure}

\begin{propriete}[Règles de dérivation usuelles (à connaître par cœur)]
Soient $u$ et $v$ des fonctions dérivables sur un intervalle $I$, $\lambda \in \mathbb{R}$ et $n \in \mathbb{N}$.
\begin{itemize}
    \item \textbf{Somme~:} $(u+v)' = u' + v'$.
    \item \textbf{Multiple par un scalaire~:} $(\lambda u)' = \lambda u'$.
    \item \textbf{Puissance entière naturelle~:} $(x^n)' = n x^{n-1}$.
    \item \textbf{Inverse~:} $\left(\dfrac{1}{u}\right)' = -\dfrac{u'}{u^2}$ (là où $u \neq 0$).
\end{itemize}
\emph{Remarque~:} La dérivée d'une constante est nulle~: $(c)' = 0$.
Pour une fonction homographique mise sous forme standard $\alpha + \frac{k}{x-\beta}$, on utilise la règle de l'inverse.
\end{propriete}

\begin{exemplebox}[Calcul de dérivée]
Soit $f(x) = 3x^2 - 4x + 1$ définie sur $[-1\,; 3]$.
\begin{align*}
f'(x) &= (3x^2)' - (4x)' + (1)' \\
      &= 3 \cdot 2x - 4 \cdot 1 + 0 \\
      &= 6x - 4.
\end{align*}
La fonction dérivée est $f'(x) = 6x - 4$ sur $[-1\,; 3]$.
\end{exemplebox}

\begin{popfigure}
\begin{tikzpicture}[scale=0.9, >=Stealth]
\begin{axis}[
    width=\linewidth, height=7cm, axis lines=middle,
    xlabel=$x$, ylabel=$y$, xlabel style={anchor=north}, ylabel style={anchor=east},
    xmin=-1.5, xmax=3.5, ymin=-5, ymax=15,
    xtick={-1,0,1,2,3}, ytick={-4,0,4,8,12},
    grid=major, grid style={popline}, tick label style={font=\footnotesize},
    clip=false
]
% f(x)
\addplot[popblue, thick, domain=-1:3, samples=100] {3*x^2 - 4*x + 1};
\node[popblue, anchor=south west] at (axis cs:3,16) {$\mathcal{C}_f : y = 3x^2-4x+1$};
% Tangente en x=1 : f'(1)=2, f(1)=0 -> y = 2(x-1) = 2x-2
\addplot[poporange, dashed, domain=-1:3, samples=2] {2*x - 2};
\node[poporange, anchor=south east] at (axis cs:3,4) {Tangente en $1$~: $y = 2x-2$};
% Point de tangence
\fill[poporange] (axis cs:1,0) circle (2pt) node[above right, black] {$A(1\,;0)$};
% Pente triangle
\draw[popmuted, dashed] (axis cs:1,0) -- (axis cs:2,0) -- (axis cs:2,2) -- cycle;
\node[popmuted, font=\scriptsize] at (axis cs:1.5,-0.5) {$1$};
\node[popmuted, font=\scriptsize, rotate=90] at (axis cs:2.15,1) {$2$};
\end{axis}
\end{tikzpicture}
\legende{Illustration géométrique~: $f'(1)=2$ est la pente de la tangente en $A(1\,;f(1))$. L'équation de la tangente est $y=2x-2$.}
\end{popfigure}

\courssub{Lien fondamental entre le signe de $f'$ et le sens de variation de $f$}

C'est le théorème «~moteur~» de tout l'étude de fonctions. Il permet de remplacer l'étude algébrique fastidieuse de $f(x_2)-f(x_1)$ par l'étude du signe de la dérivée.

\begin{propriete}[Théorème : Signe de la dérivée et monotonie]
Soit $f$ une fonction continue sur un intervalle $I$ et dérivable sur l'intérieur de $I$ (noté $\mathring{I}$).
\begin{itemize}
    \item Si $f'(x) \ge 0$ pour tout $x \in \mathring{I}$, alors $f$ est \cle{croissante} sur $I$.
    \item Si $f'(x) \le 0$ pour tout $x \in \mathring{I}$, alors $f$ est \cle{décroissante} sur $I$.
    \item Si $f'(x) = 0$ pour tout $x \in \mathring{I}$, alors $f$ est \cle{constante} sur $I$.
\end{itemize}
\end{propriete}

\begin{experience}[Démonstration guidée par le Théorème des Accroissements Finis (TAF)]
\textbf{Objectif~:} Prouver que $f' \ge 0 \implies f$ croissante.
\begin{enumerate}
    \item Soient $a, b \in I$ avec $a < b$. La fonction $f$ est continue sur $[a,b]$ et dérivable sur $]a,b[$.
    \item Appliquer le \textbf{Théorème des Accroissements Finis}~: il existe $c \in ]a,b[$ tel que
    \[
    f'(c) = \frac{f(b)-f(a)}{b-a}.
    \]
    \item Puisque $b-a > 0$ et $f'(c) \ge 0$ (hypothèse), le quotient est $\ge 0$, donc $f(b)-f(a) \ge 0$, c'est-à-dire $f(a) \le f(b)$.
    \item Conclusion~: $f$ est croissante sur $I$.
\end{enumerate}
\emph{Note~:} La réciproque est vraie si $f$ est dérivable~: $f$ croissante $\implies f' \ge 0$.
\end{experience}

\begin{popfigure}
\begin{tikzpicture}[scale=0.9, >=Stealth]
\begin{axis}[
    width=\linewidth, height=6cm, axis lines=middle,
    xlabel=$x$, ylabel=$y$, xlabel style={anchor=north}, ylabel style={anchor=east},
    xmin=-1, xmax=4, ymin=-1, ymax=4,
    xtick={0,1,2,3}, ytick={0,1,2,3},
    grid=major, grid style={popline}, tick label style={font=\footnotesize},
    clip=false
]
% f croissante
\addplot[popgreen, thick, domain=0:3, samples=100] {0.5*x + 0.5};
\node[popgreen, anchor=south west] at (axis cs:3,2) {$f$ croissante, $f'>0$};
% Tangentes
\addplot[popgreen!50, dashed, domain=0.09999999999999998:0.9, samples=2] {0.5*x + 0.5};
\fill[popgreen] (axis cs:0.5, 0.75) circle (1.5pt);
\addplot[popgreen!50, dashed, domain=1.1:1.9, samples=2] {0.5*x + 0.5};
\fill[popgreen] (axis cs:1.5, 1.25) circle (1.5pt);
\addplot[popgreen!50, dashed, domain=2.1:2.9, samples=2] {0.5*x + 0.5};
\fill[popgreen] (axis cs:2.5, 1.75) circle (1.5pt);
\end{axis}
\end{tikzpicture}
\legende{Si $f'>0$, les tangentes ont une pente positive~: la fonction monte.}
\end{popfigure}

\begin{attention}[Pièges fréquents sur le signe de $f'$]
\begin{itemize}
    \item \textbf{Confondre $f$ et $f'$~:} $f'(x) > 0$ ne veut pas dire que $f(x) > 0$~! La fonction peut être négative mais croissante (ex~: $f(x)=x-5$ sur $[0\,;4]$, $f'=1>0$, $f$ croissante mais $f(x)<0$).
    \item \textbf{Oublier les bornes de l'intervalle~:} Le théorème s'applique sur un \emph{intervalle}. Si $f'$ change de signe, on découpe en sous-intervalles.
    \item \textbf{Points non dérivables~:} Le théorème ne s'applique pas aux points où $f$ n'est pas dérivable (ex~: valeur absolue, racine carrée en $0$, asymptote verticale). On traite ces points séparément.
    \item \textbf{Intervalle borné imposé~:} En 1ère A, l'étude se fait \textbf{toujours} sur un intervalle borné $I=[a;b]$ donné. Le tableau de variations ne comporte que les valeurs $a$, $b$ et les valeurs critiques \emph{appartenant à $I$}.
\end{itemize}
\end{attention}

\courssub{Ensemble de définition $D_f$ sur un intervalle borné}

C'est la \textbf{première étape} de toute étude de fonction. On ne peut pas dériver ni tracer une courbe là où la fonction n'existe pas.

\begin{propriete}[Ensemble de définition]
\begin{itemize}
    \item \textbf{Fonction polynôme~:} $D_f = \mathbb{R}$. Si l'énoncé impose un intervalle borné $I = [a\,; b]$, alors $D_f = I$.
    \item \textbf{Fonction homographique~:} $f(x) = \dfrac{ax+b}{cx+d}$ avec $c \neq 0$ et $ad-bc \neq 0$. Le dénominateur s'annule pour $x_0 = -\dfrac{d}{c}$.
    \[
    D_f = \mathbb{R} \setminus \{x_0\}.
    \]
    Si l'énoncé impose un intervalle borné $I = [a\,; b]$, on \cle{intersecte}~:
    \[
    D_f = I \setminus \{x_0\} \quad \text{(si $x_0 \in I$)}, \qquad D_f = I \quad \text{(si $x_0 \notin I$)}.
    \]
\end{itemize}
\end{propriete}

\begin{attention}[Intersection avec l'intervalle borné : piège classique]
On ne dit jamais «~$D_f = \mathbb{R} \setminus \{3\}$~» si l'énoncé dit «~$f$ définie sur $[0\,;5]$~».
\begin{itemize}
    \item Si $f(x) = \dfrac{2x+1}{x-3}$ sur $[0\,;5]$, alors $D_f = [0\,;5] \setminus \{3\} = [0\,;3[ \cup ]3\,;5]$.
    \item Si $f(x) = \dfrac{2x+1}{x-6}$ sur $[0\,;5]$, alors $D_f = [0\,;5]$ (car $6 \notin [0\,;5]$).
\end{itemize}
\textbf{Toujours écrire explicitement $D_f$ sous forme d'intervalle(s) ou de réunion d'intervalles.}
\end{attention}

\begin{exemplebox}[Détermination de $D_f$]
Soit $f(x) = \dfrac{3x-2}{x+1}$ définie sur $[-4\,; 2]$.
\begin{enumerate}
    \item Dénominateur nul pour $x+1=0 \iff x = -1$.
    \item $-1 \in [-4\,; 2]$.
    \item $D_f = [-4\,; 2] \setminus \{-1\} = [-4\,; -1[ \cup ]-1\,; 2]$.
\end{enumerate}
\end{exemplebox}

\courssub{Forme canonique du trinôme du second degré et translation vectorielle}

Pour une fonction polynôme du second degré $f(x) = ax^2 + bx + c$ ($a \neq 0$), la forme canonique révèle le sommet du parabole et la translation par rapport à $x \mapsto ax^2$.

\begin{methode}[Mise sous forme canonique $f(x) = a(x-\alpha)^2 + \beta$]
Pour $f(x) = ax^2 + bx + c$~:
\begin{enumerate}
    \item Factoriser $a$ sur les deux premiers termes~: $f(x) = a\left(x^2 + \frac{b}{a}x\right) + c$.
    \item Compléter le carré~: ajouter et soustraire $\left(\frac{b}{2a}\right)^2$ à l'intérieur.
    \item Factoriser le trinôme carré parfait.
    \item Simplifier la constante.
\end{enumerate}
\textbf{Résultat~:} $\alpha = -\dfrac{b}{2a}$, $\beta = -\dfrac{\Delta}{4a}$ avec $\Delta = b^2 - 4ac$.
Le sommet du parabole est $S(\alpha\,; \beta)$.
La courbe $\mathcal{C}_f$ s'obtient par \cle{translation} de la courbe $\mathcal{C}_0 : y = ax^2$ par le vecteur $\vec{v}(\alpha\,; \beta)$.
\end{methode}

\begin{aretenir}[Forme canonique et translation -- Polynôme degré 2]
Pour $f(x) = ax^2+bx+c$ ($a \neq 0$)~:
\begin{itemize}
    \item Forme canonique : $f(x) = a(x-\alpha)^2 + \beta$ avec $\alpha = -\frac{b}{2a}$, $\beta = f(\alpha) = -\frac{\Delta}{4a}$.
    \item Sommet du parabole : $S(\alpha\,; \beta)$.
    \item Vecteur de translation de $\mathcal{C}_0 : y = ax^2$ vers $\mathcal{C}_f$ : $\vec{v}(\alpha\,; \beta)$.
    \item Sens de variation : $f$ décroît puis croît si $a>0$ (minimum en $\alpha$) ; $f$ croît puis décroît si $a<0$ (maximum en $\alpha$).
\end{itemize}
\end{aretenir}

\begin{popfigure}
\begin{tikzpicture}
\begin{axis}[
    width=\linewidth, height=7cm, axis lines=middle,
    xlabel=$x$, ylabel=$y$, xlabel style={anchor=north}, ylabel style={anchor=east},
    xmin=-1, xmax=5, ymin=-4, ymax=6,
    xtick={-1,0,1,2,3,4}, ytick={-3,-2,-1,0,1,2,3,4,5},
    grid=major, grid style={popline}, tick label style={font=\footnotesize},
    clip=false
]
% y = x^2 (reference)
\addplot[popblue!40, thick, dashed, domain=-1:4, samples=100] {x^2};
\node[popblue!60, anchor=south east] at (axis cs:4,16) {$\mathcal{C}_0 : y = x^2$};
% y = (x-2)^2 + 1
\addplot[poporange, very thick, domain=-0.5:4.5, samples=100] {(x-2)^2 + 1};
\node[poporange, anchor=south west] at (axis cs:4.5,7) {$\mathcal{C}_f : y = (x-2)^2+1$};
% Sommet S(2,1)
\fill[poporange] (axis cs:2,1) circle (2.5pt) node[above right, black] {$S(2\,;1)$};
% Vecteur translation
\draw[poppurple, very thick, ->, >=Stealth] (axis cs:0,0) -- (axis cs:2,1) node[midway, above left, poppurple] {$\vec{v}(2\,;1)$};
% Repères
\draw[popmuted, dashed] (axis cs:0,0) -- (axis cs:2,0) -- (axis cs:2,1);
\draw[popmuted, dashed] (axis cs:0,0) -- (axis cs:0,1) -- (axis cs:2,1);
\end{axis}
\end{tikzpicture}
\legende{Translation de $y=x^2$ par $\vec{v}(2\,;1)$ pour obtenir $y=(x-2)^2+1$. Le sommet $S$ est l'image de $O(0,0)$.}
\end{popfigure}

\begin{exemplebox}[Forme canonique : $f(x) = 2x^2 - 8x + 5$]
\begin{align*}
f(x) &= 2(x^2 - 4x) + 5 \\
     &= 2\left[(x^2 - 4x + 4) - 4\right] + 5 \quad \text{(car $(-4/2)^2 = 4$)} \\
     &= 2\left[(x-2)^2 - 4\right] + 5 \\
     &= 2(x-2)^2 - 8 + 5 \\
     &= 2(x-2)^2 - 3.
\end{align*}
Forme canonique~: $f(x) = 2(x-2)^2 - 3$.
Sommet $S(2\,; -3)$.
Translation de $y = 2x^2$ par le vecteur $\vec{v}(2\,; -3)$.
\end{exemplebox}

\courssub{Forme standard de la fonction homographique et translation vectorielle}

Pour une fonction homographique $f(x) = \dfrac{ax+b}{cx+d}$ ($c \neq 0$, $ad-bc \neq 0$), la forme standard révèle le centre de symétrie et la translation par rapport à $x \mapsto \dfrac{k}{x}$.

\begin{methode}[Mise sous forme standard $f(x) = \alpha + \dfrac{k}{x-\beta}$]
On effectue la \textbf{division euclidienne} du numérateur par le dénominateur.
Pour $f(x) = \dfrac{ax+b}{cx+d}$~:
\begin{enumerate}
    \item Diviser $ax+b$ par $cx+d$~: le quotient est $\dfrac{a}{c}$, le reste est $b - \dfrac{ad}{c} = \dfrac{bc-ad}{c}$.
    \item Écrire~: $f(x) = \dfrac{a}{c} + \dfrac{bc-ad}{c(cx+d)}$.
    \item Mettre au format $\alpha + \dfrac{k}{x-\beta}$ en factorisant $c$ au dénominateur du reste~:
    \[
    f(x) = \underbrace{\frac{a}{c}}_{\alpha} + \underbrace{\frac{bc-ad}{c^2}}_{k} \cdot \frac{1}{x - \underbrace{\left(-\frac{d}{c}\right)}_{\beta}}.
    \]
\end{enumerate}
\textbf{Résultat~:} $\alpha = \dfrac{a}{c}$, $\beta = -\dfrac{d}{c}$, $k = \dfrac{bc-ad}{c^2}$.
Le centre de symétrie de la courbe (hyperbole) est $I(\beta\,; \alpha)$.
La courbe $\mathcal{C}_f$ s'obtient par \cle{translation} de la courbe $\mathcal{C}_0 : y = \dfrac{k}{x}$ par le vecteur $\vec{v}(\beta\,; \alpha)$.
\end{methode}

\begin{aretenir}[Forme standard et translation -- Fonction homographique]
Pour $f(x) = \frac{ax+b}{cx+d}$ ($c \neq 0$, $ad-bc \neq 0$)~:
\begin{itemize}
    \item Forme standard : $f(x) = \alpha + \frac{k}{x-\beta}$ avec $\alpha = \frac{a}{c}$, $\beta = -\frac{d}{c}$, $k = \frac{bc-ad}{c^2}$.
    \item Centre de symétrie (intersection des asymptotes) : $I(\beta\,; \alpha)$.
    \item Asymptotes : verticale $x = \beta$, horizontale $y = \alpha$.
    \item Vecteur de translation de $\mathcal{C}_0 : y = \frac{k}{x}$ vers $\mathcal{C}_f$ : $\vec{v}(\beta\,; \alpha)$.
    \item Sens de variation : sur chaque intervalle de $D_f$, $f$ est strictement croissante si $k<0$, strictement décroissante si $k>0$.
\end{itemize}
\end{aretenir}

\begin{popfigure}
\begin{tikzpicture}
\begin{axis}[
    width=\linewidth, height=7cm, axis lines=middle,
    xlabel=$x$, ylabel=$y$, xlabel style={anchor=north}, ylabel style={anchor=east},
    xmin=-5, xmax=3, ymin=-3, ymax=5,
    xtick={-5,-4,-3,-2,-1,0,1,2}, ytick={-2,-1,0,1,2,3,4},
    grid=major, grid style={popline}, tick label style={font=\footnotesize},
    clip=false, restrict y to domain=-3:5
]
% y = 1/x (reference branches)
\addplot[popblue!40, thick, dashed, domain=-5:-0.3, samples=100] {1/x};
\addplot[popblue!40, thick, dashed, domain=0.3:3, samples=100] {1/x};
\node[popblue!60, anchor=south east] at (axis cs:3,0.33) {$\mathcal{C}_0 : y = 1/x$};
% y = 2 + 3/(x+1) -> k=3, beta=-1, alpha=2. Asymptotes x=-1, y=2.
\addplot[poporange, very thick, domain=-5:-1.3, samples=100] {2 + 3/(x+1)};
\addplot[poporange, very thick, domain=-0.7:3, samples=100] {2 + 3/(x+1)};
\node[poporange, anchor=south west] at (axis cs:3,2.75) {$\mathcal{C}_f : y = 2+\frac{3}{x+1}$};
% Asymptotes
\draw[popmuted, densely dashed] (axis cs:-1,-3) -- (axis cs:-1,5);
\draw[popmuted, densely dashed] (axis cs:-5,2) -- (axis cs:3,2);
% Centre I(-1,2)
\fill[poppurple] (axis cs:-1,2) circle (2.5pt) node[above right, black] {$I(-1\,;2)$};
% Vecteur translation
\draw[poppurple, very thick, ->, >=Stealth] (axis cs:0,0) -- (axis cs:-1,2) node[midway, above right, poppurple] {$\vec{v}(-1\,;2)$};
\end{axis}
\end{tikzpicture}
\legende{Translation de $y=1/x$ par $\vec{v}(-1\,;2)$ pour obtenir $y=2+\frac{3}{x+1}$. Le centre $I$ est l'image de $O$.}
\end{popfigure}

\begin{exemplebox}[{Forme standard : $f(x) = \dfrac{3x-1}{x+2}$ sur $[-5\,; 5]$}]
\begin{enumerate}
    \item \textbf{Ensemble de définition~:} Dénominateur nul pour $x=-2 \in [-5\,;5]$. $D_f = [-5\,;-2[ \cup ]-2\,;5]$.
    \item \textbf{Division euclidienne~:}
    \[
    \frac{3x-1}{x+2} = 3 - \frac{7}{x+2} = 3 + \frac{-7}{x - (-2)}.
    \]
    \item \textbf{Identification~:} $\alpha = 3$, $\beta = -2$, $k = -7$.
    \item \textbf{Centre de symétrie~:} $I(-2\,; 3)$.
    \item \textbf{Translation~:} La courbe s'obtient par translation de $y = \dfrac{-7}{x}$ par le vecteur $\vec{v}(-2\,; 3)$.
\end{enumerate}
\end{exemplebox}

\courssub{Exercices d'entraînement ciblés (Résolus)}

\begin{exoresolu}[Exercice 1 : Dérivée et variations d'un polynôme du second degré sur intervalle borné]
\textbf{Énoncé.}\quad Soit $f(x) = -2x^2 + 4x + 6$ définie sur $[-1\,; 3]$.
\begin{enumerate}
    \item Calculer $f'(x)$.
    \item Étudier le signe de $f'(x)$ sur $[-1\,; 3]$.
    \item En déduire le tableau de variations de $f$ sur $[-1\,; 3]$ et les extrema.
\end{enumerate}
\tcblower
\textbf{Solution.}
\begin{enumerate}
    \item $f'(x) = -4x + 4 = -4(x-1)$.
    \item $f'(x) = 0 \iff x=1$. Cette valeur appartient à $[-1\,;3]$.
    Tableau de signes de $f'$ (6 colonnes : 3 valeurs remarquables, 2 intervalles)~:
    \begin{center}
    \begin{tabular}{|c|ccccc|}\hline
    $x$ & $-1$ & & $1$ & & $3$ \\ \hline
    $f'(x)$ & & $+$ & $0$ & $-$ & \\ \hline
    \end{tabular}
    \end{center}
    \item Tableau de variations de $f$~:
    \begin{center}
    \begin{tabular}{|c|ccccc|}\hline
    $x$ & $-1$ & & $1$ & & $3$ \\ \hline
    $f'(x)$ & & $+$ & $0$ & $-$ & \\ \hline
    $f$ & $0$ & $\nearrow$ & $8$ & $\searrow$ & $0$ \\ \hline
    \end{tabular}
    \end{center}
    $f(-1)=0$, $f(1)=8$, $f(3)=0$.
    Maximum $8$ atteint en $1$ ; Minimum $0$ atteint en $-1$ et $3$.
\end{enumerate}
\end{exoresolu}

\begin{exoresolu}[Exercice 2 : Forme canonique et translation]
\textbf{Énoncé.}\quad Soit $f(x) = -x^2 + 6x - 5$.
\begin{enumerate}
    \item Mettre $f$ sous forme canonique.
    \item Donner les coordonnées du sommet $S$ de la courbe $\mathcal{C}_f$.
    \item Préciser le vecteur de translation permettant d'obtenir $\mathcal{C}_f$ à partir de $\mathcal{C}_0 : y = -x^2$.
\end{enumerate}
\tcblower
\textbf{Solution.}
\begin{enumerate}
    \item $f(x) = -(x^2 - 6x) - 5 = -\left[(x-3)^2 - 9\right] - 5 = -(x-3)^2 + 9 - 5 = -(x-3)^2 + 4$.
    \item Sommet $S(3\,; 4)$.
    \item Translation de $y = -x^2$ par le vecteur $\vec{v}(3\,; 4)$.
\end{enumerate}
\end{exoresolu}

\begin{exoresolu}[Exercice 3 : Fonction homographique, forme standard et $D_f$]
\textbf{Énoncé.}\quad Soit $f(x) = \dfrac{2x+3}{x-1}$ définie sur $[0\,; 4]$.
\begin{enumerate}
    \item Déterminer $D_f$.
    \item Écrire $f$ sous la forme $\alpha + \dfrac{k}{x-\beta}$.
    \item Donner le centre de symétrie $I$ et le vecteur de translation par rapport à $y = \dfrac{k}{x}$.
\end{enumerate}
\tcblower
\textbf{Solution.}
\begin{enumerate}
    \item Dénominateur nul pour $x=1$. Comme $1 \in [0\,;4]$, $D_f = [0\,;1[ \cup ]1\,;4]$.
    \item Division~: $\dfrac{2x+3}{x-1} = 2 + \dfrac{5}{x-1}$.
    Donc $\alpha = 2$, $\beta = 1$, $k = 5$.
    \item Centre de symétrie $I(1\,; 2)$.
    Vecteur de translation $\vec{v}(1\,; 2)$ par rapport à la courbe $y = \dfrac{5}{x}$.
\end{enumerate}
\end{exoresolu}

\begin{savaistu}[Le saviez-vous ?]
La forme canonique $a(x-\alpha)^2+\beta$ et la forme standard $\alpha+\frac{k}{x-\beta}$ sont les «~codes-barres~» des fonctions du second degré et homographiques. Elles disent tout~: position du sommet/centre, orientation (signe de $a$ ou $k$), et dilatation (valeur absolue de $a$ ou $k$). Au Probatoire, on vous demandera souvent : «~Écrire sous forme canonique/standard~» puis «~En déduire la translation~». C'est un savoir-faire \textbf{incontournable} et très rentable en points~!
\end{savaistu}

\coursec{Étude complète d'une fonction polynôme du second degré sur un intervalle borné}

Dans la section précédente, nous avons rappelé les outils indispensables : la dérivée d'un polynôme du second degré, sa forme canonique et la notion d'ensemble de définition. Nous savons désormais qu'une fonction $f(x)=ax^2+bx+c$ (avec $a \neq 0$) est dérivable sur $\mathbb{R}$, que sa dérivée $f'(x)=2ax+b$ s'annule en un unique point $\alpha = -\frac{b}{2a}$ (l'abscisse du sommet de la parabole), et que sa forme canonique $f(x)=a(x-\alpha)^2+\beta$ révèle immédiatement les coordonnées du sommet $S(\alpha\,;\,\beta)$.

Nous allons maintenant mettre ces outils en œuvre pour réaliser une \textbf{étude complète et rigoureuse} de cette fonction sur un \textbf{intervalle borné fermé} $I=[m\,;\,n]$. C'est une compétence centrale du programme : savoir dresser un tableau de variations, déterminer les extremums absolus, tracer la courbe et résoudre graphiquement équations et inéquations sur cet intervalle.

\courssub{Dérivée, signe et variations sur un intervalle borné}

\textbf{Intuition.}\par
La dérivée $f'(x)$ mesure le taux de variation instantané. Pour un trinôme du second degré, $f'(x)=2ax+b$ est une fonction affine. Son signe est constant de part et d'autre de son unique zéro $\alpha$. Sur un intervalle borné $[m\,;\,n]$, trois cas se présentent : $\alpha$ est à l'extérieur de l'intervalle (à gauche ou à droite), ou $\alpha$ se trouve à l'intérieur. C'est cette position relative qui dicte la forme du tableau de variations.

\begin{propriete}[Variations d'un polynôme du second degré sur $\mathbb{R}$]
Soit $f(x)=ax^2+bx+c$ avec $a \neq 0$ et $\alpha = -\frac{b}{2a}$.
\begin{itemize}
    \item Si $a > 0$ : $f'(x) < 0$ sur $]-\infty\,;\,\alpha[$, $f'(\alpha)=0$, $f'(x) > 0$ sur $]\alpha\,;\,+\infty[$. 
          La fonction est \textbf{décroissante} puis \textbf{croissante}. Elle admet un \textbf{minimum} en $\alpha$.
    \item Si $a < 0$ : $f'(x) > 0$ sur $]-\infty\,;\,\alpha[$, $f'(\alpha)=0$, $f'(x) < 0$ sur $]\alpha\,;\,+\infty[$. 
          La fonction est \textbf{croissante} puis \textbf{décroissante}. Elle admet un \textbf{maximum} en $\alpha$.
\end{itemize}
La valeur de l'extremum (relatif sur $\mathbb{R}$) est $f(\alpha) = -\frac{\Delta}{4a}$ où $\Delta = b^2-4ac$.
\end{propriete}

\begin{attention}[{Restriction à l'intervalle $I=[m\,;\,n]$}]
Le tableau de variations sur $I$ ne conserve que la portion de la droite réelle comprise entre $m$ et $n$.
\begin{itemize}
    \item Si $\alpha \notin [m\,;\,n]$ (cas $\alpha \le m$ ou $\alpha \ge n$) : la fonction est \textbf{strictement monotone} sur tout l'intervalle. Pas de case intermédiaire pour $\alpha$ dans le tableau.
    \item Si $\alpha \in ]m\,;\,n[$ : le sommet est \textbf{à l'intérieur} de l'intervalle. Le tableau comporte une case pour $\alpha$ avec une flèche changeant de sens.
    \item Si $\alpha = m$ ou $\alpha = n$ : l'extremum est atteint \textbf{au bord} de l'intervalle. On ne double pas la colonne.
\end{itemize}
\end{attention}

\begin{methode}[{Construction du tableau de variations sur $I=[m\,;\,n]$}]
\begin{enumerate}
    \item Calculer $f'(x) = 2ax+b$.
    \item Résoudre $f'(x)=0$ pour trouver $\alpha = -\frac{b}{2a}$.
    \item \textbf{Étudier le signe de $f'(x)$ sur $[m\,;\,n]$} : 
          \begin{itemize}
              \item Tracer une droite graduée pour $x$ de $m$ à $n$.
              \item Placer $\alpha$ \textbf{uniquement si} $m < \alpha < n$.
              \item Déterminer le signe de $f'$ sur chaque sous-intervalle (tester une valeur ou utiliser le signe de $a$).
          \end{itemize}
    \item \textbf{Dresser le tableau de variations de $f$} :
          \begin{itemize}
              \item Ligne $x$ : valeurs $m$, $\alpha$ (si intérieur), $n$.
              \item Ligne $f'(x)$ : signes $+$, $0$, $-$ (ou l'inverse).
              \item Ligne $f(x)$ : valeurs $f(m)$, $f(\alpha)$, $f(n)$ aux endroits correspondants ; flèches $\nearrow$ (croissant) ou $\searrow$ (décroissant) entre les valeurs.
          \end{itemize}
\end{enumerate}
\end{methode}

\courssub{Extremums absolus sur un intervalle borné fermé}

\textbf{Intuition.}\par
Sur un intervalle \textbf{fermé} et \textbf{borné}, une fonction continue (ce qui est le cas de tout polynôme) atteint \textbf{obligatoirement} un maximum et un minimum. Ce sont les \textbf{extremums absolus} (ou globaux) sur cet intervalle. Ils ne se trouvent pas nécessairement au sommet de la parabole ! Il faut comparer les valeurs aux bornes et la valeur au sommet (s'il est dans l'intervalle).

\begin{propriete}[Théorème des bornes (ou Théorème de Heine) — Admis au Probatoire]
Toute fonction continue sur un intervalle borné fermé $[m\,;\,n]$ est bornée et atteint ses bornes : il existe $x_{\min}, x_{\max} \in [m\,;\,n]$ tels que 
\[ \forall x \in [m\,;\,n], \quad f(x_{\min}) \le f(x) \le f(x_{\max}). \]
$f(x_{\min})$ est le \textbf{minimum absolu}, $f(x_{\max})$ le \textbf{maximum absolu} sur $[m\,;\,n]$.
\end{propriete}

\begin{aretenir}[{Recherche des extremums absolus sur $I=[m\,;\,n]$}]
Pour $f(x)=ax^2+bx+c$ sur $[m\,;\,n]$ :
\begin{enumerate}
    \item Calculer $f(m)$ et $f(n)$ (valeurs aux bornes).
    \item Calculer $\alpha = -\frac{b}{2a}$.
    \item \textbf{Si $\alpha \in [m\,;\,n]$} : calculer $f(\alpha)$. Comparer les trois valeurs $f(m), f(\alpha), f(n)$.
    \item \textbf{Si $\alpha \notin [m\,;\,n]$} : comparer seulement $f(m)$ et $f(n)$.
\end{enumerate}
Le plus grand est le maximum absolu, le plus petit est le minimum absolu.
\end{aretenir}

\begin{exemplebox}[{Exemple : $f(x) = x^2 - 6x + 5$ sur $I=[0\,;\,4]$}]
$f'(x) = 2x - 6$. $f'(x)=0 \Leftrightarrow x=3$. Donc $\alpha = 3 \in [0\,;\,4]$.
$a=1 > 0$ : parabole tournée vers le haut $\rightarrow$ minimum en $\alpha$.
Calculs :
$f(0) = 5$
$f(3) = 9 - 18 + 5 = -4$
$f(4) = 16 - 24 + 5 = -3$

\begin{center}
\begin{tabular}{|c|ccccc|}\hline
$x$ & $0$ & & $3$ & & $4$ \\ \hline
$f'(x)$ & & $-$ & $0$ & $+$ & \\ \hline
$f(x)$ & $5$ & $\searrow$ & $-4$ & $\nearrow$ & $-3$ \\ \hline
\end{tabular}
\end{center}

\textbf{Conclusion :} Minimum absolu $-4$ (atteint en $x=3$). Maximum absolu $5$ (atteint en $x=0$).
Remarque : Le maximum n'est pas au sommet, mais à la borne gauche car la parabole monte plus haut à gauche qu'à droite sur cet intervalle.
\end{exemplebox}

\courssub{Intersections avec les axes et construction de la courbe}

\textbf{Intersections avec l'axe des ordonnées (si $0 \in I$).}\par
Le point d'ordonnée à l'origine est $A(0\,;\,f(0)) = A(0\,;\,c)$.

\textbf{Intersections avec l'axe des abscisses (zéros de la fonction).}\par
On résout $f(x)=0 \Leftrightarrow ax^2+bx+c=0$.
Discriminant $\Delta = b^2-4ac$.
\begin{itemize}
    \item $\Delta < 0$ : Pas d'intersection (la courbe ne coupe pas l'axe $Ox$).
    \item $\Delta = 0$ : Un point de contact (sommet sur l'axe) $x_0 = -\frac{b}{2a} = \alpha$.
    \item $\Delta > 0$ : Deux intersections $x_1 = \frac{-b-\sqrt{\Delta}}{2a}$, $x_2 = \frac{-b+\sqrt{\Delta}}{2a}$ (avec $x_1 < x_2$).
          \textbf{Attention :} On ne retient que les racines appartenant à l'intervalle d'étude $I$.
\end{itemize}

\textbf{Concavité.}\par
Le signe de $a$ donne la concavité : $a>0 \Rightarrow$ convexe (en forme de $\cup$), $a<0 \Rightarrow$ concave (en forme de $\cap$).

\begin{methode}[{Tracer la courbe $\mathcal{C}_f$ de $f(x)=ax^2+bx+c$ sur $[m\,;\,n]$}]
\begin{enumerate}
    \item Choisir une échelle adaptée (repère orthonormé si possible, sinon repère quelconque avec échelles distinctes).
    \item Tracer les points caractéristiques :
          \begin{itemize}
              \item Sommet $S(\alpha\,;\,f(\alpha))$ si $\alpha \in [m\,;\,n]$.
              \item Points aux bornes $M(m\,;\,f(m))$ et $N(n\,;\,f(n))$.
              \item Intersections avec $Ox$ (racines dans $I$) et $Oy$ (si $0 \in I$).
              \item Eventuellement 1 ou 2 points supplémentaires pour la précision (ex: $x = \alpha \pm 1$).
          \end{itemize}
    \item Relier les points par un arc de parabole lisse, respectant la concavité (signe de $a$) et le sens des variations.
    \item Indiquer les coordonnées des points remarquables.
\end{enumerate}
\end{methode}

\courssub{Application complète : $f(x) = -x^2 + 4x - 1$ sur $I=[0\,;\,5]$}

Nous allons traiter cet exemple de A à Z : dérivée, variations, extremums, intersections, courbe, résolution graphique.

\begin{exoresolu}[{Étude complète de $f(x) = -x^2 + 4x - 1$ sur $[0\,;\,5]$}]
\textbf{1. Dérivée et variations.}
$f'(x) = -2x + 4$.
$f'(x)=0 \Leftrightarrow -2x+4=0 \Leftrightarrow x=2$. Donc $\alpha = 2 \in [0\,;\,5]$.
$a = -1 < 0$ : $f'$ positif puis négatif $\rightarrow$ $f$ croissante puis décroissante. Maximum en $2$.

\textbf{2. Tableau de signes de $f'$ sur $[0\,;\,5]$.}
\begin{center}
\begin{tabular}{|c|ccc|ccc|}\hline
$x$ & $0$ & & $2$ & & $5$ \\ \hline
$f'(x)$ & & $+$ & $0$ & $-$ & \\ \hline
\end{tabular}
\end{center}

\textbf{3. Valeurs aux points remarquables.}
$f(0) = -1$.
$f(2) = -4 + 8 - 1 = 3$.
$f(5) = -25 + 20 - 1 = -6$.

\textbf{4. Tableau de variations de $f$ sur $[0\,;\,5]$.}
\begin{center}
\begin{tabular}{|c|ccccc|}\hline
$x$ & $0$ & & $2$ & & $5$ \\ \hline
$f'(x)$ & & $+$ & $0$ & $-$ & \\ \hline
$f(x)$ & $-1$ & $\nearrow$ & $3$ & $\searrow$ & $-6$ \\ \hline
\end{tabular}
\end{center}

\textbf{5. Extremums absolus sur $[0\,;\,5]$.}
Comparaison : $f(0)=-1$, $f(2)=3$, $f(5)=-6$.
Maximum absolu : $3$ (atteint en $x=2$, au sommet).
Minimum absolu : $-6$ (atteint en $x=5$, à la borne droite).

\textbf{6. Intersections avec les axes.}
\begin{itemize}
    \item Axe $Oy$ : $0 \in I$, $f(0)=-1 \Rightarrow A(0\,;\,-1)$.
    \item Axe $Ox$ : $\Delta = 16 - 4 = 12 > 0$. $\sqrt{\Delta} = 2\sqrt{3}$.
          $x_1 = \frac{-4 - 2\sqrt{3}}{-2} = 2 - \sqrt{3} \approx \num{0,27} \in I$.
          $x_2 = \frac{-4 + 2\sqrt{3}}{-2} = 2 + \sqrt{3} \approx \num{3,73} \in I$.
          Points $B(2-\sqrt{3}\,;\,0)$ et $C(2+\sqrt{3}\,;\,0)$.
\end{itemize}

\textbf{7. Construction de la courbe $\mathcal{C}_f$.}
Voir la figure ci-dessous. Sommet $S(2\,;\,3)$, concavité tournée vers le bas ($a<0$).
\tcblower
\textbf{Solution détaillée.}
Tous les calculs sont montrés ci-dessus. La courbe est tracée dans la figure ci-après (générée par le code TikZ/pgfplots). On y voit clairement le sommet en $(2,3)$, les intersections avec $Ox$ vers $\num{0,27}$ et $\num{3,73}$, le point $A(0,-1)$ et le point à la borne droite $(5,-6)$.
\end{exoresolu}

\begin{popfigure}
\begin{tikzpicture}
\begin{axis}[
    width=\linewidth,
    height=9cm,
    axis lines=middle,
    xlabel=$x$, ylabel=$y$,
    xmin=-0.5, xmax=5.5,
    ymin=-7, ymax=4,
    xtick={0,1,2,3,4,5},
    ytick={-6,-4,-2,0,2,3},
    minor tick num=1,
    grid=major,
    grid style={popline, opacity=0.5},
    tick label style={font=\small},
    label style={font=\small},
    clip=false,
    domain=0:5,
    samples=200,
]
% Fonction f(x) = -x^2 + 4x - 1
\addplot[popcurve, thick, domain=0:5] {-x^2 + 4*x - 1};
% Points remarquables
\addplot[only marks, mark=*, mark size=2.5pt, popdark] coordinates {
    (0,-1)       % A
    (2,3)        % S Sommet
    (5,-6)       % N Borne droite
    (2-1.732,0)  % B approx
    (2+1.732,0)  % C approx
};
% Labels
\node[popdark, anchor=south east, font=\small] at (axis cs:0,-1) {$A(0;-1)$};
\node[popdark, anchor=south, font=\small] at (axis cs:2,3) {$S(2;3)$};
\node[popdark, anchor=north west, font=\small] at (axis cs:5,-6) {$N(5;-6)$};
\node[popdark, anchor=north, font=\small] at (axis cs:0.268,0) {$B$};
\node[popdark, anchor=north, font=\small] at (axis cs:3.732,0) {$C$};
% Droite y=3 (max) pour visualisation
\addplot[poporange, dashed, domain=0:5] {3};
\node[poporange, anchor=west, font=\tiny] at (axis cs:5,3) {$y=3$ (Max)};
% Droite y=-6 (min)
\addplot[popgreen, dashed, domain=0:5] {-6};
\node[popgreen, anchor=west, font=\tiny] at (axis cs:5,-6) {$y=-6$ (Min)};
\end{axis}
\end{tikzpicture}
\legende{Courbe $\mathcal{C}_f$ de $f(x)=-x^2+4x-1$ sur $[0;5]$. Sommet $S(2;3)$, intersections $A,B,C$, extremums absolus.}
\end{popfigure}

\courssub{Résolution graphique d'équations et d'inéquations sur $I$}

Une fois la courbe $\mathcal{C}_f$ tracée avec précision sur l'intervalle $I$, elle devient un outil puissant pour résoudre graphiquement des équations et inéquations.

\textbf{Équation $f(x) = 0$.}\par
Les solutions sont les \textbf{abscisses des points d'intersection} de $\mathcal{C}_f$ avec l'axe des abscisses ($Ox$).
Sur l'exemple précédent : $x_1 = 2-\sqrt{3}$ et $x_2 = 2+\sqrt{3}$.

\textbf{Équation $f(x) = m$ (avec $m$ réel).}\par
On trace la droite horizontale $\mathcal{D}_m : y = m$.
Les solutions sont les \textbf{abscisses des points d'intersection} de $\mathcal{C}_f$ avec $\mathcal{D}_m$.
Le nombre de solutions dépend de la position de $m$ par rapport aux extremums de $f$ sur $I$.

\begin{propriete}[{Nombre de solutions de $f(x)=m$ sur $I=[m\,;\,n]$ pour $f$ polynôme degré 2}]
Soit $m_{\min}$ le minimum absolu et $m_{\max}$ le maximum absolu de $f$ sur $I$.
\begin{itemize}
    \item Si $m < m_{\min}$ ou $m > m_{\max}$ : \textbf{0 solution}.
    \item Si $m = m_{\min}$ ou $m = m_{\max}$ : \textbf{1 solution} (si l'extremum est au sommet intérieur) ou \textbf{1 ou 2 solutions} (si l'extremum est à une borne, vérifier l'autre intersection).
    \item Si $m_{\min} < m < m_{\max}$ : \textbf{2 solutions} (généralement), sauf cas particuliers aux bornes.
\end{itemize}
\end{propriete}

\begin{attention}[Lecture graphique vs Calcul algébrique]
La résolution \textbf{graphique} donne des valeurs \textbf{approchées} (lecture sur le graphique). Pour des valeurs exactes, il faut résoudre l'équation algébriquement ($f(x)-m=0$). Au Probatoire, la question « Résoudre graphiquement » attend une lecture sur le graphique avec encadrement (ex: $x \approx \num{1,2}$), tandis que « Résoudre » ou « Déterminer » attend le calcul exact.
\end{attention}

\textbf{Inéquation $f(x) > 0$ (ou $<0$, $\ge 0$, $\le 0$).}\par
On repère sur l'axe des abscisses les portions où la courbe est \textbf{au-dessus} (pour $>0$) ou \textbf{au-dessous} (pour $<0$) de l'axe $Ox$.
\begin{itemize}
    \item $f(x) > 0$ : ensemble des $x \in I$ tels que le point de la courbe a $y>0$ (zone hachurée au-dessus de $Ox$).
    \item $f(x) < 0$ : ensemble des $x \in I$ tels que le point de la courbe a $y<0$ (zone hachurée en-dessous de $Ox$).
\end{itemize}
On écrit la solution sous forme d'intervalle(s) ou d'union d'intervalles, en tenant compte des bornes de $I$ et de l'appartenance des zéros (selon inéquation large ou stricte).

\begin{popfigure}
\begin{tikzpicture}
\begin{axis}[
    width=\linewidth,
    height=8cm,
    axis lines=middle,
    xlabel=$x$, ylabel=$y$,
    xmin=-0.5, xmax=5.5,
    ymin=-7, ymax=4,
    xtick={0,1,2,3,4,5},
    ytick={-6,-4,-2,0,2,3},
    minor tick num=1,
    grid=major,
    grid style={popline, opacity=0.3},
    tick label style={font=\small},
    label style={font=\small},
    clip=false,
    domain=0:5,
    samples=200,
]
% Courbe
\addplot[popcurve, thick, domain=0:5] {-x^2 + 4*x - 1};
% Zone f(x) > 0 (hachures) - racines approx 0.268 et 3.732
\addplot[popgreen!30, fill opacity=0.3, domain=0.268:3.732] {-x^2 + 4*x - 1} \closedcycle;
% Droite y = m (exemple m=1)
\addplot[poporange, thick, dashed, domain=0:5] {1};
\node[poporange, anchor=west, font=\small] at (axis cs:5.1,1) {$\mathcal{D}_1 : y=1$};
% Points intersection y=1
% -x^2+4x-1=1 -> -x^2+4x-2=0 -> x^2-4x+2=0 -> x = 2 +/- sqrt(2) ~ 0.59 et 3.41
\addplot[only marks, mark=*, mark size=2.5pt, poporange] coordinates {(0.586,1) (3.414,1)};
% Points intersection y=0 (racines)
\addplot[only marks, mark=*, mark size=2.5pt, popdark] coordinates {(0.268,0) (3.732,0)};
% Annotations
\node[popgreen, font=\small, rotate=0] at (axis cs:2,1.5) {$f(x) > 0$};
\node[popdark, font=\small] at (axis cs:2,-4) {$f(x) < 0$};
\end{axis}
\end{tikzpicture}
\legende{Résolution graphique : $f(x)=1$ (2 solutions $\approx \num{0,59}$ et $\num{3,41}$) ; $f(x)>0$ (zone verte hachurée entre les racines $2-\sqrt{3}$ et $2+\sqrt{3}$).}
\end{popfigure}

\courssub{Lien avec la forme canonique : Translation de la parabole de référence}

C'est une méthode alternative, souvent plus rapide pour tracer, exigée au programme : utiliser la forme canonique pour identifier la translation qui envoie la parabole de référence $y=ax^2$ sur la parabole $\mathcal{C}_f$.

\begin{propriete}[Forme canonique et translation]
Toute fonction $f(x)=ax^2+bx+c$ s'écrit sous la forme canonique :
\[ f(x) = a(x-\alpha)^2 + \beta \quad \text{avec} \quad \alpha = -\frac{b}{2a},\quad \beta = f(\alpha) = -\frac{\Delta}{4a}. \]
La courbe $\mathcal{C}_f$ est l'image de la parabole de référence $\mathcal{P}_0 : y = ax^2$ par la \textbf{translation} de vecteur $\vec{v}(\alpha\,;\,\beta)$.
Le sommet $S$ de $\mathcal{C}_f$ est l'image de l'origine $O(0,0)$ par cette translation.
\end{propriete}

\begin{savaistu}[Pourquoi la forme canonique est-elle « canonique » ?]
Le terme vient du grec \emph{kanon} (règle, modèle). La forme $a(x-\alpha)^2+\beta$ est la « règle standard » qui exhibe immédiatement le sommet $S(\alpha,\beta)$ et le coefficient directeur $a$ (ouverture et concavité). C'est la forme « normale » pour les translations. Au Cameroun, on la rencontre aussi en Terminale pour l'étude des fonctions polynomiales de degré supérieur (factorisation par $a$ et mise en évidence du sommet).
\end{savaistu}

\begin{experience}[Découverte : De $y=ax^2$ à $f(x)=a(x-\alpha)^2+\beta$ par translation]
Prenons $f(x) = -x^2 + 4x - 1$.
\begin{enumerate}
    \item Forme canonique : $f(x) = -(x^2 - 4x) - 1 = -[(x-2)^2 - 4] - 1 = -(x-2)^2 + 4 - 1 = -(x-2)^2 + 3$.
    \item On a $a=-1$, $\alpha=2$, $\beta=3$.
    \item Parabole de référence : $\mathcal{P}_0 : y = -x^2$ (sommet à l'origine $O$, branches vers le bas).
    \item Translation $t_{\vec{v}}$ de vecteur $\vec{v}(2\,;\,3)$.
    \item Image de $O(0,0)$ : $S(2,3)$.
    \item Image d'un point $P(1,-1)$ de $\mathcal{P}_0$ : $P'(1+2, -1+3) = P'(3,2)$. Vérification : $f(3) = -9+12-1=2$. Correct.
    \item Image de $P(-1,-1)$ : $P'(1,2)$. Vérification : $f(1) = -1+4-1=2$. Correct.
\end{enumerate}
\emph{Conclusion :} Pour tracer $\mathcal{C}_f$, il suffit de tracer $y=ax^2$ (facile : points $(0,0), (\pm 1, a), (\pm 2, 4a)$...) et de décaler tout le tracé de $\alpha$ en abscisse et $\beta$ en ordonnée.
\end{experience}

\begin{methode}[Tracer $\mathcal{C}_f$ par translation (Méthode canonique)]
\begin{enumerate}
    \item Mettre $f$ sous forme canonique $f(x)=a(x-\alpha)^2+\beta$ (ou lire $\alpha, \beta$ via $\alpha=-b/2a, \beta=-\Delta/4a$).
    \item Tracer la parabole de référence $y=ax^2$ dans un repère (point sommet $O$, points $(\pm 1, a), (\pm 2, 4a)$).
    \item Effectuer la translation de vecteur $\vec{v}(\alpha\,;\,\beta)$ : déplacer le repère ou reporter les points en ajoutant $\alpha$ aux abscisses et $\beta$ aux ordonnées.
    \item Ne garder que la portion correspondant à l'intervalle d'étude $I=[m\,;\,n]$.
\end{enumerate}
\end{methode}

\begin{attention}[Piège : Intervalle d'étude et translation]
La translation s'applique sur \textbf{tout} $\mathbb{R}$. Si l'exercice demande l'étude sur $[m\,;\,n]$, on trace la parabole complète par translation, puis on \textbf{tronque} la courbe aux bornes $m$ et $n$ (points pleins si bornes incluses, creux si exclues — ici fermées donc pleins). N'oubliez pas de calculer $f(m)$ et $f(n)$ pour les points extrêmes du tracé.
\end{attention}

\courssub{Exercices résolus et entraînement}

\begin{exoresolu}[{Exercice type Probatoire : $g(x) = 2x^2 - 8x + 5$ sur $I=[1\,;\,4]$}]
\textbf{1. Forme canonique et translation.}
$g(x) = 2(x^2 - 4x) + 5 = 2[(x-2)^2 - 4] + 5 = 2(x-2)^2 - 8 + 5 = 2(x-2)^2 - 3$.
$a=2>0$, $\alpha=2$, $\beta=-3$. Sommet $S(2\,;\,-3)$. Translation de $y=2x^2$ par $\vec{v}(2\,;\,-3)$.

\textbf{2. Dérivée et variations sur $[1\,;\,4]$.}
$g'(x) = 4x - 8 = 4(x-2)$.
$g'(x)=0 \Leftrightarrow x=2 \in [1\,;\,4]$.
Signe de $g'$ : $-$ sur $[1\,;\,2[$, $+$ sur $]2\,;\,4]$.
$g$ décroissante sur $[1\,;\,2]$, croissante sur $[2\,;\,4]$. Minimum en $2$.

\textbf{3. Tableau de variations.}
$g(1) = 2 - 8 + 5 = -1$.
$g(2) = -3$.
$g(4) = 32 - 32 + 5 = 5$.

\begin{center}
\begin{tabular}{|c|ccccc|}\hline
$x$ & $1$ & & $2$ & & $4$ \\ \hline
$g'(x)$ & & $-$ & $0$ & $+$ & \\ \hline
$g(x)$ & $-1$ & $\searrow$ & $-3$ & $\nearrow$ & $5$ \\ \hline
\end{tabular}
\end{center}

\textbf{4. Extremums absolus sur $[1\,;\,4]$.}
Min absolu : $-3$ (en $x=2$). Max absolu : $5$ (en $x=4$).

\textbf{5. Intersections.}
$Oy$ : $0 \notin I$, pas d'intersection dans l'intervalle.
$Ox$ : $\Delta = 64 - 40 = 24$. $\sqrt{\Delta}=2\sqrt{6}$.
$x_1 = \frac{8-2\sqrt{6}}{4} = 2 - \frac{\sqrt{6}}{2} \approx \num{0,77} \notin I$.
$x_2 = 2 + \frac{\sqrt{6}}{2} \approx \num{3,22} \in I$.
Un seul zéro dans $I$ : $x_0 = 2 + \frac{\sqrt{6}}{2}$.

\textbf{6. Résolution graphique sur $I$ :}
\begin{itemize}
    \item $g(x)=0 \Rightarrow x \approx \num{3,22}$.
    \item $g(x) < 0 \Rightarrow x \in [1\,;\, 2+\frac{\sqrt{6}}{2}[$ (courbe sous $Ox$).
    \item $g(x) = 2 \Rightarrow$ Droite $y=2$. $2 \in ]-3\,;\,5[$ $\Rightarrow$ 2 solutions. 
          $2(x-2)^2-3=2 \Leftrightarrow (x-2)^2=2,5 \Leftrightarrow x=2\pm\sqrt{2,5}$. 
          $2-\sqrt{2,5} \approx \num{0,42} \notin I$. $2+\sqrt{2,5} \approx \num{3,58} \in I$.
          \textbf{Une seule solution dans $I$ : $x \approx \num{3,58}$}.
\end{itemize}
\tcblower
\textbf{Solution détaillée.}
Voir les étapes ci-dessus. L'analyse du nombre de solutions pour $g(x)=m$ montre l'importance de croiser l'intervalle d'étude $I$ avec les solutions algébriques. Ici, pour $m=2$, une seule des deux solutions algébriques tombe dans $[1\,;\,4]$.
\end{exoresolu}

\begin{exercice}[{Entraînement : $h(x) = -2x^2 + 4x + 6$ sur $I=[-1\,;\,3]$}]
\begin{enumerate}
    \item Écrire $h$ sous forme canonique. En déduire le vecteur de translation depuis $y=-2x^2$.
    \item Étudier les variations de $h$ sur $I$ et dresser le tableau de variations.
    \item Déterminer les extremums absolus de $h$ sur $I$.
    \item Tracer la courbe $\mathcal{C}_h$ dans un repère (échelle : 2 cm par unité).
    \item Résoudre graphiquement sur $I$ : $h(x)=0$ ; $h(x)=6$ ; $h(x) > 0$.
\end{enumerate}
\end{exercice}

\begin{corrige}[Exercice Entraînement]
\textbf{1.} $h(x) = -2(x^2-2x) + 6 = -2[(x-1)^2-1] + 6 = -2(x-1)^2 + 2 + 6 = -2(x-1)^2 + 8$.
$a=-2$, $\alpha=1$, $\beta=8$. Sommet $S(1,8)$. Translation de $y=-2x^2$ par $\vec{v}(1\,;\,8)$.

\textbf{2.} $h'(x) = -4x+4 = -4(x-1)$. $h'(x)=0 \Leftrightarrow x=1 \in I$.
$a<0 \Rightarrow h'$ $+$ puis $-$ $\Rightarrow$ $h$ croissante puis décroissante. Max en $1$.
$h(-1) = -2-4+6=0$. $h(1)=8$. $h(3) = -18+12+6=0$.
Tableau :
\begin{center}
\begin{tabular}{|c|ccccc|}\hline
$x$ & $-1$ & & $1$ & & $3$ \\ \hline
$h'(x)$ & & $+$ & $0$ & $-$ & \\ \hline
$h(x)$ & $0$ & $\nearrow$ & $8$ & $\searrow$ & $0$ \\ \hline
\end{tabular}
\end{center}

\textbf{3.} Min absolu : $0$ (atteint en $-1$ et $3$, aux bornes). Max absolu : $8$ (atteint en $1$, au sommet).

\textbf{4.} Courbe : Parabole tournée vers le bas, sommet $S(1,8)$, passant par $A(-1,0)$ et $B(3,0)$ (racines doubles aux bornes ! $\Delta=16+48=64$, racines $-1$ et $3$). Concavité $a=-2$.

\textbf{5.} 
\begin{itemize}
    \item $h(x)=0 \Rightarrow x=-1$ et $x=3$ (les bornes).
    \item $h(x)=6 \Rightarrow -2(x-1)^2+8=6 \Rightarrow (x-1)^2=1 \Rightarrow x=0$ ou $x=2$. Les deux dans $I$.
    \item $h(x)>0 \Rightarrow$ Courbe au-dessus de $Ox$. Comme les zéros sont aux bornes et le sommet au-dessus, $h(x)>0$ sur $]-1\,;\,3[$.
\end{itemize}
\end{corrige}

\begin{aretenir}[{Check-list pour l'étude complète sur $I=[m\,;\,n]$ (Polynôme degré 2)}]
$\square$ 1. Forme canonique $a(x-\alpha)^2+\beta$ (sommet $S$, translation $\vec{v}(\alpha,\beta)$).
$\square$ 2. Dérivée $f'(x)=2ax+b$, zéro $\alpha$, signe de $f'$ sur $I$.
$\square$ 3. Tableau de variations sur $I$ (bornes $m,n$ + $\alpha$ si intérieur).
$\square$ 4. Extremums absolus : comparer $f(m), f(n), f(\alpha)$ (si $\alpha\in I$).
$\square$ 5. Intersections axes : $f(0)$ si $0\in I$ ; racines de $f(x)=0$ dans $I$.
$\square$ 6. Tracé courbe : points caractéristiques + concavité ($a$) + restriction à $I$.
$\square$ 7. Résolution graphique : $f(x)=m$ (droite $y=m$) ; $f(x)>0$ (au-dessus $Ox$).
\end{aretenir}

\coursec{Étude complète d'une fonction homographique sur un intervalle borné}

Après avoir maîtrisé l'étude des fonctions polynomiales du second degré, nous abordons maintenant les fonctions homographiques. Ces fonctions, quotients de deux polynômes du premier degré, possèdent une géométrie fascinante : leurs courbes sont des \textbf{hyperboles} dont les asymptotes guident le tracé. Comme pour les polynômes, nous travaillerons sur un \textbf{intervalle borné} ne contenant pas la valeur interdite, ce qui garantit l'existence des extrema et simplifie le tableau de variations. Vous verrez que la méthode est systématique : réécriture, asymptotes, dérivée, variations, courbe.

\courssub{Réécriture canonique par division euclidienne : la clé de tout}

Toute fonction homographique s'écrit $f(x) = \dfrac{ax+b}{cx+d}$ avec $c \neq 0$ (sinon c'est une fonction affine) et $ad-bc \neq 0$ (sinon c'est une fonction constante). La première étape \textbf{obligatoire} est de la mettre sous la forme $f(x) = \alpha + \dfrac{k}{x-\beta}$. Cette forme révèle immédiatement les asymptotes, le centre de symétrie et le sens de variation.

\begin{methode}[Division euclidienne pour obtenir la forme $\alpha + \frac{k}{x-\beta}$]
Soit $f(x) = \dfrac{ax+b}{cx+d}$ ($c \neq 0$).
\begin{enumerate}
    \item Diviser le numérateur par le dénominateur (division euclidienne de polynômes de degré 1).
    \item Le \textbf{quotient} est le nombre $\alpha = \dfrac{a}{c}$ (c'est l'ordonnée de l'asymptote horizontale).
    \item Le \textbf{reste} est un nombre $R = b - \alpha d = b - \frac{a}{c}d = \frac{bc-ad}{c}$.
    \item On factorise le dénominateur : $cx+d = c(x + \frac{d}{c}) = c(x-\beta)$ avec $\beta = -\frac{d}{c}$ (c'est l'abscisse de l'asymptote verticale, la valeur interdite).
    \item On obtient : $f(x) = \alpha + \dfrac{R}{c(x-\beta)} = \alpha + \dfrac{k}{x-\beta}$ où \textbf{$k = \dfrac{R}{c} = \dfrac{bc-ad}{c^2}$}.
\end{enumerate}
\textbf{Astuce mémoire :} $\alpha = \frac{a}{c}$, $\beta = -\frac{d}{c}$, $k = \frac{bc-ad}{c^2} = \frac{\det\begin{pmatrix}a & b \\ c & d\end{pmatrix}}{c^2}$.
\end{methode}

\begin{exemplebox}[Mise en forme canonique]
Mettons $f(x) = \dfrac{2x+1}{x-3}$ sous forme canonique.
Ici $a=2, b=1, c=1, d=-3$.
\begin{itemize}
    \item $\alpha = \frac{a}{c} = 2$.
    \item $\beta = -\frac{d}{c} = 3$. (Valeur interdite : $x \neq 3$).
    \item $k = \frac{bc-ad}{c^2} = \frac{1\cdot1 - 2\cdot(-3)}{1^2} = \frac{1+6}{1} = 7$.
\end{itemize}
On vérifie : $f(x) = 2 + \dfrac{7}{x-3}$.
\emph{Vérification rapide :} $2 + \frac{7}{x-3} = \frac{2(x-3)+7}{x-3} = \frac{2x-6+7}{x-3} = \frac{2x+1}{x-3}$. C'est bon !
\end{exemplebox}

\courssub{Asymptotes et centre de symétrie : la géométrie de l'hyperbole}

La forme $f(x) = \alpha + \frac{k}{x-\beta}$ est une mine d'informations géométriques.

\begin{definition}[Asymptotes d'une fonction homographique]
Soit $f(x) = \alpha + \frac{k}{x-\beta}$ définie sur $\mathbb{R}^*_\beta = \mathbb{R} \setminus \{\beta\}$.
\begin{itemize}
    \item La droite \textbf{verticale} d'équation $\mathcal{V} : x = \beta$ est une \textbf{asymptote verticale (AV)}.
      $\lim\limits_{x \to \beta^-} f(x) = \operatorname{signe}(k) \cdot (-\infty)$ et $\lim\limits_{x \to \beta^+} f(x) = \operatorname{signe}(k) \cdot (+\infty)$.
    \item La droite \textbf{horizontale} d'équation $\mathcal{H} : y = \alpha$ est une \textbf{asymptote horizontale (AH)}.
      $\lim\limits_{x \to -\infty} f(x) = \alpha$ et $\lim\limits_{x \to +\infty} f(x) = \alpha$.
\end{itemize}
\end{definition}

\begin{propriete}[Centre de symétrie]
Les deux asymptotes se coupent en le point $I(\beta\,;\, \alpha)$. Ce point est le \textbf{centre de symétrie} de la courbe $\mathcal{C}_f$.
La courbe est invariante par la symétrie centrale de centre $I$ (rotation d'angle $\pi$ autour de $I$).
\emph{Preuve (exigible en rédaction) :} Pour tout $x \neq \beta$, le point symétrique de $M(x\,;\, f(x))$ par rapport à $I$ est $M'(2\beta-x\,;\, 2\alpha-f(x))$.
On vérifie que $f(2\beta-x) = \alpha + \frac{k}{2\beta-x-\beta} = \alpha + \frac{k}{\beta-x} = \alpha - \frac{k}{x-\beta} = 2\alpha - f(x)$.
Donc $M'$ appartient bien à la courbe.
\end{propriete}

\begin{attention}[Piège : Domaine de définition et intervalle d'étude]
Au Probatoire, l'étude se fait \textbf{toujours sur un intervalle borné $[m, n]$ NE CONTENANT PAS $\beta$}.
Si $\beta \in [m, n]$, la fonction n'est pas définie sur tout l'intervalle (elle « explose » en $\beta$). On ne dresse \textbf{qu'un seul tableau de variations} sur un intervalle où $f$ est continue et dérivable. Vérifiez toujours que $\beta \notin [m, n]$ avant de commencer le tableau !
\end{attention}

\courssub{Dérivée, sens de variation et tableau de variations}

La forme canonique rend le calcul de la dérivée immédiat.

\begin{propriete}[Dérivée et variations]
Soit $f(x) = \alpha + \frac{k}{x-\beta}$ sur un intervalle $I$ ne contenant pas $\beta$.
\[ f'(x) = -\frac{k}{(x-\beta)^2} \]
Le dénominateur $(x-\beta)^2$ est \textbf{strictement positif} sur $I$.
Le signe de $f'(x)$ est donc \textbf{l'opposé du signe de $k$} sur tout l'intervalle $I$.
\begin{itemize}
    \item Si $k > 0$ : $f'(x) < 0$ $\Rightarrow$ $f$ est \textbf{strictement décroissante} sur $I$.
    \item Si $k < 0$ : $f'(x) > 0$ $\Rightarrow$ $f$ est \textbf{strictement croissante} sur $I$.
\end{itemize}
Il n'y a \textbf{jamais d'extremum relatif} à l'intérieur de l'intervalle (pas de $f'(x)=0$). Les extrema (min et max) sont atteints aux \textbf{bords} de l'intervalle $[m, n]$.
\end{propriete}

\begin{exemplebox}[{Tableau de variations sur $[4, 6]$ pour $f(x) = 2 + \frac{7}{x-3}$}]
Rappel : $\alpha=2, \beta=3, k=7 > 0$. Intervalle $[4, 6]$. $\beta=3 \notin [4, 6]$ : c'est valide.
$f'(x) = -\frac{7}{(x-3)^2} < 0$ sur $[4, 6]$. $f$ est strictement décroissante.
Valeurs aux bornes :
$f(4) = 2 + \frac{7}{1} = 9$.
$f(6) = 2 + \frac{7}{3} = \frac{13}{3} \approx 4,33$.

\begin{center}
\textbf{Tableau de variations sur $[4, 6]$ :}
\begin{tabular}{|c|ccc|}\hline
$x$ & $4$ & & $6$ \\ \hline
$f'(x)$ & & $-$ & \\ \hline
$f$ & $9$ & $\searrow$ & $\frac{13}{3}$ \\ \hline
\end{tabular}
\end{center}
\emph{Note :} Comme la fonction est strictement décroissante, le maximum est $9$ (en $x=4$) et le minimum est $\frac{13}{3}$ (en $x=6$).
\end{exemplebox}

\courssub{Construction de la courbe $\mathcal{C}_f$ par translation}

C'est une compétence clé du programme : passer de la courbe de référence $y = \frac{k}{x}$ à celle de $f$ par une translation.

\begin{propriete}[Construction par translation]
La courbe $\mathcal{C}_f$ de $f(x) = \alpha + \frac{k}{x-\beta}$ s'obtient par translation de la courbe $\mathcal{C}_0$ d'équation $y = \frac{k}{x}$ (hyperbole de centre $O$, asymptotes les axes) par le \textbf{vecteur $\vec{v}(\beta\,;\, \alpha)$}.
\begin{itemize}
    \item Le centre $O(0,0)$ de $\mathcal{C}_0$ devient le centre $I(\beta, \alpha)$ de $\mathcal{C}_f$.
    \item Les asymptotes $x=0$ et $y=0$ deviennent $x=\beta$ et $y=\alpha$.
    \item Les branches gardent leur forme (convexité) et leur position relative par rapport aux asymptotes.
\end{itemize}
\end{propriete}

\begin{experience}[Découverte : Translation de $y = 7/x$ vers $y = 2 + 7/(x-3)$]
Prenons la courbe $\mathcal{C}_0 : y = \frac{7}{x}$ (branches dans les quadrants 1 et 3 car $k=7>0$).
Appliquons la translation de vecteur $\vec{v}(3\,;\, 2)$.
\begin{enumerate}
    \item Le point $A(1, 7)$ de $\mathcal{C}_0$ va en $A'(1+3, 7+2) = A'(4, 9)$. C'est bien $f(4)=9$.
    \item Le point $B(3, 7/3)$ va en $B'(6, 7/3+2) = B'(6, 13/3)$. C'est bien $f(6)=13/3$.
    \item L'asymptote verticale $x=0$ devient $x=3$.
    \item L'asymptote horizontale $y=0$ devient $y=2$.
\end{enumerate}
C'est la méthode la plus rapide et la plus sûre pour tracer la courbe proprement.
\end{experience}

\courssub{Intersections avec les axes et résolution graphique}

\begin{itemize}
    \item \textbf{Axe des ordonnées ($x=0$) :} Si $0 \in [m, n]$, $f(0) = \frac{b}{d}$ (ou $\alpha - \frac{k}{\beta}$). Point $A(0, f(0))$.
    \item \textbf{Axe des abscisses ($y=0$) :} Résoudre $f(x)=0 \Leftrightarrow \alpha + \frac{k}{x-\beta} = 0 \Leftrightarrow \frac{k}{x-\beta} = -\alpha \Leftrightarrow x-\beta = -\frac{k}{\alpha}$ (si $\alpha \neq 0$).
      Soit $x = \beta - \frac{k}{\alpha}$. Vérifier si cette solution appartient à $[m, n]$.
      \emph{Remarque :} C'est exactement la résolution du numérateur $ax+b=0$ : $x = -\frac{b}{a}$.
\end{itemize}

\begin{propriete}[{Résolution graphique de $f(x) = m$ et $f(x) > 0$ sur $[m, n]$}]
\begin{itemize}
    \item \textbf{Équation $f(x) = m$ :} Tracer la droite horizontale $y = m$. Les abscisses des points d'intersection avec $\mathcal{C}_f$ (sur la portion tracée) sont les solutions. Le nombre de solutions est $0, 1$ ou $2$ (une par branche au maximum, mais ici une seule branche est visible sur $[m,n]$).
    \item \textbf{Inéquation $f(x) > 0$ :} Repérer sur l'intervalle $[m, n]$ les portions de courbe \textbf{strictement au-dessus} de l'axe des abscisses (axe $x$).
\end{itemize}
\end{propriete}

\courssub{Exemple complet d'étude : $f(x) = \frac{3x-2}{x+1}$ sur $[0, 5]$}

Appliquons la méthode complète sur cet exemple type Probatoire.

\begin{exoresolu}[{Étude complète de $f(x) = \frac{3x-2}{x+1}$ sur $[0, 5]$}]
\textbf{Énoncé :} On considère la fonction $f$ définie sur $[0, 5]$ par $f(x) = \frac{3x-2}{x+1}$.
\begin{enumerate}
    \item Mettre $f$ sous la forme $\alpha + \frac{k}{x-\beta}$.
    \item Déterminer les asymptotes et le centre de symétrie.
    \item Étudier les variations de $f$ sur $[0, 5]$ et dresser son tableau de variations.
    \item Déterminer les intersections avec les axes.
    \item Résoudre graphiquement $f(x) = 2$ et $f(x) > 0$ sur $[0, 5]$.
\end{enumerate}

\textbf{Solution détaillée :}

\noindent\textbf{1. Forme canonique.}
$a=3, b=-2, c=1, d=1$.
$\alpha = \frac{a}{c} = 3$.
$\beta = -\frac{d}{c} = -1$. (Valeur interdite $x \neq -1$).
$k = \frac{bc-ad}{c^2} = \frac{(-2)(1) - (3)(1)}{1} = -5$.
$f(x) = 3 - \frac{5}{x+1} = 3 + \frac{-5}{x-(-1)}$. \quad \checkmark\ ($3 - \frac{5}{x+1} = \frac{3x+3-5}{x+1} = \frac{3x-2}{x+1}$).

\noindent\textbf{2. Asymptotes et centre.}
\begin{itemize}
    \item Asymptote verticale (AV) : $x = -1$. (Hors intervalle $[0, 5]$, pas de problème).
    \item Asymptote horizontale (AH) : $y = 3$.
    \item Centre de symétrie : $I(-1\,;\, 3)$.
\end{itemize}

\noindent\textbf{3. Variations sur $[0, 5]$.}
$k = -5 < 0 \Rightarrow f'(x) = \frac{5}{(x+1)^2} > 0$ sur $[0, 5]$.
$f$ est \textbf{strictement croissante} sur $[0, 5]$.
Valeurs aux bornes :
$f(0) = \frac{-2}{1} = -2$.
$f(5) = \frac{15-2}{6} = \frac{13}{6} \approx 2,17$.

\begin{center}
\textbf{Tableau de variations sur $[0, 5]$ :}
\begin{tabular}{|c|ccc|}\hline
$x$ & $0$ & & $5$ \\ \hline
$f'(x)$ & & $+$ & \\ \hline
$f$ & $-2$ & $\nearrow$ & $\frac{13}{6}$ \\ \hline
\end{tabular}
\end{center}
Minimum $m = -2$ (en $0$), Maximum $M = \frac{13}{6}$ (en $5$).

\noindent\textbf{4. Intersections avec les axes.}
\begin{itemize}
    \item Axe des ordonnées ($x=0 \in [0,5]$) : $A(0\,;\, -2)$.
    \item Axe des abscisses ($y=0$) : $3x-2=0 \Leftrightarrow x = \frac{2}{3}$.
      $\frac{2}{3} \in [0, 5]$. Point $B(\frac{2}{3}\,;\, 0)$.
\end{itemize}

\noindent\textbf{5. Résolution graphique.}
\begin{itemize}
    \item $f(x) = 2$ : Droite $y=2$. La courbe est croissante, passe par $A(0,-2)$ et $B(5, 13/6 \approx 2,17)$. Elle coupe $y=2$ une fois.
      Calcul exact : $3 - \frac{5}{x+1} = 2 \Leftrightarrow \frac{5}{x+1} = 1 \Leftrightarrow x+1 = 5 \Leftrightarrow x = 4$.
      Solution unique : $S = \{4\}$.
    \item $f(x) > 0$ : La courbe est au-dessus de l'axe $Ox$ pour $x > \frac{2}{3}$.
      Sur $[0, 5]$, l'ensemble des solutions est $]\frac{2}{3}, 5]$.
\end{itemize}
\end{exoresolu}

\begin{popfigure}
\begin{tikzpicture}[scale=0.8, >=Stealth]
% Axes
\draw[->] (-2,0) -- (6.5,0) node[right] {$x$};
\draw[->] (0,-3) -- (0,4.5) node[above] {$y$};
% Graduations
\foreach \x in {-1,1,2,3,4,5,6} \draw (\x,0.1) -- (\x,-0.1) node[below] {\small $\x$};
\foreach \y in {-2,-1,1,2,3,4} \draw (0.1,\y) -- (-0.1,\y) node[left] {\small $\y$};
% Asymptotes
\draw[popmuted, dashed, thick] (-1,-3) -- (-1,4.5) node[above] {$x=-1$ (AV)};
\draw[popmuted, dashed, thick] (-2,3) -- (6.5,3) node[right] {$y=3$ (AH)};
% Centre I
\fill[poppurple] (-1,3) circle (2.5pt) node[above left] {$I(-1,3)$};
% Courbe sur [0,5] (branche droite de l'hyperbole)
\draw[popblue, very thick, domain=0:5, smooth, variable=\x, samples=100] plot ({\x}, {3 - 5/(\x+1)});
% Points caractéristiques
\fill[popblue] (0,-2) circle (2pt) node[below left] {$A(0,-2)$};
\fill[popblue] (0.666,0) circle (2pt) node[below right] {$B(2/3,0)$};
\fill[popblue] (4,2) circle (2pt) node[above right] {$M(4,2)$};
\fill[popblue] (5,2.166) circle (2pt) node[above right] {$N(5,13/6)$};
% Droite y=2 pour f(x)=2
\draw[poporange, dashed] (-1,2) -- (6,2) node[right] {$y=2$};
% Zone f(x)>0
\draw[popgreen!50!black, thick, domain=0.666:5, smooth, variable=\x, samples=50] plot ({\x}, {3 - 5/(\x+1)});
\end{tikzpicture}
\legende{Courbe de $f(x)=\frac{3x-2}{x+1}$ sur $[0,5]$. Asymptotes pointillées, centre $I$, intersections axes, résolution $f(x)=2$ (point $M$) et $f(x)>0$ (portion verte).}
\end{popfigure}

\courssub{Cas particulier : Fonction homographique et translation de $y = k/x$ (Savoir-faire programme)}

Le programme insiste sur l'écriture $f(x) = \alpha + \frac{k}{x-\beta}$ pour identifier le vecteur de translation $\vec{v}(\beta, \alpha)$ depuis la courbe de référence $y = \frac{k}{x}$ (ou $y = \frac{k}{x}$ si on part de $x \mapsto k/x$). Attention : la courbe de référence du programme est souvent $x \mapsto kx$ pour les affines, mais pour les homographiques, c'est bien $x \mapsto \frac{k}{x}$ (ou $x \mapsto \frac{1}{x}$ puis homothétie de rapport $k$). Le programme dit : « Utiliser l'écriture $x \mapsto \alpha + \frac{k}{x-\beta}$ pour déterminer le vecteur de la translation permettant de construire la courbe d'une fonction homographique à partir de celle de la fonction $x \mapsto \frac{k}{x}$ ».

\begin{exemplebox}[Identification du vecteur de translation]
Soit $f(x) = \frac{4x+5}{2x-1}$.
Forme canonique : $\alpha = \frac{4}{2} = 2$, $\beta = -\frac{-1}{2} = 0,5$, $k = \frac{5\cdot2 - 4\cdot(-1)}{4} = \frac{10+4}{4} = \frac{14}{4} = 3,5$.
$f(x) = 2 + \frac{3,5}{x-0,5}$.
La courbe $\mathcal{C}_f$ s'obtient par translation de la courbe $\mathcal{C}_0 : y = \frac{3,5}{x}$ par le vecteur $\vec{v}(0,5\,;\, 2)$.
Le centre $O(0,0)$ va en $I(0,5\,;\, 2)$.
\end{exemplebox}

\begin{savaistu}[Le déterminant caché : $k = \frac{\det}{c^2}$]
Le coefficient $k = \frac{bc-ad}{c^2}$ est (à un facteur près) le \textbf{déterminant} de la matrice associée à l'homographie $\begin{pmatrix} a & b \\ c & d \end{pmatrix}$.
\begin{itemize}
    \item Si $\det = 0$ ($ad=bc$), alors $k=0$ et $f(x) = \alpha$ constante. La « courbe » se réduit à une droite horizontale (cas dégénéré).
    \item Le signe de $k$ (donc du déterminant) dicte la position des branches par rapport aux asymptotes :
    \begin{itemize}
        \item $k > 0$ (det $> 0$) : Branches dans les quadrants « haut-droite » et « bas-gauche » par rapport à $I$ (comme $y=1/x$). Fonction décroissante.
        \item $k < 0$ (det $< 0$) : Branches dans les quadrants « haut-gauche » et « bas-droite » par rapport à $I$ (comme $y=-1/x$). Fonction croissante.
    \end{itemize}
    C'est une belle unification algébrique-géométrique !
\end{itemize}
\end{savaistu}

\begin{attention}[Erreurs fatales à éviter au Probatoire]
\begin{enumerate}
    \item \textbf{Oublier la division euclidienne :} Tenter d'étudier le signe de $f'(x)$ sur la forme $\frac{ax+b}{cx+d}$ sans la simplifier en $\alpha + \frac{k}{x-\beta}$. C'est possible mais beaucoup plus long et source d'erreurs de signe.
    \item \textbf{Confondre $\beta$ et $-\beta$ :} $\beta = -\frac{d}{c}$. L'asymptote verticale est $x = \beta$. Si $f(x) = \frac{2x+1}{x-3}$, $d=-3$, $\beta = -(-3)/1 = 3$. AV : $x=3$. Ne pas écrire $x=-3$ !
    \item \textbf{Faire un tableau de variations sur un intervalle contenant $\beta$ :} Ex: $f(x)=\frac{1}{x-2}$ sur $[0, 4]$. $\beta=2 \in [0,4]$. Interdit de faire un seul tableau. (Hors programme 1ère A, on vous donnera un intervalle ne contenant pas $\beta$, ex $[0,1]$ ou $[3,4]$).
    \item \textbf{Oublier de vérifier que les solutions (équations, inéquations) appartiennent à l'intervalle d'étude $[m, n]$ :} Une solution $x = -5$ est refusée si on étudie sur $[0, 10]$.
    \item \textbf{Confondre « croissante » et « décroissante » par rapport au signe de $k$ :} $k>0 \Rightarrow$ Décroissante. $k<0 \Rightarrow$ Croissante. Répétez-le : « $k$ positif, fonction négative (décroissante) ».
\end{enumerate}
\end{attention}

\begin{exercice}[Entraînement : Étude autonome]
On considère la fonction $g$ définie sur $[2, 8]$ par $g(x) = \frac{5x-4}{x-1}$.
\begin{enumerate}
    \item Écrire $g$ sous la forme $\alpha + \frac{k}{x-\beta}$. Préciser $\alpha, \beta, k$.
    \item Donner les équations des asymptotes et les coordonnées du centre de symétrie $I$.
    \item Étudier le sens de variation de $g$ sur $[2, 8]$. Dresser le tableau de variations.
    \item Déterminer les intersections de la courbe avec les axes (si elles existent dans l'intervalle).
    \item Résoudre graphiquement (et vérifier par le calcul) l'équation $g(x) = 4$ sur $[2, 8]$.
    \item Résoudre l'inéquation $g(x) \leq 5$ sur $[2, 8]$.
\end{enumerate}
\emph{Indications :} $\beta=1 \notin [2,8]$. $k = \frac{-4\cdot1 - 5\cdot(-1)}{1} = 1 > 0$. Donc décroissante. $g(2)=6$, $g(8)=\frac{36}{7}\approx 5,14$.
\end{exercice}

\begin{corrige}[Exercice n°1]
\begin{enumerate}
    \item $g(x) = 5 + \frac{1}{x-1}$. $\alpha=5, \beta=1, k=1$.
    \item AV : $x=1$. AH : $y=5$. Centre $I(1, 5)$.
    \item $k=1>0 \Rightarrow g$ strictement décroissante sur $[2, 8]$.
    Tableau :
    \begin{center}
    \begin{tabular}{|c|ccc|}\hline
    $x$ & $2$ & & $8$ \\ \hline
    $g'(x)$ & & $-$ & \\ \hline
    $g$ & $6$ & $\searrow$ & $\frac{36}{7}$ \\ \hline
    \end{tabular}
    \end{center}
    \item Axe $Oy$ : $x=0 \notin [2,8]$ $\rightarrow$ pas d'intersection dans l'intervalle.
    Axe $Ox$ : $5x-4=0 \Rightarrow x=0,8 \notin [2,8]$ $\rightarrow$ pas d'intersection.
    \item $g(x)=4 \Leftrightarrow 5 + \frac{1}{x-1} = 4 \Leftrightarrow \frac{1}{x-1} = -1 \Leftrightarrow x-1 = -1 \Leftrightarrow x=0$.
    $0 \notin [2, 8]$. \textbf{Aucune solution} dans l'intervalle d'étude. (Graphiquement, la droite $y=4$ est sous la courbe qui varie de $6$ à $5,14$).
    \item $g(x) \leq 5 \Leftrightarrow 5 + \frac{1}{x-1} \leq 5 \Leftrightarrow \frac{1}{x-1} \leq 0$.
    Sur $[2, 8]$, $x-1 > 0$, donc $\frac{1}{x-1} > 0$. L'inéquation n'a \textbf{aucune solution} sur $[2, 8]$.
    (La courbe est toujours strictement au-dessus de l'asymptote horizontale $y=5$ car $k>0$).
\end{enumerate}
\end{corrige}

\courssub{Synthèse de la procédure d'étude (Fiche méthode)}

Pour briller à l'examen, suivez ce rituel immuable pour toute fonction homographique $f(x) = \frac{ax+b}{cx+d}$ sur $[m, n]$ :

\begin{methode}[Rituel « Homographique sur intervalle borné » - 6 étapes]
\begin{enumerate}
    \item \textbf{Domaine / Forme canonique :} Calculer $\alpha = a/c$, $\beta = -d/c$, $k = (bc-ad)/c^2$. Écrire $f(x) = \alpha + \frac{k}{x-\beta}$. Vérifier que $\beta \notin [m, n]$.
    \item \textbf{Asymptotes / Centre :} AV : $x=\beta$. AH : $y=\alpha$. Centre $I(\beta, \alpha)$. Tracer les asymptotes en pointillés.
    \item \textbf{Dérivée / Variations :} $f'(x) = -k/(x-\beta)^2$. Signe de $f'$ = opposé de signe($k$).
        \begin{itemize}
            \item $k>0 \Rightarrow$ Décroissante.
            \item $k<0 \Rightarrow$ Croissante.
        \end{itemize}
        Dresser le tableau de variations sur $[m, n]$ avec $f(m)$ et $f(n)$.
    \item \textbf{Intersections :} $A(0, f(0))$ si $0 \in [m, n]$. $B(x_0, 0)$ avec $x_0 = -b/a$ (ou $\beta - k/\alpha$) si $x_0 \in [m, n]$.
    \item \textbf{Construction de la courbe :} Tracer le repère, asymptotes, centre $I$, points $A, B, (m, f(m)), (n, f(n))$. Tracer la branche unique (hyperbole) approchant les asymptotes. Utiliser la translation de $y=k/x$ par $\vec{v}(\beta, \alpha)$ pour guider la main.
    \item \textbf{Résolution graphique :} $f(x)=m$ $\rightarrow$ intersection avec $y=m$. $f(x)>0$ $\rightarrow$ portion au-dessus de $Ox$. Lire les abscisses sur l'axe $x$.
\end{enumerate}
\end{methode}

\begin{savaistu}[Application concrète : Modélisation d'un coût moyen]
Au Cameroun, une usine de transformation de cacao à Douala a un coût fixe mensuel de $F = 500\,000$ FCFA (loyer, salaires fixes) et un coût variable de $v = 200$ FCFA par kg de cacao traité. Le coût total pour $x$ kg est $C_T(x) = 500\,000 + 200x$.
Le \textbf{coût moyen} par kg est une fonction homographique :
\[ C_m(x) = \frac{C_T(x)}{x} = \frac{200x + 500\,000}{x} = 200 + \frac{500\,000}{x} \]
Ici $\alpha = 200$ (coût variable, limite inférieure), $\beta = 0$, $k = 500\,000 > 0$.
La courbe est décroissante : plus on produit, plus le coût moyen baisse (économies d'échelle), mais il ne descendra jamais sous $200$ FCFA/kg (asymptote horizontale).
Si l'usine a une capacité max de $10\,000$ kg/mois, on étudie $C_m$ sur $[1000, 10\,000]$ (on évite $0$).
$C_m(1000) = 700$ FCFA/kg. $C_m(10\,000) = 250$ FCFA/kg.
C'est la puissance des fonctions homographiques en économie réelle !
\end{savaistu}

Vous maîtrisez maintenant l'étude complète des fonctions homographiques sur intervalle borné. La prochaine étape (Section 4) sera passionnante : nous verrons comment, à partir de la courbe d'une fonction $f$ déjà connue, construire les courbes de $-f$, $|f|$, $f(x-a)$, $f(x)+b$... sans refaire tout le calcul ! C'est la magie des \textbf{fonctions associées}.

\coursec{Fonctions associées : Transformations géométriques des courbes}

Après avoir maîtrisé l'étude complète des fonctions polynômes du second degré et homographiques sur des intervalles bornés (Sections 2 et 3), nous disposons désormais d'une « boîte à outils » graphique puissante. Mais en mathématiques, comme dans la construction d'une maison à Yaoundé ou Douala, on gagne un temps précieux en réutilisant des plans types plutôt qu'en reconstruisant tout depuis les fondations. C'est tout l'objet de cette section : comprendre comment modifier une courbe connue $\mathcal{C}_f$ pour obtenir la courbe d'une fonction « cousine » sans refaire tout le travail de dérivation et de tableau de variations.

\courssub{Rappel des transformations élémentaires du plan}

Avant d'agir sur les fonctions, rappelons les transformations géométriques du plan qui seront nos « briques de base ». Soit un repère orthonormé $(O;\vec{\imath},\vec{\jmath})$.

\begin{definition}[Translation]
La translation de vecteur $\vec{v}(a;b)$ est l'application qui à tout point $M(x;y)$ associe le point $M'(x';y')$ tel que $\overrightarrow{MM'}=\vec{v}$. Ses équations cartésiennes sont :
\[
\begin{cases}
x' = x + a \\
y' = y + b
\end{cases}
\quad\Leftrightarrow\quad
\begin{cases}
x = x' - a \\
y = y' - b
\end{cases}
\]
\end{definition}

\begin{definition}[Symétrie axiale par rapport à l'axe des abscisses (Ox)]
C'est l'application qui à $M(x;y)$ associe $M'(x;-y)$. L'axe $Ox$ est l'axe de symétrie ; les points de $Ox$ sont invariants.
\end{definition}

\begin{definition}[Symétrie axiale par rapport à l'axe des ordonnées (Oy)]
C'est l'application qui à $M(x;y)$ associe $M'(-x;y)$. L'axe $Oy$ est l'axe de symétrie ; les points de $Oy$ sont invariants.
\end{definition}

\courssub{Effet sur la courbe $\mathcal{C}_f$ : Les quatre transformations fondamentales}

Soit $f$ une fonction définie sur un intervalle $I$ et $\mathcal{C}_f$ sa courbe représentative dans un repère orthonormé. Nous étudions l'effet de quatre modifications algébriques sur la courbe.

\begin{propriete}[Translation verticale : $x \mapsto f(x) + b$]
La courbe $\mathcal{C}_{g}$ de la fonction $g : x \mapsto f(x) + b$ (définie sur $I$) s'obtient en translatant $\mathcal{C}_f$ du vecteur $\vec{v}(0;b)$.
\begin{itemize}
    \item Si $b > 0$ : translation vers le haut.
    \item Si $b < 0$ : translation vers le bas.
\end{itemize}
\emph{Justification :} Pour tout $x \in I$, le point $M(x;f(x))$ de $\mathcal{C}_f$ correspond au point $M'(x;f(x)+b)$ de $\mathcal{C}_g$. On a $\overrightarrow{MM'}=(0;b)$.
\end{propriete}

\begin{propriete}[Translation horizontale : $x \mapsto f(x - a)$]
La courbe $\mathcal{C}_{h}$ de la fonction $h : x \mapsto f(x - a)$ (définie sur $I+a = \{x \in \mathbb{R} \mid x-a \in I\}$) s'obtient en translatant $\mathcal{C}_f$ du vecteur $\vec{v}(a;0)$.
\begin{itemize}
    \item Si $a > 0$ : translation vers la \textbf{droite}.
    \item Si $a < 0$ : translation vers la \textbf{gauche}.
\end{itemize}
\emph{Justification :} Soit $X = x-a$. Le point $M(X;f(X))$ de $\mathcal{C}_f$ correspond au point $M'(x;f(x-a))$ de $\mathcal{C}_h$ avec $x = X+a$. On a $\overrightarrow{MM'}=(a;0)$.
\end{propriete}

\begin{attention}[Le piège du signe horizontal]
C'est l'erreur n°1 au Probatoire : \textbf{« $x \mapsto f(x-2)$ décale vers la droite »}. Pourquoi ? Pour retrouver la valeur $f(0)$ sur la nouvelle courbe, il faut prendre $x=2$. Le graphique « suit le signe opposé » de ce qui est dans la parenthèse.
\begin{center}
\textbf{Mémo :} $x \mapsto f(x-a)$ $\rightarrow$ Translation de vecteur $\vec{v}(a;0)$.
\end{center}
\end{attention}

\begin{propriete}[Translation composée : $x \mapsto f(x-a) + b$]
La courbe de $x \mapsto f(x-a)+b$ s'obtient par translation de $\mathcal{C}_f$ du vecteur $\vec{v}(a;b)$. L'ordre des deux translations (horizontal puis vertical, ou l'inverse) est indifférent car les translations commutent.
\end{propriete}

\begin{propriete}[Symétrie par rapport à $Ox$ : $x \mapsto -f(x)$]
La courbe $\mathcal{C}_{k}$ de la fonction $k : x \mapsto -f(x)$ (définie sur $I$) est le symétrique de $\mathcal{C}_f$ par rapport à l'axe des abscisses ($Ox$).
\emph{Justification :} Le point $M(x;f(x))$ devient $M'(x;-f(x))$. $Ox$ est le médiateur de $[MM']$.
\end{propriete}

\begin{propriete}[Valeur absolue : $x \mapsto |f(x)|$]
La courbe $\mathcal{C}_{\ell}$ de la fonction $\ell : x \mapsto |f(x)|$ (définie sur $I$) se construit ainsi :
\begin{enumerate}
    \item On conserve la partie de $\mathcal{C}_f$ située dans le demi-plan $y \ge 0$ (au-dessus ou sur $Ox$).
    \item On supprime la partie de $\mathcal{C}_f$ située dans le demi-plan $y < 0$ (sous $Ox$).
    \item On reporte le symétrique (par rapport à $Ox$) de cette partie supprimée.
\end{enumerate}
Les points d'intersection de $\mathcal{C}_f$ avec $Ox$ (zéros de $f$) restent fixes.
\end{propriete}

\begin{aretenir}[Résumé visuel des transformations]
\begin{center}
\begin{tabular}{|c|c|c|}\hline
\textbf{Fonction $g(x)$} & \textbf{Transformation géométrique} & \textbf{Vecteur / Axe} \\ \hline
$f(x) + b$ & Translation verticale & $\vec{v}(0;b)$ \\ \hline
$f(x-a)$ & Translation horizontale & $\vec{v}(a;0)$ \\ \hline
$f(x-a)+b$ & Translation composée & $\vec{v}(a;b)$ \\ \hline
$-f(x)$ & Symétrie axiale & Axe $Ox$ \\ \hline
$|f(x)|$ & « Miroir » sur $Ox$ pour $y<0$ & Axe $Ox$ (plissage) \\ \hline
\end{tabular}
\end{center}
\end{aretenir}

\courssub{Ordre des opérations : La règle d'or pour les compositions}

Lorsque plusieurs transformations s'enchaînent (ex: $x \mapsto -f(x-1)+2$), l'ordre est \textbf{crucial} si l'on mélange translations et symétrie/homothétie. La règle universelle est : \textbf{On transforme d'abord la variable $x$ (horizontal), puis les images $y$ (vertical).}

\begin{methode}[Ordre de construction pour $g(x) = A \cdot f(Bx + C) + D$ (cas général) / Ici $A=\pm1, B=1$]
Pour tracer $\mathcal{C}_g$ à partir de $\mathcal{C}_f$ avec $g(x) = \pm f(x-a) + b$ :
\begin{enumerate}
    \item \textbf{Horizontal (sur $x$) :} Appliquer $x \mapsto x-a$ (Translation $\vec{v}(a;0)$). On obtient $\mathcal{C}_{f(x-a)}$.
    \item \textbf{Symétrie / Homothétie (sur $y$, coefficient de $f$) :} Appliquer le signe $-$ si présent (Symétrie $Ox$). On obtient $\mathcal{C}_{-f(x-a)}$.
    \item \textbf{Vertical (sur $y$, constante ajoutée) :} Appliquer $+b$ (Translation $\vec{v}(0;b)$). On obtient $\mathcal{C}_{-f(x-a)+b}$.
\end{enumerate}
\emph{Note :} Les translations entre elles commutent (étape 1 et 3 peuvent s'inverser), mais la symétrie $Ox$ ne commute pas avec la translation verticale.
\end{methode}

\begin{exemplebox}[Ordre des opérations : $g(x) = -f(x-1)+2$]
Soit $f(x)=x^2$ sur $[-2;2]$. On veut tracer $g(x)=-(x-1)^2+2$.
\begin{enumerate}
    \item \textbf{Départ :} $\mathcal{C}_f$ : Parabole sommet $S_f(0;0)$, passant par $A(-2;4), B(2;4)$.
    \item \textbf{Étape 1 (Horizontal $x \to x-1$) :} Translation $\vec{v}(1;0)$. Nouveau sommet $S_1(1;0)$. Points $A_1(-1;4), B_1(3;4)$. Courbe : $y=(x-1)^2$.
    \item \textbf{Étape 2 (Symétrie $Ox$ : $-f$) :} Symétrie par rapport à $Ox$. Sommet $S_2(1;0)$ (inchangé car sur $Ox$). Points $A_2(-1;-4), B_2(3;-4)$. Courbe : $y=-(x-1)^2$. \emph{Attention : si on avait fait le $+2$ avant le $-$, le sommet serait monté à $y=2$ puis descendu à $y=-2$ : erreur !}
    \item \textbf{Étape 3 (Vertical $+2$) :} Translation $\vec{v}(0;2)$. Sommet final $S_g(1;2)$. Points $A_g(-1;-2), B_g(3;-2)$. Courbe finale : $y=-(x-1)^2+2$.
\end{enumerate}
\end{exemplebox}

\courssub{Application aux fonctions de référence : Formes canoniques et standards}

C'est ici que tout se relie. Les formes canoniques (polynôme degré 2) et standards (homographique) ne sont rien d'autre que des écritures qui font \emph{apparaître} les vecteurs de translation par rapport aux courbes de base $x \mapsto ax^2$ et $x \mapsto \frac{k}{x}$.

\begin{propriete}[Fonction polynôme du second degré : Forme canonique]
Soit $f(x) = ax^2+bx+c$ ($a \neq 0$). Sa forme canonique est $f(x) = a(x-\alpha)^2 + \beta$ avec $\alpha = -\frac{b}{2a}$ et $\beta = f(\alpha) = \frac{4ac-b^2}{4a}$.
La courbe $\mathcal{C}_f$ s'obtient par translation de la parabole $\mathcal{C}_0 : y = ax^2$ (sommet à l'origine $O$) du vecteur $\vec{v}(\alpha;\beta)$.
Le sommet de $\mathcal{C}_f$ est $S(\alpha;\beta)$. L'axe de symétrie est la droite $x=\alpha$.
\end{propriete}

\begin{propriete}[Fonction homographique : Forme standard]
Soit $f(x) = \frac{ax+b}{cx+d}$ ($c \neq 0, ad-bc \neq 0$). On effectue la division euclidienne du numérateur par le dénominateur pour l'écrire sous la forme :
\[
f(x) = \alpha + \frac{k}{x-\beta} \quad \text{avec} \quad \alpha = \frac{a}{c},\quad \beta = -\frac{d}{c},\quad k = \frac{bc-ad}{c^2}.
\]
La courbe $\mathcal{C}_f$ s'obtient par translation de l'hyperbole $\mathcal{H}_0 : y = \frac{k}{x}$ (centre à l'origine $O$, asymptotes $Ox$ et $Oy$) du vecteur $\vec{v}(\beta;\alpha)$.
Le centre de symétrie de $\mathcal{C}_f$ est $I(\beta;\alpha)$. Les asymptotes sont les droites $x=\beta$ (verticale) et $y=\alpha$ (horizontale).
\end{propriete}

\begin{exoresolu}[{De la forme canonique au tracé : $f(x) = 2x^2 - 8x + 5$ sur $[-1;5]$}]
\textbf{Énoncé :} Étudier $f$ et tracer $\mathcal{C}_f$ par translation.
\textbf{Solution :}
\begin{enumerate}
    \item \textbf{Forme canonique :} $f(x) = 2(x^2 - 4x) + 5 = 2[(x-2)^2 - 4] + 5 = 2(x-2)^2 - 8 + 5 = 2(x-2)^2 - 3$.
    \item \textbf{Identification :} $a=2$, $\alpha=2$, $\beta=-3$. Vecteur de translation $\vec{v}(2;-3)$ par rapport à $y=2x^2$.
    \item \textbf{Courbe de base $y=2x^2$ :} Parabole sommet $O(0;0)$, branche vers le haut ($a>0$). Points : $(-1;2), (1;2), (-2;8), (2;8)$.
    \item \textbf{Translation $\vec{v}(2;-3)$ :} Sommet $S(2;-3)$. Points : $A(1;-1), B(3;-1), C(0;5), D(4;5)$.
    \item \textbf{Ensemble de définition :} $[-1;5]$. Bornes : $f(-1)=15$, $f(5)=15$.
    \item \textbf{Tableau de variations :} Décroissante sur $[-1;2]$, croissante sur $[2;5]$. Minimum $m=-3$ en $x=2$. Maximum $M=15$ aux bornes.
\end{enumerate}
\begin{popfigure}
\begin{tikzpicture}[scale=0.6, >=Stealth]
    % Axes
    \draw[->] (-2,0) -- (7,0) node[right] {$x$};
    \draw[->] (0,-4) -- (0,17) node[above] {$y$};
    \foreach \x in {-1,1,2,3,4,5,6} \draw (\x,0.1) -- (\x,-0.1) node[below] {\small $\x$};
    \foreach \y in {-3,-1,1,3,5,7,9,11,13,15} \draw (0.1,\y) -- (-0.1,\y) node[left] {\small $\y$};
    % Base y=2x^2 (pointillés)
    \draw[popmuted, dashed, thick, domain=-1.5:3.5, samples=100] plot (\x, {2*\x*\x});
    \node[popmuted] at (3.5,12) {\small $y=2x^2$};
    % Translated f(x)
    \draw[popblue, very thick, domain=-1:5, samples=100] plot (\x, {2*(\x-2)*(\x-2)-3});
    % Points
    \fill[popblue] (2,-3) circle (3pt) node[below right] {$S(2;-3)$};
    \fill[popblue] (-1,15) circle (3pt) node[above left] {$A(-1;15)$};
    \fill[popblue] (5,15) circle (3pt) node[above right] {$B(5;15)$};
    \fill[popblue] (1,-1) circle (3pt) node[below] {$(1;-1)$};
    \fill[popblue] (3,-1) circle (3pt) node[below] {$(3;-1)$};
    % Vector
    \draw[poporange, ->, thick] (0,0) -- (2,-3) node[midway, above right] {$\vec{v}(2;-3)$};
    \fill[poporange] (0,0) circle (2pt) node[below left] {$O$};
\end{tikzpicture}
\legende{Construction de $\mathcal{C}_f$ par translation de $y=2x^2$ du vecteur $\vec{v}(2;-3)$.}
\end{popfigure}
\end{exoresolu}

\begin{exoresolu}[{Fonction homographique : $f(x) = \frac{3x-1}{x-2}$ sur $]{- \infty}; 2[ \cup ]2; +\infty[$}]
\textbf{Énoncé :} Écrire sous forme standard, donner centre et asymptotes, tracer $\mathcal{C}_f$.
\textbf{Solution :}
\begin{enumerate}
    \item \textbf{Division :} $3x-1 = 3(x-2) + 5$. Donc $f(x) = \frac{3(x-2)+5}{x-2} = 3 + \frac{5}{x-2}$.
    \item \textbf{Identification :} $\alpha=3$, $\beta=2$, $k=5$. Vecteur $\vec{v}(2;3)$ par rapport à $y=\frac{5}{x}$.
    \item \textbf{Courbe de base $y=\frac{5}{x}$ :} Hyperbole centre $O$, asymptotes $Ox, Oy$, branches dans quadrants 1 et 3 ($k>0$).
    \item \textbf{Translation $\vec{v}(2;3)$ :} Centre $I(2;3)$. Asymptotes : $x=2$ (verticale), $y=3$ (horizontale). Branches dans quadrants relatifs à $I$ : haut-droite et bas-gauche.
    \item \textbf{Points caractéristiques :} $f(0)=0,5 \to A(0;0,5)$. $f(1)=-2 \to B(1;-2)$. $f(3)=8 \to C(3;8)$. $f(4)=5,5 \to D(4;5,5)$.
\end{enumerate}
\begin{popfigure}
\begin{tikzpicture}[scale=0.7, >=Stealth]
    % Axes
    \draw[->] (-2,0) -- (7,0) node[right] {$x$};
    \draw[->] (0,-5) -- (0,9) node[above] {$y$};
    \foreach \x in {-1,1,2,3,4,5,6} \draw (\x,0.1) -- (\x,-0.1) node[below] {\small $\x$};
    \foreach \y in {-4,-2,1,3,5,7} \draw (0.1,\y) -- (-0.1,\y) node[left] {\small $\y$};
    % Asymptotes
    \draw[popline, dashed] (2,-5) -- (2,9) node[above] {$x=2$};
    \draw[popline, dashed] (-2,3) -- (7,3) node[right] {$y=3$};
    % Base y=5/x (pointillés)
    \draw[popmuted, dashed, thick, domain=-2:-0.5, samples=50] plot (\x, {5/\x});
    \draw[popmuted, dashed, thick, domain=0.5:4, samples=50] plot (\x, {5/\x});
    \node[popmuted] at (-1.5,-3) {\small $y=\frac{5}{x}$};
    % Translated f(x)
    \draw[popblue, very thick, domain=-2:1.5, samples=100] plot (\x, {3+5/(\x-2)});
    \draw[popblue, very thick, domain=2.5:6, samples=100] plot (\x, {3+5/(\x-2)});
    % Center
    \fill[poporange] (2,3) circle (3pt) node[above right] {$I(2;3)$};
    % Points
    \fill[popblue] (0,0.5) circle (2pt) node[left] {$A(0;0,5)$};
    \fill[popblue] (1,-2) circle (2pt) node[below] {$B(1;-2)$};
    \fill[popblue] (3,8) circle (2pt) node[above] {$C(3;8)$};
    \fill[popblue] (4,5.5) circle (2pt) node[right] {$D(4;5,5)$};
    % Vector
    \draw[poporange, ->, thick] (0,0) -- (2,3) node[midway, above left] {$\vec{v}(2;3)$};
    \fill[poporange] (0,0) circle (2pt) node[below left] {$O$};
\end{tikzpicture}
\legende{Hyperbole $f(x)=3+\frac{5}{x-2}$ obtenue par translation de $y=\frac{5}{x}$.}
\end{popfigure}
\end{exoresolu}

\courssub{La valeur absolue : $x \mapsto |f(x)|$ — Le « pliage » sur $Ox$}

C'est une transformation très particulière, non bijective, qui « redresse » la courbe.

\begin{propriete}[Construction de $\mathcal{C}_{|f|}$]
Soit $f$ définie sur $I$. Pour tracer $\mathcal{C}_g$ avec $g(x)=|f(x)|$ :
\begin{enumerate}
    \item Tracer $\mathcal{C}_f$ (pointillés ou trait fin).
    \item Repérer les intervalles où $f(x) \ge 0$ et où $f(x) < 0$ (zéros de $f$).
    \item \textbf{Conserver} l'arc de $\mathcal{C}_f$ où $f(x) \ge 0$.
    \item \textbf{Symétriser} par rapport à $Ox$ l'arc de $\mathcal{C}_f$ où $f(x) < 0$ (on « replie » vers le haut).
    \item Les points d'intersection de $\mathcal{C}_f$ avec $Ox$ (zéros de $f$) sont des points anguleux (souvent des minima locaux) sauf si $f$ s'annule avec une tangente horizontale (ex: $x^2$ en 0).
\end{enumerate}
\end{propriete}

\begin{exemplebox}[{Construction pas à pas : $g(x) = |x^2 - 4|$ sur $[-3;3]$}]
On part de $f(x)=x^2-4$ (parabole sommet $S_f(0;-4)$, zéros $-2$ et $2$).
\begin{enumerate}
    \item \textbf{Sur $[-3;-2] \cup [2;3]$ :} $f(x) \ge 0$. On garde la parabole. Points : $A(-3;5), B(-2;0), C(2;0), D(3;5)$.
    \item \textbf{Sur $[-2;2]$ :} $f(x) \le 0$. On symétrise par rapport à $Ox$. Le sommet $S_f(0;-4)$ devient $S_g(0;4)$. La courbe devient une « bosse » au-dessus de $Ox$.
    \item \textbf{Résultat :} $\mathcal{C}_g$ a une forme de « W » arrondi. Minima en $-2$ et $2$ (valeur 0). Maximum en $0$ (valeur 4). Maximum aux bornes $\pm 3$ (valeur 5).
\end{enumerate}
\begin{popfigure}
\begin{tikzpicture}[scale=0.7, >=Stealth]
    \draw[->] (-3.5,0) -- (3.5,0) node[right] {$x$};
    \draw[->] (0,-5) -- (0,6) node[above] {$y$};
    \foreach \x in {-3,-2,-1,1,2,3} \draw (\x,0.1) -- (\x,-0.1) node[below] {\small $\x$};
    \foreach \y in {-4,-2,2,4} \draw (0.1,\y) -- (-0.1,\y) node[left] {\small $\y$};
    % f(x) = x^2-4 (dashed)
    \draw[popmuted, dashed, thick, domain=-3:3, samples=100] plot (\x, {\x*\x-4});
    \node[popmuted] at (3,4) {\small $y=x^2-4$};
    % |f(x)| (solid)
    \draw[popblue, very thick, domain=-3:-2, samples=50] plot (\x, {\x*\x-4});
    \draw[popblue, very thick, domain=-2:2, samples=100] plot (\x, {4-\x*\x});
    \draw[popblue, very thick, domain=2:3, samples=50] plot (\x, {\x*\x-4});
    % Key points
    \fill[popblue] (-2,0) circle (2pt) node[below left] {$-2$};
    \fill[popblue] (2,0) circle (2pt) node[below right] {$2$};
    \fill[popblue] (0,4) circle (2pt) node[above] {$S_g(0;4)$};
    \fill[popblue] (-3,5) circle (2pt) node[above left] {$5$};
    \fill[popblue] (3,5) circle (2pt) node[above right] {$5$};
    % Arrows showing folding
    \draw[poporange, ->, thick] (0,-4) -- (0,4) node[midway, right] {\small Pliage};
    \fill[poporange] (0,-4) circle (2pt) node[below] {$S_f$};
\end{tikzpicture}
\legende{Construction de $y=|x^2-4|$ par pliage de la partie négative de $y=x^2-4$ sur $Ox$.}
\end{popfigure}
\end{exemplebox}

\courssub{Méthode unifiée : Tracer par points caractéristiques}

Au Probatoire, on ne trace pas « au feeling ». On utilise la méthode des points caractéristiques transformés.

\begin{methode}[Tracer $\mathcal{C}_g$ pour $g(x) = f(x-a)+b$ ou $g(x) = |f(x-a)+b|$ etc.]
\begin{enumerate}
    \item \textbf{Identifier la fonction de base $f_0$ et sa courbe $\mathcal{C}_0$ connue} (ex: $x \mapsto ax^2$, $x \mapsto k/x$, $x \mapsto x$).
    \item \textbf{Décomposer $g$ en transformations successives} dans l'ordre : Horizontal $\to$ Symétrie/Homothétie $\to$ Vertical $\to$ Valeur absolue (si présent, \textbf{toujours en dernier}).
    \item \textbf{Lister les points caractéristiques de $\mathcal{C}_0$} : Sommet / Centre, Intersections avec axes, Points aux bornes de l'intervalle d'étude, Points de passage commodes.
    \item \textbf{Appliquer les transformations à ces points} (coordonnées par coordonnées) pour obtenir les points de $\mathcal{C}_g$.
    \item \textbf{Transformer les asymptotes} (si homographique) : $x=\beta \to x=\beta+a$, $y=\alpha \to y=\alpha+b$ (et symétrie si $-f$).
    \item \textbf{Tracer} : Reporter les nouveaux points, tracer les nouvelles asymptotes (trait pointillé), relier les points en respectant la forme de la courbe de base (convexité, branches).
\end{enumerate}
\end{methode}

\begin{exoresolu}[{Synthèse : $g(x) = \left| \frac{2}{x-1} - 3 \right|$ sur $]{- \infty}; 1[ \cup ]1; 3]$}]
\textbf{Énoncé :} Tracer $\mathcal{C}_g$ par transformations successives.
\textbf{Solution :}
\begin{enumerate}
    \item \textbf{Base :} $f_0(x) = \frac{2}{x}$ ($k=2>0$, quadrants 1 et 3). Asymptotes $Ox, Oy$.
    \item \textbf{Décomposition :}
        \begin{itemize}
            \item $f_1(x) = f_0(x-1) = \frac{2}{x-1}$ (Translation $\vec{v}(1;0)$). Asymptotes $x=1, y=0$. Centre $I_1(1;0)$.
            \item $f_2(x) = f_1(x) - 3 = \frac{2}{x-1} - 3$ (Translation $\vec{v}(0;-3)$). Asymptotes $x=1, y=-3$. Centre $I_2(1;-3)$.
            \item $g(x) = |f_2(x)|$ (Valeur absolue / Pliage sur $Ox$).
        \end{itemize}
    \item \textbf{Points caractéristiques de $f_0$ (sur domaine adapté) :} $A_0(-2;-1), B_0(-1;-2), C_0(1;2), D_0(2;1)$. (On évite 0).
    \item \textbf{Suivi des points :}
    \begin{center}
    \begin{tabular}{|c|c|c|c|c|}\hline
    Point & $f_0$ & $f_1(x)=f_0(x-1)$ & $f_2(x)=f_1(x)-3$ & $g(x)=|f_2(x)|$ \\ \hline
    $A$ & $(-2;-1)$ & $(-1;-1)$ & $(-1;-4)$ & $(-1;4)$ \\ \hline
    $B$ & $(-1;-2)$ & $(0;-2)$ & $(0;-5)$ & $(0;5)$ \\ \hline
    $C$ & $(1;2)$ & $(2;2)$ & $(2;-1)$ & $(2;1)$ \\ \hline
    $D$ & $(2;1)$ & $(3;1)$ & $(3;-2)$ & $(3;2)$ \\ \hline
    \end{tabular}
    \end{center}
    \item \textbf{Zéros de $f_2$ (points de pliage) :} $\frac{2}{x-1}-3=0 \Rightarrow 2=3x-3 \Rightarrow 3x=5 \Rightarrow x=\frac{5}{3} \approx 1,67$. Point $Z(\frac{5}{3};0)$.
    \item \textbf{Asymptotes finales :} $x=1$ (inchangée). L'asymptote horizontale $y=-3$ (branche du bas) est pliée vers $y=3$. La branche du haut (vers $x=1^+$) tend vers $+\infty$. Vers $x \to +\infty$, $f_2 \to -3^-$, $|f_2| \to 3^+$. Donc asymptote horizontale $y=3$ pour la branche pliée.
    \item \textbf{Tracé final :} Branche gauche ($x<1$) : $f_2 < 0$ tout le temps, tout se plie au-dessus : branche dans quadrant 2 relatif à $(1;3)$, asymptotes $x=1$ et $y=3$. Branche droite ($1<x\le 3$) : Partie haute conservée ($1<x<5/3$), partie basse pliée ($5/3<x\le 3$). Point $Z(1,67;0)$ anguleux.
\end{enumerate}
\begin{popfigure}
\begin{tikzpicture}[scale=0.7, >=Stealth]
    \draw[->] (-2,0) -- (4,0) node[right] {$x$};
    \draw[->] (0,-6) -- (0,6) node[above] {$y$};
    \foreach \x in {-1,1,2,3} \draw (\x,0.1) -- (\x,-0.1) node[below] {\small $\x$};
    \foreach \y in {-5,-3,-1,1,3,5} \draw (0.1,\y) -- (-0.1,\y) node[left] {\small $\y$};
    % Asymptotes
    \draw[popline, dashed] (1,-6) -- (1,6) node[above] {$x=1$};
    \draw[popline, dashed] (-2,3) -- (4,3) node[right] {$y=3$};
    \draw[popline, dashed] (-2,-3) -- (1,-3) node[left] {$y=-3$ (ref.)};
    % f2(x) = 2/(x-1) - 3 (dashed)
    \draw[popmuted, dashed, thick, domain=-2:0.8, samples=100] plot (\x, {2/(\x-1)-3});
    \draw[popmuted, dashed, thick, domain=1.2:3.5, samples=100] plot (\x, {2/(\x-1)-3});
    % g(x) = |f2(x)| (solid)
    % Left branch x<1 : f2 < -3 < 0, folded up. y = -f2 = 3 - 2/(x-1) = 3 + 2/(1-x). Asymptote y=3.
    \draw[popblue, very thick, domain=-2:0.8, samples=100] plot (\x, {3 - 2/(\x-1)}); 
    % Right branch 1 < x < 5/3 : f2 > 0, kept.
    \draw[popblue, very thick, domain=1.2:1.66, samples=100] plot (\x, {2/(\x-1)-3});
    % Right branch 5/3 < x < 3 : f2 < 0, folded up. y = -f2 = 3 - 2/(x-1).
    \draw[popblue, very thick, domain=1.68:3, samples=100] plot (\x, {3 - 2/(\x-1)});
    % Points
    \fill[popblue] (-1,4) circle (2pt) node[above left] {$(-1;4)$};
    \fill[popblue] (0,5) circle (2pt) node[above] {$(0;5)$};
    \fill[popblue] (2,1) circle (2pt) node[below right] {$(2;1)$};
    \fill[popblue] (3,2) circle (2pt) node[right] {$(3;2)$};
    \fill[popblue] (1.666,0) circle (2pt) node[below] {$Z(5/3;0)$};
    % Center f2
    \fill[poporange] (1,-3) circle (2pt) node[below right] {$I_2(1;-3)$};
\end{tikzpicture}
\legende{Tracé de $g(x)=\left|\frac{2}{x-1}-3\right|$ par transformations successives.}
\end{popfigure}
\end{exoresolu}

\courssub{Complément (Terminale) : Parité et symétries}

\begin{savaistu}[Fonctions paires et impaires : Les symétries « naturelles »]
En Terminale, on formalise ce qu'on utilise implicitement ici.
\begin{itemize}
    \item \textbf{Fonction paire :} $f(-x)=f(x)$. $\mathcal{C}_f$ symétrique par rapport à $Oy$ (axe des ordonnées). Ex: $x \mapsto ax^2$, $x \mapsto |x|$, $x \mapsto \cos x$.
    \item \textbf{Fonction impaire :} $f(-x)=-f(x)$. $\mathcal{C}_f$ symétrique par rapport à l'origine $O$ (rotation de $180^\circ$). Ex: $x \mapsto k/x$, $x \mapsto x^3$, $x \mapsto \sin x$.
\end{itemize}
Cela explique pourquoi $y=ax^2$ a un sommet sur $Oy$ et pourquoi $y=k/x$ a son centre en $O$. Les translations brisent ces symétries par rapport aux axes d'origine.
\end{savaistu}

\courssub{Exercices d'entraînement (Probatoire Style)}

\begin{exercice}[Transformation de parabole]
Soit $f(x) = -x^2 + 2x + 3$ définie sur $[-2;4]$.
\begin{enumerate}
    \item Écrire $f$ sous forme canonique. En déduire le vecteur de translation par rapport à $y=-x^2$.
    \item Tracer $\mathcal{C}_f$.
    \item Tracer la courbe $\mathcal{C}_g$ de $g(x) = |f(x)|$ sur $[-2;4]$.
    \item Déterminer graphiquement le nombre de solutions de $|f(x)| = 2$.
\end{enumerate}
\end{exercice}

\begin{exercice}[Transformation d'hyperbole]
Soit $f(x) = \frac{x+1}{x-3}$ définie sur $]{- \infty}; 3[ \cup ]3; 5]$.
\begin{enumerate}
    \item Écrire $f$ sous la forme $\alpha + \frac{k}{x-\beta}$. Préciser $\alpha, \beta, k$.
    \item Tracer $\mathcal{C}_f$ par translation de $y=\frac{k}{x}$.
    \item Soit $h(x) = -f(x-2)$. Décrire les transformations successives pour obtenir $\mathcal{C}_h$ à partir de $\mathcal{C}_f$. Donner les nouvelles asymptotes et le centre de symétrie de $\mathcal{C}_h$.
\end{enumerate}
\end{exercice}

\begin{exercice}[Problème concret : Couverture Wi-Fi]
Une antenne-relais MTN émet un signal dont la puissance $P$ (en dBm) en fonction de la distance $x$ (en km) est modélisée par $P(x) = 20 - \frac{50}{x+1}$ pour $x \in [0; 10]$.
\begin{enumerate}
    \item Écrire $P$ sous forme standard. Identifier la translation par rapport à $y=-\frac{50}{x}$.
    \item Tracer $\mathcal{C}_P$. Interpréter l'asymptote horizontale.
    \item Une mise à jour logicielle augmente la puissance de 5 dBm partout. Nouvelle fonction $P_{new}$. Tracer $\mathcal{C}_{P_{new}}$ par translation. Quelle transformation géométrique ?
    \item On installe un répéteur qui décale la zone de couverture de 2 km vers l'est. Nouvelle fonction $P_{rep}(x) = P_{new}(x-2)$. Tracer $\mathcal{C}_{P_{rep}}$. Donner l'équation de la nouvelle asymptote verticale.
\end{enumerate}
\end{exercice}

\begin{corrige}[Exercice 1]
\begin{enumerate}
    \item $f(x) = -(x^2 - 2x) + 3 = -[(x-1)^2 - 1] + 3 = -(x-1)^2 + 4$. Forme canonique : $a=-1, \alpha=1, \beta=4$. Translation $\vec{v}(1;4)$ par rapport à $y=-x^2$.
    \item Sommet $S(1;4)$ (maximum). Zéros : $-(x-1)^2+4=0 \Rightarrow (x-1)^2=4 \Rightarrow x-1=\pm 2 \Rightarrow x=-1, x=3$. Bornes : $f(-2)=-5, f(4)=-5$. \textbf{Croissante sur $[-2;1]$, décroissante sur $[1;4]$}. Maximum $M=4$ en $x=1$. Minimum $m=-5$ aux bornes.
    \item $g(x)=|f(x)|$. $f \ge 0$ sur $[-1;3]$, $f \le 0$ sur $[-2;-1] \cup [3;4]$. Pliage des deux bouts vers le haut. Minima en $-1$ et $3$ (valeur 0). Maximum central $4$ en $1$. Maxima aux bornes $5$ en $-2$ et $4$.
    \item Équation $|f(x)|=2 \Leftrightarrow f(x)=2$ ou $f(x)=-2$.
        $f(x)=2 \Rightarrow -(x-1)^2+4=2 \Rightarrow (x-1)^2=2 \Rightarrow x=1\pm\sqrt{2} \approx -0,4 ; 2,4$ (dans $[-1;3]$, 2 solutions).
        $f(x)=-2 \Rightarrow -(x-1)^2+4=-2 \Rightarrow (x-1)^2=6 \Rightarrow x=1\pm\sqrt{6} \approx -1,4 ; 3,4$ (dans $[-2;-1] \cup [3;4]$, 2 solutions).
        Total : \textbf{4 solutions}.
\end{enumerate}
\end{corrige}

\begin{corrige}[Exercice 2]
\begin{enumerate}
    \item $x+1 = 1(x-3) + 4$. $f(x) = 1 + \frac{4}{x-3}$. $\alpha=1, \beta=3, k=4$.
    \item Base $y=4/x$ (quadrants 1,3). Translation $\vec{v}(3;1)$. Centre $I(3;1)$. Asymptotes $x=3, y=1$. Branches haut-droite / bas-gauche par rapport à $I$. Domaine $]{- \infty};3[ \cup ]3;5]$. En $x=5, f(5)=3$. En $x \to {- \infty}, f \to 1^+$.
    \item $h(x) = -f(x-2)$.
        Ordre : 1. Horizontal $x \to x-2$ : Translation $\vec{v}(2;0)$. Centre $I_1(5;1)$. Asymptotes $x=5, y=1$.
        2. Symétrie $Ox$ (coeff $-1$) : Centre $I_h(5;-1)$. Asymptotes $x=5, y=-1$. Branches inversées (haut-gauche / bas-droite par rapport à $I_h$).
\end{enumerate}
\end{corrige}

\begin{corrige}[Exercice 3]
\begin{enumerate}
    \item $P(x) = 20 - \frac{50}{x+1} = 20 + \frac{-50}{x-(-1)}$. $\alpha=20, \beta=-1, k=-50$. Translation $\vec{v}(-1;20)$ par rapport à $y=-50/x$ (quadrants 2 et 4).
    \item Centre $I(-1;20)$ (hors domaine $x\ge 0$). Asymptotes $x=-1$ (verticale, hors domaine), $y=20$ (horizontale). Interprétation : Puissance maximale théorique 20 dBm à distance infinie. Sur $[0;10]$, $P$ croissante (car $k<0$ sur branche droite). $P(0)=-30$, $P(10)=20-50/11 \approx 15,45$.
    \item $P_{new}(x) = P(x) + 5 = 25 - \frac{50}{x+1}$. Translation verticale $\vec{v}(0;5)$. Asymptote horizontale $y=25$. Courbe décalée vers le haut de 5 unités.
    \item $P_{rep}(x) = P_{new}(x-2) = 25 - \frac{50}{(x-2)+1} = 25 - \frac{50}{x-1}$. Translation horizontale $\vec{v}(2;0)$. Asymptote verticale : $x=1$. Domaine : $[2;12]$. La zone de couverture a glissé de 2 km vers la droite.
\end{enumerate}
\end{corrige}

\begin{attention}[Check-list avant de rendre la copie (Probatoire)]
\begin{itemize}
    \item \textbf{Domaine :} As-tu reporté les bornes du domaine sur l'axe $x$ ? As-tu calculé les images des bornes ?
    \item \textbf{Asymptotes (homographique) :} Sont-elles tracées en \textbf{pointillés} ? Sont-elles \textbf{étiquetées} ($x=...$, $y=...$) ? Ont-elles été \textbf{translatées} correctement ?
    \item \textbf{Valeur absolue :} As-tu bien gardé la partie $y\ge 0$ et \textbf{seulement} symétrisé la partie $y<0$ ? Les zéros sont-ils des points anguleux (souvent) ?
    \item \textbf{Ordre :} Pour $-f(x-a)+b$, as-tu fait : 1. Translation $a$ (horizontal), 2. Symétrie $-$ (vertical), 3. Translation $b$ (vertical) ?
    \item \textbf{Unités :} Échelles indiquées ? Repère orthonormé respecté (même unité sur les deux axes) ?
\end{itemize}
\end{attention}

\coursec{Résolution graphique d'équations et d'inéquations ; Paramètre m}

\courssub{Principes fondamentaux : Intersection courbe -- droite horizontale}

Après avoir maîtrisé la construction des courbes par transformations (section 4), nous abordons maintenant l'une des applications les plus puissantes de la représentation graphique : la résolution d'équations et d'inéquations. L'idée directrice est d'une simplicité trompeuse : \emph{une équation $f(x)=m$ demande les abscisses des points où la courbe $\mathcal{C}_f$ rencontre la droite horizontale $y=m$}.

\begin{definition}[Solutions graphiques d'une équation]
Soit $f$ une fonction définie sur un intervalle borné $I$ et $m\in\mathbb{R}$. Les solutions de l'équation $f(x)=m$ sur $I$ sont exactement les abscisses des points d'intersection de la courbe représentative $\mathcal{C}_f$ (restreinte à $I$) avec la droite d'équation $y=m$.
\end{definition}

\begin{propriete}[Nombre de solutions pour un polynôme du second degré]
Soit $f(x)=ax^2+bx+c$ définie sur un intervalle borné $I=[\alpha;\beta]$. La courbe $\mathcal{C}_f$ est un arc de parabole. Pour tout réel $m$, l'équation $f(x)=m$ admet sur $I$ :
\begin{itemize}
\item \textbf{0 solution} si la droite $y=m$ ne coupe pas l'arc de parabole ;
\item \textbf{1 solution} si la droite $y=m$ est tangente à la parabole (sommet intérieur à $I$) ou passe par un extremum aux bornes, ou ne coupe qu'une seule branche accessible ;
\item \textbf{2 solutions} si la droite $y=m$ coupe l'arc de parabole en deux points distincts.
\end{itemize}
\end{propriete}

\begin{attention}[Piège classique : abscisses vs ordonnées]
\emph{Ne confondez jamais :} on lit des \textbf{abscisses} (valeurs de $x$) sur l'axe horizontal. L'ordonnée $m$ est fixée \emph{avant} de tracer la droite horizontale. Pour $f(x)>0$, on cherche où la courbe est \textbf{strictement au-dessus} de l'axe $Ox$ (droite $y=0$).
\end{attention}

\courssub{Méthodologie pas à pas}

\begin{methode}[Résolution graphique de $f(x)=m$ sur un intervalle borné $I$]
\begin{enumerate}
\item \textbf{Tracer} soigneusement la courbe $\mathcal{C}_f$ sur $I$ (tableau de variations, points remarquables, asymptotes si fonction homographique).
\item \textbf{Tracer} la droite horizontale $y=m$ (pointillés, couleur distincte).
\item \textbf{Repérer} les points d'intersection $A_1, A_2, \dots$ entre $\mathcal{C}_f$ et $y=m$.
\item \textbf{Projeter} ces points sur l'axe des abscisses (droites verticales pointillées) pour lire les valeurs $x_1, x_2, \dots$
\item \textbf{Écrire} l'ensemble solution : $S=\{x_1, x_2, \dots\}$ ou $S=\varnothing$ si aucune intersection.
\end{enumerate}
\end{methode}

\begin{methode}[Résolution graphique de $f(x) > m$ (ou $\ge, <, \le$) sur $I$]
\begin{enumerate}
\item Tracer $\mathcal{C}_f$ et la droite $y=m$.
\item Identifier les \textbf{portions de courbe} situées \emph{strictement au-dessus} (pour $>$) ou \emph{au-dessus ou sur} (pour $\ge$) de la droite $y=m$.
\item Projeter ces portions sur l'axe des abscisses : on obtient un ou plusieurs intervalles.
\item \textbf{Écrire l'ensemble solution} en intervalle(s) :
\begin{itemize}
\item Parenthèse $]$ ou $[ $ si la borne correspond à une intersection avec $y=m$ (valeur exclue pour $>$, incluse pour $\ge$) ;
\item Crochet $[$ ou $]$ si la borne est une extrémité de $I$ \emph{et} que l'inéquation y est vérifiée.
\end{itemize}
\end{enumerate}
\end{methode}

\courssub{Étude du paramètre $m$ : Partition de l'ensemble image}

C'est le cœur du savoir-faire attendu au Probatoire : \emph{donner le nombre de solutions de $f(x)=m$ en fonction du paramètre $m$}.

\begin{propriete}[Nombre de solutions selon $m$ -- Polynôme degré 2 sur intervalle borné]
Soit $f(x)=ax^2+bx+c$ sur $I=[\alpha;\beta]$. Notons $x_0=-\frac{b}{2a}$ l'abscisse du sommet. On détermine l'ensemble image $f(I)=[m_{\min}; m_{\max}]$ à partir du tableau de variations.
\begin{itemize}
\item Si $m < m_{\min}$ ou $m > m_{\max}$ : \textbf{0 solution} (la droite $y=m$ est en dehors de l'ensemble image).
\item Si $m = m_{\min}$ ou $m = m_{\max}$ : \textbf{1 solution} si l'extremum est atteint en un point unique de $I$ (sommet intérieur ou borne unique), \textbf{2 solutions} si la valeur $m$ est atteinte à la fois au sommet (ou une borne) \emph{et} sur l'autre branche (cas $m$ égal à la valeur à l'autre borne).
\item Si $m_{\min} < m < m_{\max}$ : le nombre de solutions est \textbf{1 ou 2} selon que la droite $y=m$ coupe une ou deux branches de l'arc de parabole \emph{accessibles dans $I$}. Il faut comparer $m$ aux valeurs aux bornes $f(\alpha)$ et $f(\beta)$.
\end{itemize}
\end{propriete}

\begin{exemplebox}[{Exemple complet : $f(x)=-x^2+4x-1$ sur $I=[0;5]$}]
\emph{Étude préliminaire :} $f'(x)=-2x+4$. $f'(x)=0 \Leftrightarrow x=2 \in I$. $f(2)=3$ (maximum). $f(0)=-1$, $f(5)=-6$. Minimum sur $I$ : $m_{\min}=-6$ (en $x=5$). Image : $f(I)=[-6;3]$.

\medskip
\noindent\textbf{Équation $f(x)=m$ selon $m$ :}
\begin{itemize}
\item $m=4$ : $4>3=m_{\max}$ $\Rightarrow$ \textbf{0 solution} (droite au-dessus du sommet).
\item $m=3$ : $m=m_{\max}$ atteint en $x=2$ (intérieur) $\Rightarrow$ \textbf{1 solution} : $S=\{2\}$ (tangence).
\item $m=0$ : $-6<0<3$ et $0 \in ]-1;3[$ $\Rightarrow$ la droite coupe les deux branches $\Rightarrow$ \textbf{2 solutions}. Algébriquement : $-x^2+4x-1=0 \Leftrightarrow x^2-4x+1=0$, $\Delta=12$, $x=2\pm\sqrt{3}$. Les deux appartiennent à $[0;5]$.
\item $m=-1$ : $m=f(0)$ (borne gauche) et $-6 < -1 < 3$. La droite coupe en $x=0$ et sur la branche droite $\Rightarrow$ \textbf{2 solutions} ($x=0$ et $x=4$).
\item $m=-6$ : $m=m_{\min}$ atteint en $x=5$ (borne) $\Rightarrow$ \textbf{1 solution} : $S=\{5\}$.
\item $m=-7$ : $m<m_{\min}$ $\Rightarrow$ \textbf{0 solution}.
\end{itemize}

\medskip
\noindent\textbf{Inéquation $f(x)>0$ :} On cherche où la courbe est au-dessus de $Ox$. Les racines sont $2-\sqrt{3}\approx 0,27$ et $2+\sqrt{3}\approx 3,73$. Sur $[0;5]$, la courbe est au-dessus entre ces deux racines. $S = ]2-\sqrt{3}\,;\,2+\sqrt{3}[$.
\end{exemplebox}

\begin{popfigure}
\begin{tikzpicture}
\begin{axis}[
    width=\linewidth,
    height=8cm,
    axis lines=middle,
    xlabel=$x$, ylabel=$y$,
    xmin=-0.5, xmax=5.5,
    ymin=-7, ymax=4,
    xtick={0,1,2,3,4,5},
    ytick={-6,-4,-2,0,2,3},
    grid=major, grid style={gray!30},
    clip=false,
    domain=0:5,
    samples=200
]
% Parabole
\addplot[popcurve, thick, domain=0:5] {-x^2+4*x-1};
% Droites horizontales
\addplot[poporange, dashed, thick, domain=-0.5:5.5] {3} node[right, pos=0.95] {$y=3$ (1 sol)};
\addplot[popgreen, dashed, thick, domain=-0.5:5.5] {1} node[right, pos=0.95] {$y=1$ (2 sol)};
\addplot[poppink, dashed, thick, domain=-0.5:5.5] {-2} node[right, pos=0.95] {$y=-2$ (2 sol)};
\addplot[popblue, dashed, thick, domain=-0.5:5.5] {-6} node[right, pos=0.95] {$y=-6$ (1 sol)};
\addplot[poppurple, dashed, thick, domain=-0.5:5.5] {4} node[right, pos=0.95] {$y=4$ (0 sol)};
% Points d'intersection pour m=1
\addplot[only marks, mark=*, mark size=2.5, popgreen] coordinates {(0.268,1) (3.732,1)};
% Points pour m=3
\addplot[only marks, mark=*, mark size=2.5, poporange] coordinates {(2,3)};
% Points pour m=-6
\addplot[only marks, mark=*, mark size=2.5, popblue] coordinates {(5,-6)};
% Projections sur Ox pour m=1
\draw[popgreen, dashed, thin] (0.268,1) -- (0.268,0) node[below, font=\small] {$2-\sqrt{3}$};
\draw[popgreen, dashed, thin] (3.732,1) -- (3.732,0) node[below, font=\small] {$2+\sqrt{3}$};
% Sommet
\addplot[only marks, mark=*, mark size=3, popdark] coordinates {(2,3)};
\end{axis}
\end{tikzpicture}
\legende{Parabole $f(x)=-x^2+4x-1$ sur $[0;5]$ et droites $y=m$ pour divers $m$. En vert : 2 solutions ($m=1$). En orange : 1 solution par tangence au sommet ($m=3$). En bleu : 1 solution à la borne ($m=-6$). En violet : 0 solution ($m=4$). En rose : 2 solutions ($m=-2$).}
\end{popfigure}

\courssub{Cas des fonctions homographiques : Monotonie stricte}

Pour une fonction homographique $f(x)=\alpha+\frac{k}{x-\beta}$ sur un intervalle borné $I$ ne contenant pas $\beta$, la courbe est une branche d'hyperbole \emph{strictement monotone} (croissante si $k<0$, décroissante si $k>0$). Conséquence majeure :

\begin{propriete}[Nombre de solutions pour une homographique sur intervalle borné]
Soit $f(x)=\alpha+\frac{k}{x-\beta}$ définie sur $I=[a;b]$ avec $\beta\notin I$. $f$ est strictement monotone sur $I$. Pour tout $m\in\mathbb{R}$, l'équation $f(x)=m$ admet \textbf{au plus une solution} sur $I$ :
\begin{itemize}
\item \textbf{1 solution} si $m$ appartient à l'ensemble image $f(I)$ (intervalle ouvert ou fermé selon les bornes) ;
\item \textbf{0 solution} si $m\notin f(I)$.
\end{itemize}
\end{propriete}

\begin{exemplebox}[{Homographique : $f(x)=2+\frac{3}{x-1}$ sur $I=[2;4]$}]
$f'(x)=-\frac{3}{(x-1)^2}<0$ sur $I$ : $f$ est strictement décroissante.
$f(2)=2+\frac{3}{1}=5$, $f(4)=2+\frac{3}{3}=3$. Image : $f(I)=[3;5]$.
\begin{itemize}
\item $m=4 \in [3;5]$ $\Rightarrow$ \textbf{1 solution} : $2+\frac{3}{x-1}=4 \Rightarrow \frac{3}{x-1}=2 \Rightarrow x-1=\frac{3}{2} \Rightarrow x=2,5 \in I$.
\item $m=6 \notin [3;5]$ $\Rightarrow$ \textbf{0 solution}.
\item $m=3$ (borne) $\Rightarrow$ \textbf{1 solution} : $x=4$.
\end{itemize}
\end{exemplebox}

\begin{popfigure}
\begin{tikzpicture}
\begin{axis}[
    width=\linewidth,
    height=7cm,
    axis lines=middle,
    xlabel=$x$, ylabel=$y$,
    xmin=1.5, xmax=4.5,
    ymin=2, ymax=6,
    xtick={2,2.5,3,3.5,4},
    ytick={3,4,5},
    grid=major, grid style={gray!30},
    clip=false,
    domain=2:4,
    samples=200
]
% Hyperbole
\addplot[popcurve, thick, domain=2:4] {2+3/(x-1)};
% Droites horizontales
\addplot[popgreen, dashed, thick, domain=1.5:4.5] {4} node[right, pos=0.95] {$y=4$ (1 sol)};
\addplot[poporange, dashed, thick, domain=1.5:4.5] {6} node[right, pos=0.95] {$y=6$ (0 sol)};
\addplot[popblue, dashed, thick, domain=1.5:4.5] {3} node[right, pos=0.95] {$y=3$ (1 sol, borne)};
% Intersection m=4
\addplot[only marks, mark=*, mark size=3, popgreen] coordinates {(2.5,4)};
\draw[popgreen, dashed, thin] (2.5,4) -- (2.5,0) node[below, font=\small] {$2,5$};
% Bornes
\addplot[only marks, mark=*, mark size=3, popdark] coordinates {(2,5) (4,3)};
\end{axis}
\end{tikzpicture}
\legende{Hyperbole $f(x)=2+\frac{3}{x-1}$ sur $[2;4]$. Strictement décroissante : pour tout $m$, au plus une intersection. En vert : $m=4$ donne $x=2,5$. En orange : $m=6$ hors image, 0 solution. En bleu : $m=3$ donne la borne $x=4$.}
\end{popfigure}

\courssub{Résolution graphique des inéquations : Lecture par hachures}

\begin{exoresolu}[{Résolution graphique de $f(x) > 2$ pour $f(x)=-x^2+4x-1$ sur $[0;5]$}]
\emph{Énoncé :} Résoudre graphiquement l'inéquation $f(x) > 2$ sur $[0;5]$.

\medskip
\noindent\textbf{Solution détaillée :}
\begin{enumerate}
\item On trace $\mathcal{C}_f$ (parabole, sommet $(2;3)$, $f(0)=-1$, $f(5)=-6$).
\item On trace la droite $y=2$ (horizontal).
\item On repère les intersections : $-x^2+4x-1=2 \Leftrightarrow x^2-4x+3=0 \Leftrightarrow (x-1)(x-3)=0$. Abscisses : $x=1$ et $x=3$.
\item On identifie où la courbe est \textbf{strictement au-dessus} de $y=2$ : entre $x=1$ et $x=3$.
\item Comme l'inéquation est \emph{stricte} ($>$), les bornes $1$ et $3$ sont \textbf{exclues}. L'intervalle $I=[0;5]$ ne coupe pas la solution aux bornes.
\item Ensemble solution : $S = ]1;3[$.
\end{enumerate}

\medskip
\noindent\textbf{Vérification algébrique :} $-x^2+4x-1 > 2 \Leftrightarrow -x^2+4x-3 > 0 \Leftrightarrow x^2-4x+3 < 0$. Trinôme $a=1>0$, racines $1$ et $3$, négatif entre les racines : $]1;3[$. Cohérent.
\tcblower
\begin{tikzpicture}
\begin{axis}[
    width=\linewidth,
    height=6cm,
    axis lines=middle,
    xlabel=$x$, ylabel=$y$,
    xmin=-0.5, xmax=5.5,
    ymin=-2, ymax=4,
    xtick={0,1,2,3,4,5},
    ytick={2},
    grid=major, grid style={gray!30},
    clip=false,
    domain=0:5,
    samples=200
]
% Parabole
\addplot[popcurve, thick, domain=0:5] {-x^2+4*x-1};
% Droite y=2
\addplot[poporange, dashed, thick, domain=-0.5:5.5] {2} node[right, pos=0.95] {$y=2$};
% Hachure zone solution (au-dessus de y=2)
\addplot[popgreen!30, fill opacity=0.3, domain=1:3] {-x^2+4*x-1} \closedcycle;
\addplot[poporange, dashed, thick, domain=-0.5:5.5] {2};
% Intersections
\addplot[only marks, mark=*, mark size=3, poporange] coordinates {(1,2) (3,2)};
% Projections et hachure sur Ox
\draw[popgreen, line width=3pt] (1,0) -- (3,0) node[midway, above=2pt, font=\small\bfseries] {$S=]1;3[$};
\draw[popgreen, thick] (1,0.1) -- (1,-0.1) node[below, font=\small] {$1$};
\draw[popgreen, thick] (3,0.1) -- (3,-0.1) node[below, font=\small] {$3$};
% Cercles vides aux bornes
\draw[popgreen, thick] (1,2) circle (3pt);
\draw[popgreen, thick] (3,2) circle (3pt);
\end{axis}
\end{tikzpicture}
\legende{Résolution de $f(x)>2$ : la zone hachée en vert (courbe au-dessus de $y=2$) se projette sur l'intervalle $]1;3[$ (trait vert épais sur $Ox$). Les cercles vides rappellent l'exclusion des bornes (inéquation stricte).}
\end{exoresolu}

\begin{attention}[Écriture ensembliste rigoureuse : crochets vs parenthèses]
\begin{itemize}
\item \textbf{Intersection avec $y=m$ (équation ou borne d'inéquation stricte)} $\rightarrow$ \textbf{parenthèse} $]$ ou $[ $ (valeur exclue).
\item \textbf{Borne de l'intervalle de définition $I$ incluse dans la solution} $\rightarrow$ \textbf{crochet} $[$ ou $]$ (valeur incluse).
\item \textbf{Inéquation large ($\ge$ ou $\le$) et intersection} $\rightarrow$ \textbf{crochet} (valeur incluse).
\end{itemize}
\emph{Exemple :} $f(x)\ge 2$ sur $[0;5]$ donne $S=[1;3]$. $f(x)>2$ sur $[0;2]$ donne $S=]1;2]$ (borne $2$ incluse car $f(2)=3>2$ et $2\in I$).
\end{attention}

\courssub{Complément algébrique : Discriminant et signe du trinôme}

La méthode graphique donne la réponse ; l'algèbre la prouve et la précise.

\begin{propriete}[Discriminant et nombre de solutions de $f(x)=m$]
Pour $f(x)=ax^2+bx+c$ ($a\neq 0$), l'équation $f(x)=m$ s'écrit $ax^2+bx+(c-m)=0$. Son discriminant est $\Delta(m)=b^2-4a(c-m)=\Delta_0+4am$ où $\Delta_0=b^2-4ac$.
\begin{itemize}
\item $\Delta(m) < 0$ $\Leftrightarrow$ 0 solution réelle (sur $\mathbb{R}$, donc a fortiori sur $I$).
\item $\Delta(m) = 0$ $\Leftrightarrow$ 1 solution réelle (double) $x_0=-\frac{b}{2a}$. Elle appartient à $I$ \textbf{si et seulement si} $x_0\in I$.
\item $\Delta(m) > 0$ $\Leftrightarrow$ 2 solutions réelles distinctes $x_1,x_2$. Elles appartiennent à $I$ selon leur position par rapport aux bornes.
\end{itemize}
\end{propriete}

\begin{propriete}[Signe du trinôme et résolution de $f(x) > m$]
Pour $f(x)=ax^2+bx+c$ et $m\in\mathbb{R}$, soit $g(x)=f(x)-m=ax^2+bx+(c-m)$.
\begin{itemize}
\item Si $a>0$ (parabole tournée vers le haut) : $g(x)>0$ \textbf{en dehors} des racines (si $\Delta>0$), ou $\forall x\in\mathbb{R}$ (si $\Delta<0$), ou $\mathbb{R}\setminus\{x_0\}$ (si $\Delta=0$).
\item Si $a<0$ (parabole tournée vers le bas) : $g(x)>0$ \textbf{entre} les racines (si $\Delta>0$), ou $\varnothing$ (si $\Delta\le 0$).
\end{itemize}
\emph{Sur un intervalle borné $I$, on intersecte ces ensembles avec $I$.}
\end{propriete}

\begin{exemplebox}[{Vérification algébrique de l'exemple $f(x)=-x^2+4x-1$ sur $[0;5]$}]
$f(x)=m \Leftrightarrow -x^2+4x-1-m=0 \Leftrightarrow x^2-4x+(1+m)=0$.
$\Delta(m)=16-4(1+m)=12-4m=4(3-m)$.
\begin{itemize}
\item $\Delta(m)<0 \Leftrightarrow m>3$ : 0 solution (confirmé : $m_{\max}=3$).
\item $\Delta(m)=0 \Leftrightarrow m=3$ : 1 solution double $x=2\in I$ (confirmé : sommet).
\item $\Delta(m)>0 \Leftrightarrow m<3$ : 2 solutions réelles $x=2\pm\sqrt{3-m}$.
  \begin{itemize}
  \item Pour $m=0$ : $x=2\pm\sqrt{3}\approx 0,27$ et $3,73$ $\in [0;5]$ : 2 solutions.
  \item Pour $m=-6$ : $x=2\pm\sqrt{9}=2\pm 3 \Rightarrow -1$ et $5$. Seul $5\in I$ : 1 solution.
  \item Pour $m=-7$ : $x=2\pm\sqrt{10}\approx -1,16$ et $5,16$. Aucune dans $I$ : 0 solution ($m<m_{\min}$).
  \end{itemize}
\end{itemize}
\end{exemplebox}

\courssub{Méthode rigoureuse pour discuter le nombre de solutions selon $m$}

\begin{methode}[Discussion complète du nombre de solutions de $f(x)=m$ (Polynôme degré 2)]
\begin{enumerate}
\item \textbf{Étudier $f$ sur $I$ :} Tableau de variations, extremums $m_{\min}, m_{\max}$, valeurs aux bornes $f(\alpha), f(\beta)$, abscisse du sommet $x_0$.
\item \textbf{Partitionner $\mathbb{R}$ par rapport à $m_{\min}$ et $m_{\max}$ :} Intervales $]-\infty; m_{\min}[$, $\{m_{\min}\}$, $]m_{\min}; m_{\max}[$, $\{m_{\max}\}$, $]m_{\max}; +\infty[$.
\item \textbf{Analyser chaque cas :}
\begin{itemize}
\item \textbf{Hors image ($m < m_{\min}$ ou $m > m_{\max}$) :} 0 solution.
\item \textbf{Sur les bornes de l'image ($m = m_{\min}$ ou $m = m_{\max}$) :} Compter le nombre de préimages dans $I$ (1 si extremum au sommet intérieur ou borne unique ; 2 si la valeur est atteinte au sommet \emph{et} à une borne, ou aux deux bornes).
\item \textbf{À l'intérieur de l'image ($m_{\min} < m < m_{\max}$) :} Comparer $m$ aux valeurs aux bornes $f(\alpha)$ et $f(\beta)$.
  \begin{itemize}
  \item Si $m$ est \textbf{strictement entre} les deux valeurs aux bornes (ou égal à l'une si l'autre est l'extremum) : la droite coupe les deux branches accessibles $\Rightarrow$ \textbf{2 solutions}.
  \item Si $m$ est \textbf{entre un extremum intérieur et la valeur à la borne opposée} (cas intervalle asymétrique) : la droite ne coupe qu'une branche $\Rightarrow$ \textbf{1 solution}.
  \end{itemize}
\end{itemize}
\item \textbf{Rédiger la réponse finale} sous forme de tableau de valeurs ou phrases structurées couvrant tout $\mathbb{R}$.
\end{enumerate}
\end{methode}

\courssub{Exercices types Probatoire}

\begin{exoresolu}[Type Probatoire : Paramètre $m$ et inéquation]
\emph{Soit $f(x)=x^2-4x+2$ définie sur $I=[-1;4]$.}
\begin{enumerate}
\item Étudier les variations de $f$ sur $I$ et dresser son tableau de variations.
\item Déterminer l'ensemble image $f(I)$.
\item Discuter le nombre de solutions de l'équation $f(x)=m$ sur $I$ selon les valeurs de $m$.
\item Résoudre graphiquement l'inéquation $f(x) \le 1$ sur $I$.
\end{enumerate}

\medskip
\noindent\textbf{Solution :}
\begin{enumerate}
\item $f'(x)=2x-4$. $f'(x)=0 \Leftrightarrow x=2 \in I$. $f'(-1)=-6<0$, $f'(4)=4>0$.
  Tableau de variations :
  \begin{center}
  \begin{tabular}{|c|ccccc|}\hline
  $x$ & $-1$ & & $2$ & & $4$ \\ \hline
  $f'(x)$ & & $-$ & $0$ & $+$ & \\ \hline
  $f$ & $7$ & $\searrow$ & $-2$ & $\nearrow$ & $2$ \\ \hline
  \end{tabular}
  \end{center}
\item Minimum $m_{\min}=f(2)=-2$. Maximum $m_{\max}=f(-1)=7$ (car $f(-1)=7 > f(4)=2$). Image : $f(I)=[-2;7]$.
\item Discussion selon $m$ :
  \begin{itemize}
  \item $m < -2$ ou $m > 7$ : 0 solution.
  \item $m = -2$ : 1 solution ($x=2$, minimum intérieur).
  \item $m = 7$ : 1 solution ($x=-1$, maximum à la borne).
  \item $-2 < m < 2$ : 2 solutions (droite coupe les deux branches entre le minimum et la valeur à la borne droite $f(4)=2$).
  \item $m = 2$ : 2 solutions ? $f(x)=2 \Leftrightarrow x^2-4x=0 \Leftrightarrow x=0$ ou $x=4$. Les deux dans $I$ : \textbf{2 solutions} ($x=0$ et $x=4$).
  \item $2 < m < 7$ : 1 solution (seule la branche gauche est atteinte, la branche droite s'arrête à $f(4)=2$).
  \end{itemize}
  \textbf{Attention au cas $m=2$ :} bien que ce soit une valeur à la borne $x=4$, l'autre branche donne $x=0$, donc 2 solutions au total.
\item Inéquation $f(x) \le 1$ : On cherche où la courbe est \textbf{en-dessous ou sur} $y=1$.
  $f(x)=1 \Leftrightarrow x^2-4x+1=0 \Leftrightarrow \Delta=12$, $x=2\pm\sqrt{3}$.
  $2-\sqrt{3}\approx 0,27 \in I$, $2+\sqrt{3}\approx 3,73 \in I$.
  Comme $a=1>0$, $f(x)\le 1$ \textbf{entre} les racines. Intersection avec $I=[-1;4]$ :
  $S = [2-\sqrt{3}\,;\,2+\sqrt{3}] \subset I$.
\end{enumerate}
\tcblower
\begin{tikzpicture}
\begin{axis}[
    width=\linewidth,
    height=7cm,
    axis lines=middle,
    xlabel=$x$, ylabel=$y$,
    xmin=-1.5, xmax=4.5,
    ymin=-3, ymax=8,
    xtick={-1,0,1,2,3,4},
    ytick={-2,1,2,7},
    grid=major, grid style={gray!30},
    clip=false,
    domain=-1:4,
    samples=200
]
\addplot[popcurve, thick, domain=-1:4] {x^2-4*x+2};
\addplot[poporange, dashed, thick, domain=-1.5:4.5] {1} node[right, pos=0.95] {$y=1$};
\addplot[popgreen!30, fill opacity=0.3, domain=0.268:3.732] {x^2-4*x+2} \closedcycle;
\addplot[poporange, dashed, thick, domain=-1.5:4.5] {1};
\addplot[only marks, mark=*, mark size=3, poporange] coordinates {(0.268,1) (3.732,1)};
\draw[popgreen, line width=3pt] (0.268,0) -- (3.732,0) node[midway, above=2pt, font=\small\bfseries] {$S=[2-\sqrt{3};2+\sqrt{3}]$};
\draw[popgreen, thick] (0.268,0.1) -- (0.268,-0.1) node[below, font=\small] {$2-\sqrt{3}$};
\draw[popgreen, thick] (3.732,0.1) -- (3.732,-0.1) node[below, font=\small] {$2+\sqrt{3}$};
\addplot[only marks, mark=*, mark size=3, popdark] coordinates {(-1,7) (2,-2) (4,2)};
\end{axis}
\end{tikzpicture}
\legende{Résolution de $f(x)\le 1$ pour $f(x)=x^2-4x+2$ sur $[-1;4]$. Zone hachée = courbe sous $y=1$. Solution : intervalle fermé $[2-\sqrt{3};2+\sqrt{3}]$ (inéquation large $\Rightarrow$ bornes incluses).}
\end{exoresolu}

\begin{savaistu}[Le "m" du Probatoire : un classique]
Au Probatoire, l'exercice type présente une fonction (polynôme degré 2 ou homographique) sur un intervalle borné, demande le tableau de variations, puis : \og Discuter le nombre de solutions de $f(x)=m$ selon $m $\fg. La réponse attendue est un \textbf{tableau de valeurs} ou une \textbf{phrase structurée} donnant les intervalles de $m$ et le nombre de solutions associé. Ne vous contentez pas de dire \og ça dépend de $m $\fg : donnez la partition complète de $\mathbb{R}$ (ex : \og Pour $m\in]-2;2[$, 2 solutions ; pour $m=2$, 2 solutions ; pour $m\in]2;7[$, 1 solution ; etc.\fg). C'est là que se jouent les points de rigueur.
\end{savaistu}

\begin{exercice}[Entraînement : Homographique et paramètre]
Soit $f(x)=\frac{2x+1}{x-3}$ définie sur $I=[4;6]$.
\begin{enumerate}
\item Écrire $f$ sous la forme $\alpha+\frac{k}{x-\beta}$.
\item Étudier les variations de $f$ sur $I$ et donner $f(I)$.
\item Discuter le nombre de solutions de $f(x)=m$ sur $I$ selon $m$.
\item Résoudre $f(x) \ge 5$ sur $I$.
\end{enumerate}
\end{exercice}

\begin{corrige}[Exercice : Homographique et paramètre]
\begin{enumerate}
\item Division euclidienne : $2x+1 = 2(x-3)+7$. Donc $f(x)=2+\frac{7}{x-3}$. $\alpha=2$, $k=7$, $\beta=3$.
\item $k=7>0$, $\beta=3\notin I$. $f$ strictement décroissante sur $[4;6]$.
  $f(4)=2+\frac{7}{1}=9$, $f(6)=2+\frac{7}{3}=\frac{13}{3}\approx 4,33$.
  $f(I)=[\frac{13}{3};9]$.
\item Discussion :
  \begin{itemize}
  \item $m < \frac{13}{3}$ ou $m > 9$ : 0 solution.
  \item $m = \frac{13}{3}$ : 1 solution ($x=6$, borne).
  \item $m = 9$ : 1 solution ($x=4$, borne).
  \item $\frac{13}{3} < m < 9$ : 1 solution (monotonie stricte).
  \end{itemize}
  \emph{Résumé :} $\forall m\in f(I)$, 1 solution ; $\forall m\notin f(I)$, 0 solution.
\item $f(x) \ge 5 \Leftrightarrow 2+\frac{7}{x-3} \ge 5 \Leftrightarrow \frac{7}{x-3} \ge 3$.
  Sur $I$, $x-3>0$, on multiplie : $7 \ge 3(x-3) \Leftrightarrow 7 \ge 3x-9 \Leftrightarrow 16 \ge 3x \Leftrightarrow x \le \frac{16}{3}\approx 5,33$.
  Intersection avec $I=[4;6]$ : $S = [4;\frac{16}{3}]$.
  \emph{Vérification graphique :} droite $y=5$ coupe la courbe en $x=\frac{16}{3}$. Courbe au-dessus de $y=5$ pour $x\in[4;\frac{16}{3}]$.
\end{enumerate}
\end{corrige}

\begin{aretenir}[Check-list finale pour la résolution graphique]
\begin{itemize}
\item \textbf{Toujours} commencer par l'étude complète de $f$ sur $I$ (variations, extremums, ensemble image).
\item Pour $f(x)=m$ : tracer $y=m$, compter intersections $\rightarrow$ abscisses $\rightarrow$ ensemble $S$.
\item Pour $f(x) > m$ : hachurer zone \textbf{au-dessus} de $y=m$, projeter sur $Ox$, écrire intervalles (attention strict/large, bornes de $I$).
\item Paramètre $m$ : partitionner $\mathbb{R}$ d'après $m_{\min}$ et $m_{\max}$ (et valeurs aux bornes si monotonie par morceaux).
\item Vérifier algébriquement (discriminant, signe trinôme) si possible.
\item \textbf{Relire l'énoncé :} "sur $I$" signifie qu'on ne garde que les solutions \textbf{dans $I$}.
\end{itemize}
\end{aretenir}

\coursec{Synthèse et stratégie globale d'étude de fonction}

Après avoir maîtrisé, dans les sections précédentes, l'étude détaillée des fonctions polynômes du second degré et des fonctions homographiques sur un intervalle borné, ainsi que la construction des courbes des fonctions associées et la résolution graphique d'équations et d'inéquations, il est temps de rassembler toutes ces compétences dans une \textbf{stratégie unifiée}. Au Probatoire, un exercice d'étude de fonction ne se résout pas par morceaux disjoints ; il demande une démarche logique, rigoureuse et complète, où chaque étape prépare la suivante. Cette section te donne la \emph{feuille de route} infaillible pour ne rien oublier, gagner du temps et présenter une copie impeccable.

\courssub{L'algorithme complet : Les 9 étapes de l'étude standardisée}

Quelle que soit la fonction ($f$ polynôme de degré 2 ou homographique) et l'intervalle borné $I = [a ; b]$ donné, suis \textbf{systématiquement} cet ordre. C'est l'enchaînement logique : on ne peut pas tracer une courbe sans connaître ses variations, ni déterminer ses variations sans la dérivée (ou la forme canonique/standard), ni calculer la dérivée sans connaître l'ensemble de définition.

\begin{methode}[{Algorithme d'étude complète d'une fonction sur $I=[a;b]$}]
\begin{enumerate}
    \item \textbf{Ensemble de définition restreint :} Déterminer $D_f \cap I$.
    \begin{itemize}
        \item Polynôme degré 2 : $D_f = \mathbb{R}$, donc $D_f \cap I = I$.
        \item Homographique $f(x) = \frac{ax+b}{cx+d}$ : $D_f = \mathbb{R} \setminus \{-\frac{d}{c}\}$. Si la valeur interdite $x_0 = -\frac{d}{c}$ appartient à $I$, l'intervalle d'étude est \textbf{morcelé} en deux sous-intervalles : $[a ; x_0[$ et $]x_0 ; b]$. \emph{Note soigneusement la valeur interdite.}
    \end{itemize}
    \item \textbf{Forme repère (Identification des translations) :}
    \begin{itemize}
        \item \textbf{Degré 2 :} Mettre sous \textbf{forme canonique} $f(x) = a(x-\alpha)^2 + \beta$. Le sommet du parabole est $S(\alpha ; \beta)$. Vecteur de translation par rapport à $x \mapsto ax^2$ : $\vec{v}(\alpha ; \beta)$.
        \item \textbf{Homographique :} Effectuer la division euclidienne pour écrire sous \textbf{forme standard} $f(x) = \alpha + \frac{k}{x-\beta}$. Asymptotes : $x = \beta$ (verticale), $y = \alpha$ (horizontale). Centre de symétrie : $I(\beta ; \alpha)$. Vecteur de translation par rapport à $x \mapsto \frac{k}{x}$ : $\vec{v}(\beta ; \alpha)$.
    \end{itemize}
    \item \textbf{Dérivée et signe sur chaque sous-intervalle :} Calculer $f'(x)$.
    \begin{itemize}
        \item Dégré 2 : $f'(x) = 2a(x-\alpha)$. Signe trivial : opposé à $a$ avant $\alpha$, signe de $a$ après $\alpha$.
        \item Homographique : $f'(x) = -\frac{k}{(x-\beta)^2}$. Signe constant sur chaque intervalle (signe opposé à $k$).
    \end{itemize}
    \item \textbf{Tableau de variations :} Un tableau \textbf{par sous-intervalle} de définition. Colonnes : valeurs remarquables (bornes $a, b$, $\alpha$ ou $\beta$, valeurs interdites). Lignes : $x$, $f'(x)$ (signe), $f(x)$ (valeurs numériques exactes + flèches $\nearrow$/$\searrow$).
    \item \textbf{Extremums (absolus sur $I$) :} Lire sur le tableau de variations.
    \begin{itemize}
        \item Maximum absolu = plus grande valeur de $f$ sur $D_f \cap I$.
        \item Minimum absolu = plus petite valeur de $f$ sur $D_f \cap I$.
        \item \textbf{Attention :} Les bornes $a$ et $b$ sont \textbf{toujours candidates} aux extremums absolus sur un intervalle borné, même si la fonction n'y a pas d'extremum relatif.
    \end{itemize}
    \item \textbf{Intersections avec les axes (si utiles pour le tracé) :}
    \begin{itemize}
        \item Axe des ordonnées : $x=0 \in I \Rightarrow$ calculer $f(0)$.
        \item Axe des abscisses : Résoudre $f(x)=0$ sur $I$ (discriminant pour degré 2, numérateur pour homographique).
    \end{itemize}
    \item \textbf{Asymptotes (Homographique uniquement) :} Rappeler les équations $x=\beta$ et $y=\alpha$. Préciser la position relative des branches par rapport aux asymptotes (au-dessus/en-dessous, à gauche/à droite) grâce au signe de $k$ et au sens de variation.
    \item \textbf{Tracé de la courbe $\mathcal{C}_f$ :} Repère orthonormé, échelle adaptée, pointillés pour asymptotes, points pleins/creux aux bornes et valeur interdite, flèches aux extrémités des branches.
    \item \textbf{Réponses aux questions posées :} Résolution graphique de $f(x)=m$, $f(x)=0$, $f(x) > m$, nombre de solutions selon $m$, etc. \textbf{Encadrer} les solutions lues graphiquement.
\end{enumerate}
\end{methode}

\begin{attention}[Extremum relatif vs Extremum absolu sur $I$]
Ne confonds jamais ces deux notions !
\begin{itemize}
    \item \textbf{Extremum relatif (local) :} C'est un sommet de parabole ($\alpha \in I$) ou... \textbf{jamais} pour une homographique sur $\mathbb{R}$ (pas d'extremum relatif). Il se voit sur $\mathbb{R}$.
    \item \textbf{Extremum absolu (global) sur $I$ :} C'est la valeur max ou min \textbf{effectivement atteinte} sur l'intervalle \textbf{donné} $I$. Les bornes $a$ et $b$ sont \textbf{toujours} à tester. Exemple : $f(x)=x^2$ sur $I=[2;5]$. Pas d'extremum relatif dans $I$ (sommet en 0 hors intervalle). Mais minimum absolu $f(2)=4$, maximum absolu $f(5)=25$.
\end{itemize}
Sur un intervalle borné, \textbf{les bornes sont des candidates obligatoires} pour les extremums absolus.
\end{attention}

\begin{attention}[Oubli de la restriction à l'intervalle $I$]
L'erreur n°1 au Probatoire : Écrire $D_f = \mathbb{R}$ ou $D_f = \mathbb{R} \setminus \{\beta\}$ et faire le tableau de variations sur $\mathbb{R}$.
\begin{itemize}
    \item \textbf{Consigne :} "Étudier $f$ sur $I = [-2 ; 3]$".
    \item \textbf{Bon réflexe :} La première ligne de ta copie doit être : "On étudie $f$ sur $I = [-2 ; 3]$. $D_f \cap I = [-2 ; 3]$ (ou $[-2 ; \beta[ \cup ]\beta ; 3]$)".
    \item Le tableau de variations ne va que de $-2$ à $3$. Les flèches s'arrêtent aux bornes. Pas de $-\infty$ ni $+\infty$ en abscisse.
\end{itemize}
\end{attention}

\courssub{Deux approches complémentaires : Construction vs Justification}

Le programme attend la maîtrise de deux démarches qui ne s'opposent pas mais s'\textbf{articulent}.

\begin{propriete}[Les deux méthodes officielles]
\begin{enumerate}
    \item \textbf{Méthode par translations (Géométrique / Visuelle / Rapide) :}
    \begin{itemize}
        \item Utilise la forme canonique (deg 2) ou standard (homographique).
        \item Identifie le vecteur de translation $\vec{v}$ par rapport à la courbe de référence ($x \mapsto ax^2$ ou $x \mapsto \frac{k}{x}$).
        \item Trace la courbe de référence, puis applique la translation.
        \item \textbf{Idéale pour :} Construire la courbe rapidement (Étape 8), visualiser les asymptotes, le centre de symétrie, le sommet.
    \end{itemize}
    \item \textbf{Méthode par tableau de variations (Algébrique / Rigoureuse / Universelle) :}
    \begin{itemize}
        \item Utilise la dérivée $f'$ et son signe.
        \item Produit le tableau de variations complet.
        \item \textbf{Idéale pour :} Justifier rigoureusement les sens de variation (Étape 4), trouver les valeurs exactes des extremums (Étape 5), résoudre $f(x)=m$ par l'algèbre si besoin, garantir la cohérence du tracé.
    \end{itemize}
\end{enumerate}
\textbf{Stratégie gagnante :} Fais les \textbf{deux} en brouillon. La forme canonique/standard te donne la "photo" finale (courbe, asymptotes, sommet). Le tableau de variations te donne la "démonstration" (valeurs exactes, variations prouvées). Sur la copie, présente le \textbf{tableau de variations} comme justification principale, et trace la courbe en t'appuyant sur la vision translation.
\end{propriete}

\courssub{Articulation Algèbre / Graphique : La cohérence obligatoire}

Le passage de l'algèbre (tableau) au graphique (courbe) et inversement est le cœur de la compétence "Représentation".

\begin{itemize}
    \item \textbf{Tableau $\to$ Courbe :} Le tableau donne les points "clés" (bornes, extremums, intersections axes) et la \textbf{forme} des branches (croissante/décroissante, convexe/concave — bien que convexité non au programme 1ère A, la forme de la parabole ou de l'hyperbole est connue). \emph{Ne trace jamais "au feeling" : chaque segment de courbe correspond à une ligne du tableau.}
    \item \textbf{Courbe $\to$ Lecture graphique :} Une fois la courbe tracée \textbf{avec précision} (échelle respectée, points clés placés exactement), elle devient un outil de mesure.
    \begin{itemize}
        \item $f(x) = m$ : Intersections avec la droite $y=m$. Compter les points.
        \item $f(x) > m$ : Portions de courbe \textbf{strictement au-dessus} de $y=m$. Lire les abscisses sur l'axe $Ox$.
        \item Nombre de solutions selon $m$ : Analyser la position de la droite $y=m$ par rapport aux extremums absolus et aux asymptotes.
    \end{itemize}
    \item \textbf{Cohérence :} Si ton tableau dit "max 5 en $x=1$" et que ton courbe passe visiblement au-dessus de $y=5$, \textbf{il y a une erreur}. Vérifie tes calculs d'images ($f(a), f(b), f(\alpha)$) et tes échelles.
\end{itemize}

\begin{methode}[Tracé rapide d'une homographique $f(x) = \alpha + \frac{k}{x-\beta}$]
En 30 secondes, tu as le squelette parfait :
\begin{enumerate}
    \item Trace les asymptotes : $x=\beta$ (verticale, pointillés), $y=\alpha$ (horizontale, pointillés). Elles se coupent en $I(\beta ; \alpha)$.
    \item Place le centre $I$.
    \item Choisis deux abscisses faciles de part et d'autre de $\beta$ : $x_1 = \beta - 1$ et $x_2 = \beta + 1$ (si dans $I$).
    \item Calcule les images : $f(\beta-1) = \alpha - k$, $f(\beta+1) = \alpha + k$.
    \item Trace les deux points $A(\beta-1 ; \alpha-k)$ et $B(\beta+1 ; \alpha+k)$.
    \item Trace les deux branches hyperboliques passant par $A$ et $B$, asymptotes aux droites, \textbf{en respectant le sens de variation} (décroissante si $k>0$, croissante si $k<0$ sur chaque branche).
\end{enumerate}
\end{methode}

\courssub{Gestion des cas limites et pièges classiques}

\begin{itemize}
    \item \textbf{Valeur interdite aux bornes de $I$ :} Si $\beta = a$ ou $\beta = b$ pour une homographique. L'intervalle d'étude n'est pas morcelé (ex: $I=[\beta ; b]$). L'asymptote verticale $x=\beta$ coïncide avec une borne. Le tableau de variations se fait sur $]\beta ; b]$. La courbe a une branche unique. \textbf{Point creux} à la borne $\beta$ (valeur non atteinte).
    \item \textbf{Sommet aux bornes (Deg 2) :} Si $\alpha = a$ ou $\alpha = b$. La fonction est monotone sur $I$ (strictement croissante ou décroissante). L'extremum absolu (min ou max selon $a$) est atteint \textbf{à la borne} qui est le sommet. L'autre extremum est à l'autre borne.
    \item \textbf{Intervalle réduit à un point :} Exclu par définition d'un intervalle borné ($a < b$).
    \item \textbf{Simplification d'homographique (Trou) :} Si $f(x) = \frac{(x-x_0)P(x)}{Q(x)}$ avec $Q(x_0)=0$ et simplification possible. $f$ se prolonge par continuité. \textbf{Étudier la fonction simplifiée} sur $I$, mais \textbf{noter le trou} (point creux) en $x_0$ sur la courbe et dans le tableau de variations (valeur interdite). Voir l'exemple complet ci-dessous.
\end{itemize}

\courssub{Rédaction mathématique : Le langage du Probatoire}

Une copie de qualité utilise un vocabulaire précis. Bannis le langage approximatif.

\begin{center}
\begin{tabularx}{\linewidth}{|l|X|X|}
\hline
\rowcolor{popblueL}
\textbf{À éviter (Familier / Flou)} & \textbf{À écrire (Rigoureux / Attendu)} & \textbf{Pourquoi ?} \\
\hline
"La fonction monte" & "La fonction $f$ est \textbf{strictement croissante} sur $[a;b]$" & Terminologie officielle. \\
\hline
"Le max c'est 5" & "Le \textbf{maximum absolu} de $f$ sur $I$ est $5$, atteint en $x=2$" & Précision : valeur + abscisse + nature (absolu). \\
\hline
"L'asymptote c'est $x=2$" & "La droite d'équation $x=2$ est une \textbf{asymptote verticale} (parallèle à l'axe des ordonnées) à $\mathcal{C}_f$" & Définition complète. \\
\hline
"On décale le graphique" & "La courbe $\mathcal{C}_f$ s'obtient par \textbf{translation} de vecteur $\vec{v}(\alpha ; \beta)$ de la courbe de référence" & Vocabulaire transformation. \\
\hline
"Il y a 2 solutions" & "L'équation $f(x)=m$ admet \textbf{deux solutions distinctes} sur $I$, encadrées par $x_1 \in ]1,2; 1,3[$ et $x_2 \in ]2,8; 2,9[$" & Encadrement obligatoire en lecture graphique. \\
\hline
\end{tabularx}
\end{center}

\courssub{Exemple complet corrigé : Le cas "piège" de la simplification}

Étudions la fonction $f(x) = \frac{x^2 - 2x - 3}{x+1}$ sur l'intervalle $I = [-2 ; 2]$.

\begin{exoresolu}[Étude complète avec simplification et trou]
\textbf{Énoncé :} Soit $f(x) = \frac{x^2 - 2x - 3}{x+1}$. Étudier $f$ sur $I = [-2 ; 2]$. Tracer $\mathcal{C}_f$. Résoudre graphiquement $f(x) = -2$.

\textbf{Solution détaillée :}

\textbf{1. Ensemble de définition restreint.}
$D_f = \mathbb{R} \setminus \{-1\}$. Sur $I = [-2 ; 2]$, la valeur interdite $-1$ appartient à $I$.
L'étude se fait sur deux sous-intervalles : $I_1 = [-2 ; -1[$ et $I_2 = ]-1 ; 2]$.

\textbf{2. Simplification et Forme standard.}
Numérateur : $x^2 - 2x - 3 = (x-3)(x+1)$.
Pour $x \neq -1$ : $f(x) = \frac{(x-3)\cancel{(x+1)}}{\cancel{x+1}} = x - 3$.
$f$ coïncide avec la fonction affine $g(x) = x - 3$ sur $D_f$.
\emph{Attention :} $f$ n'est \textbf{pas} définie en $-1$. C'est un \textbf{trou} (discontinuité évitable).
La courbe de $f$ est donc la droite d'équation $y=x-3$, privée du point d'abscisse $-1$ : il n'y a pas d'asymptote (ni verticale, ni horizontale), seulement un point manquant sur la droite.
\textit{Note :} Au programme 1ère A, on traite les asymptotes parallèles aux axes. Ici, pas d'asymptote parallèle aux axes. La courbe est une droite percée.

\textbf{3. Dérivée et signe.}
$f'(x) = 1$ sur $D_f$.
$f'(x) > 0$ sur $I_1$ et $I_2$.
$f$ est \textbf{strictement croissante} sur $[-2 ; -1[$ et sur $]-1 ; 2]$.

\textbf{4. Tableaux de variations (un par sous-intervalle).}

\emph{Sur $I_1 = [-2 ; -1[$ :}
\begin{center}
\begin{tabular}{|c|ccc|}\hline
$x$ & $-2$ & & $-1$ \\ \hline
$f'(x)$ & & $+$ & \\ \hline
$f(x)$ & $-5$ & $\nearrow$ & \multicolumn{1}{c|}{$\begin{array}{c}\text{Valeur}\\\text{interdite}\\(-4)\end{array}$} \\ \hline
\end{tabular}
\end{center}

\emph{Sur $I_2 = ]-1 ; 2]$ :}
\begin{center}
\begin{tabular}{|c|ccc|}\hline
$x$ & $-1$ & & $2$ \\ \hline
$f'(x)$ & & $+$ & \\ \hline
$f(x)$ & \multicolumn{1}{c|}{$\begin{array}{c}\text{Valeur}\\\text{interdite}\\(-4)\end{array}$} & $\nearrow$ & $-1$ \\ \hline
\end{tabular}
\end{center}
\emph{Légende :} La valeur $-4$ (image de $-1$ par $g$) est notée entre parenthèses pour indiquer le trou. Les flèches s'arrêtent avant/après le trou.

\textbf{5. Extremums absolus sur $I$.}
$f$ croissante par morceaux.
Minimum absolu : $f(-2) = -5$ (atteint en $-2$).
Maximum absolu : $f(2) = -1$ (atteint en $2$).
\emph{Aucun extremum en $-1$ car $f(-1)$ n'existe pas.}

\textbf{6. Intersections axes.}
$y = f(0) = -3$ $\Rightarrow$ Point $B(0 ; -3)$.
$f(x) = 0 \Leftrightarrow x - 3 = 0 \Leftrightarrow x = 3 \notin I$. Pas d'intersection avec $Ox$ sur $I$.

\textbf{7. Asymptotes.}
Aucune asymptote parallèle aux axes. (Asymptote oblique $y=x-3$ non demandée).

\textbf{8. Tracé de $\mathcal{C}_f$.}
\begin{popfigure}
\begin{tikzpicture}[scale=0.8, >=Stealth]
% Axes
\draw[->] (-3,0) -- (3,0) node[right] {$x$};
\draw[->] (0,-6) -- (0,1) node[above] {$y$};
% Graduations
\foreach \x in {-2,-1,1,2} \draw (\x,0.1) -- (\x,-0.1) node[below] {$\x$};
\foreach \y in {-5,-4,-3,-2,-1} \draw (0.1,\y) -- (-0.1,\y) node[left] {$\y$};
% Droite y = x-3
\draw[popblue, thick, domain=-2.5:2.5] plot (\x, \x-3);
% Trou en (-1, -4)
\draw[popblue, thick, fill=white] (-1,-4) circle (3pt);
% Points aux bornes
\draw[popblue, thick, fill=popblue] (-2,-5) circle (2.5pt) node[below left] {$A(-2;-5)$};
\draw[popblue, thick, fill=popblue] (2,-1) circle (2.5pt) node[above right] {$B(2;-1)$};
% Point intersection Oy
\draw[popblue, thick, fill=popblue] (0,-3) circle (2.5pt) node[right] {$C(0;-3)$};
% Intervalle I
\draw[poporange, ultra thick] (-2,-5.5) -- (2,-5.5) node[midway, below] {$I = [-2;2]$};
\end{tikzpicture}
\legende{Courbe $\mathcal{C}_f$ : droite $y=x-3$ avec un trou en $(-1;-4)$. Bornes $A$ et $B$ pleines.}
\end{popfigure}

\textbf{9. Résolution graphique de $f(x) = -2$.}
On trace la droite $y = -2$ (pointillés). Intersection avec $\mathcal{C}_f$ (la droite pleine) : un unique point d'abscisse $x = 1$.
Vérification algébrique : $x - 3 = -2 \Rightarrow x = 1 \in I$.
\textbf{Réponse :} L'équation admet une unique solution $x=1$ sur $I$.
\end{exoresolu}

\courssub{Entraînement croisé : Modélisation, Transformations et Paramètre}

Les sujets de Probatoire mélangent souvent tout le programme. Voici le type de problème "synthèse" à maîtriser.

\begin{exemplebox}[Problème type Probatoire : Bénéfice d'une coopérative agricole]
Une coopérative de caféiculteurs du Littoral modélise son bénéfice mensuel (en milliers de FCFA) en fonction de la quantité $x$ (en tonnes) vendue, par la fonction $f(x) = -2x^2 + 12x - 10$ sur l'intervalle $I = [1 ; 5]$.

\begin{enumerate}
    \item Étudier $f$ sur $I$ (tableau de variations, extremums, courbe).
    \item La coopérative veut un bénéfice d'au moins 8 millions FCFA. Déterminer graphiquement les quantités à vendre.
    \item On considère la fonction $g(x) = |f(x)|$. Construire $\mathcal{C}_g$ à partir de $\mathcal{C}_f$. Interpréter $g(x)$ dans le contexte.
\end{enumerate}

\textbf{Guidage pour la résolution :}
\begin{enumerate}
    \item \textbf{Forme canonique :} $f(x) = -2(x^2 - 6x) - 10 = -2[(x-3)^2 - 9] - 10 = -2(x-3)^2 + 18 - 10 = -2(x-3)^2 + 8$.
    Sommet $S(3 ; 8)$. $a=-2 < 0 \Rightarrow$ concave, max en $S$.
    \textbf{Tableau sur $I=[1;5]$ :} $f(1)=0$, $f(3)=8$, $f(5)=-2(4)+8=0$.
    Variations : $\nearrow$ sur $[1;3]$, $\searrow$ sur $[3;5]$.
    Max absolu = 8 (en $x=3$). Min absolu = 0 (en $x=1$ et $x=5$).
    \item \textbf{Bénéfice $\ge 8$ :} $f(x) \ge 8$. Sur le graphique, portion au-dessus de $y=8$. Comme max = 8 atteint seulement en $x=3$, la solution est \textbf{$x = 3$ tonnes exactement}.
    \item \textbf{Fonction $g(x) = |f(x)|$ :} $f(x) \ge 0$ sur $[1;5]$ (min = 0). Donc $g(x) = f(x)$ sur tout $I$. $\mathcal{C}_g = \mathcal{C}_f$. Interprétation : Le bénéfice est toujours positif ou nul sur cette plage de production, la valeur absolue ne change rien. Si l'intervalle était $[0;6]$, $f$ deviendrait négatif après 5, et $g$ refléterait la perte en gain (valeur absolue du déficit).
\end{enumerate}
\end{exemplebox}

\courssub{Fiche de révision visuelle (Mindmap)}

Voici une carte mentale pour réviser la structure globale. Visualise-la avant l'examen.

\begin{popfigure}
\begin{tikzpicture}[
    mindmap,
    grow cyclic,
    every node/.style=concept,
    concept color=popblue,
    level 1/.append style={level distance=4.5cm, sibling angle=90},
    level 2/.append style={level distance=3cm, sibling angle=45},
    level 3/.append style={level distance=2.5cm, sibling angle=30}
]
\node[concept, text width=2.5cm, font=\bfseries] {Étude Fonction\\ sur $I=[a;b]$}
    child[concept color=popgreen] { node[concept, text width=2.2cm] {Fonctions\\ concernées}
        child { node[concept, text width=1.8cm] {Polynôme\\ Deg 2 : $ax^2+bx+c$} }
        child { node[concept, text width=1.8cm] {Homographique\\ $\frac{ax+b}{cx+d}$} }
    }
    child[concept color=poporange] { node[concept, text width=2.2cm] {Outils\\ Algébriques}
        child { node[concept, text width=1.8cm] {Forme Canonique\\ $a(x-\alpha)^2+\beta$} }
        child { node[concept, text width=1.8cm] {Forme Standard\\ $\alpha+\frac{k}{x-\beta}$} }
        child { node[concept, text width=1.8cm] {Dérivée $f'$\\ \& Signe} }
    }
    child[concept color=poppurple] { node[concept, text width=2.2cm] {Productions\\ Obligatoires}
        child { node[concept, text width=1.8cm] {$D_f \cap I$\\ (Morcellement)} }
        child { node[concept, text width=1.8cm] {Tableau de\\ Variations} }
        child { node[concept, text width=1.8cm] {Extremums\\ Absolus sur $I$} }
        child { node[concept, text width=1.8cm] {Asymptotes\\ (Homographique)} }
        child { node[concept, text width=1.8cm] {Courbe $\mathcal{C}_f$\\ Soignée} }
    }
    child[concept color=poppink] { node[concept, text width=2.2cm] {Exploitation\\ Graphique}
        child { node[concept, text width=1.8cm] {$f(x)=m$\\ Nb solutions} }
        child { node[concept, text width=1.8cm] {$f(x) > m$\\ Ensemble sol.} }
        child { node[concept, text width=1.8cm] {Fonctions\\ Associées\\ $-f, |f|, f(x-a)+b$} }
    };
\end{tikzpicture}
\legende{Mindmap de synthèse : Structure globale de l'étude de fonction en 1ère A.}
\end{popfigure}

\courssub{Organigramme de la stratégie d'examen (Flowchart)}

Garde cet organigramme en tête pendant l'épreuve. Il gère les embranchements (homographique vs polynôme, valeur interdite dans $I$ ou non).

\begin{popfigure}
\begin{tikzpicture}[
    node distance=1.2cm and 2cm,
    startstop/.style={rectangle, rounded corners, minimum width=3cm, minimum height=1cm, text centered, draw=popblue, fill=popblueL, thick},
    process/.style={rectangle, minimum width=3.5cm, minimum height=1cm, text centered, draw=popgreen, fill=popgreenL, thick},
    decision/.style={diamond, minimum width=3cm, minimum height=1.5cm, text centered, draw=poporange, fill=poporangeL, thick, aspect=2},
    arrow/.style={thick,->,>=Stealth, popdark}
]
% Nodes
\node (start) [startstop] {DÉBUT : Expression $f(x)$, Intervalle $I=[a;b]$};
\node (def) [process, below of=start] {1. Déterminer $D_f \cap I$};
\node (dec1) [decision, below of=def, align=center]{$f$ homographique ET\\ valeur interdite $\in I$ ?};
\node (morcel) [process, left of=dec1, xshift=-3.5cm, align=center]{Morceler $I$ en $I_1, I_2$\\ (borne exclue)};
\node (form) [process, right of=dec1, xshift=3.5cm, align=center]{2. Forme Canonique (deg 2)\\ ou Standard (homog.)\\ $\to$ Translations, Asymptotes};
\node (der) [process, below of=dec1, yshift=-1.5cm, align=center]{3. Calculer $f'(x)$ et signe\\ sur CHAQUE sous-intervalle};
\node (tab) [process, below of=der, align=center]{4. Tableau de variations\\ (un par sous-intervalle)};
\node (ext) [process, below of=tab, align=center]{5. Extremums ABSOLUS sur $I$\\ (Bornes + Extremums relatifs internes)};
\node (inter) [process, below of=ext, align=center]{6. Intersections axes\\ (si simple/utile)};
\node (asym) [decision, below of=inter] {Homographique ?};
\node (drawasym) [process, left of=asym, xshift=-3cm, align=center]{7. Tracer Asymptotes\\ $x=\beta, y=\alpha$};
\node (curve) [process, below of=asym, align=center]{8. TRACER $\mathcal{C}_f$\\ (Repère, échelles, points pleins/creux)};
\node (solve) [process, below of=curve, align=center]{9. RÉPONDRE aux questions\\ (Équations, Inéquations, $m$)\\ \textbf{Encadrer solutions graphiques}};
\node (end) [startstop, below of=solve, align=center]{FIN : Relecture cohérence\\ Tableau $\leftrightarrow$ Courbe};

% Arrows
\draw [arrow] (start) -- (def);
\draw [arrow] (def) -- (dec1);
\draw [arrow] (dec1) -- node[anchor=east, pos=0.5, font=\small] {OUI} (morcel);
\draw [arrow] (morcel) |- (form);
\draw [arrow] (dec1) -- node[anchor=west, pos=0.5, font=\small] {NON} (form);
\draw [arrow] (form) -- (der);
\draw [arrow] (der) -- (tab);
\draw [arrow] (tab) -- (ext);
\draw [arrow] (ext) -- (inter);
\draw [arrow] (inter) -- (asym);
\draw [arrow] (asym) -- node[anchor=south, pos=0.5, font=\small] {OUI} (drawasym);
\draw [arrow] (drawasym) |- (curve);
\draw [arrow] (asym) -- node[anchor=north, pos=0.5, font=\small] {NON} (curve);
\draw [arrow] (curve) -- (solve);
\draw [arrow] (solve) -- (end);
\end{tikzpicture}
\legende{Organigramme de la stratégie d'étude complète (Flowchart).}
\end{popfigure}

\courssub{Complément Terminale (Convexité) — Culture mathématique}

\begin{savaistu}[Convexité et Dérivée seconde ($f''$) — Hors programme 1ère A, Base Terminale]
En 1ère A, on utilise la forme de la courbe (parabole, hyperbole) connue a priori. En Terminale, on \textbf{prouve} cette forme par la dérivée seconde $f''$.
\begin{itemize}
    \item $f''(x) > 0$ sur $I \Rightarrow f$ \textbf{convexe} sur $I$ (courbe au-dessus de ses tangentes, "creuse vers le haut").
    \item $f''(x) < 0$ sur $I \Rightarrow f$ \textbf{concave} sur $I$ (courbe en-dessous de ses tangentes, "creuse vers le bas").
\end{itemize}
\textbf{Applications aux fonctions du programme :}
\begin{itemize}
    \item $f(x) = ax^2 + bx + c \Rightarrow f''(x) = 2a$. Signe constant.
    \begin{itemize}
        \item $a > 0 \Rightarrow f'' > 0 \Rightarrow$ \textbf{Convexe} (parabole tournée vers le haut).
        \item $a < 0 \Rightarrow f'' < 0 \Rightarrow$ \textbf{Concave} (parabole tournée vers le bas).
    \end{itemize}
    \item $f(x) = \alpha + \frac{k}{x-\beta} \Rightarrow f'(x) = -\frac{k}{(x-\beta)^2} \Rightarrow f''(x) = \frac{2k}{(x-\beta)^3}$.
    Le signe de $f''$ dépend du signe de $k$ ET du côté de l'asymptote ($x-\beta$).
    \begin{itemize}
        \item Branche de droite ($x > \beta$) : signe de $f''$ = signe de $k$.
        \item Branche de gauche ($x < \beta$) : signe de $f''$ = signe opposé de $k$.
    \end{itemize}
    \emph{Exemple :} $f(x) = \frac{1}{x}$ ($k=1>0$). Branche droite ($x>0$) : $f''>0$ (Convexe). Branche gauche ($x<0$) : $f''<0$ (Concave). Cela explique la forme caractéristique de l'hyperbole.
\end{itemize}
\textbf{Pourquoi c'est utile maintenant ?} Cela justifie \textbf{rigoureusement} pourquoi tu traces les branches "courbées" correctement (pas des segments de droite !), même sans calculer $f''$.
\end{savaistu}

\begin{savaistu}[La dichotomie : De la lecture graphique à l'algorithme numérique]
La résolution graphique $f(x)=m$ (intersection courbe / droite horizontale) est l'ancêtre visuel de la \textbf{méthode de dichotomie} (ou bisection), algorithme fondamental en informatique et analyse numérique pour approcher les solutions d'équations impossibles à résoudre exactement (ex: $x^3+x-1=0$, $e^x = 5x$, $\cos x = x$).
\begin{enumerate}
    \item On cherche un intervalle $[a;b]$ où $f$ est continue et $f(a)$ et $f(b)$ ont des signes contraires (Théorème des valeurs intermédiaires $\Rightarrow$ au moins une solution).
    \item On calcule le milieu $m = \frac{a+b}{2}$.
    \item Si $f(m)=0$, c'est gagné. Sinon, on regarde le signe de $f(m)$ : il a le même signe que $f(a)$ ou $f(b)$. On remplace la borne correspondante par $m$.
    \item On recommence. L'intervalle se divise par 2 à chaque étape. En 10 étapes, précision $\approx 10^{-3}$ ; en 20 étapes, $\approx 10^{-6}$.
\end{enumerate}
C'est exactement ce que fait ta calculatrice (ou Python, Matlab, les GPS, les simulateurs financiers) quand tu lui demandes "Solve". Ton travail graphique en 1ère A, c'est \textbf{l'initialisation} de cet algorithme : trouver l'intervalle de départ $[a;b]$ où la solution est piégée.
\end{savaistu}

\begin{aretenir}[Check-list finale pour l'examen (À mémoriser)]
Avant de rendre ta copie, coche mentalement chaque case :
\begin{itemize}
    \item $\Box$ $D_f \cap I$ écrit explicitement (morcellement si homog + valeur interdite dans $I$).
    \item $\Box$ Forme canonique (deg 2) ou Standard (homog) \textbf{avec identification du vecteur de translation}.
    \item $\Box$ Dérivée $f'(x)$ calculée et signe étudié sur \textbf{chaque} sous-intervalle.
    \item $\Box$ Tableau de variations : \textbf{3 lignes} ($x$, $f'$, $f$), \textbf{même nombre de colonnes}, bornes $a,b$, valeurs remarquables ($\alpha, \beta$), flèches $\nearrow/\searrow$ UNIQUEMENT sur ligne $f$, valeurs exactes (fractions, radicaux).
    \item $\Box$ Extremums absolus : Valeur + Abscisse + "absolu sur $I$". Bornes vérifiées.
    \item $\Box$ Asymptotes (homog) : Équations + Position branches (au-dessus/en-dessous).
    \item $\Box$ Courbe : Repère orthonormé, échelles, pointillés asymptotes, points pleins (atteints) / creux (interdits/non atteints), flèches aux bouts.
    \item $\Box$ Réponses questions : Solutions encadrées (ex: $x \in ]1,2; 1,3[$), notation ensembliste pour inéquations.
\end{itemize}
\end{aretenir}

\courssub{Exercices d'entraînement synthétique}

\begin{exercice}[Entraînement 1 : Homographique avec paramètre]
Soit $f_m(x) = \frac{2x + m}{x - 1}$ définie sur $I = [0 ; 3]$, $m \in \mathbb{R}$.
\begin{enumerate}
    \item Déterminer $m$ pour que $f_m$ ait une asymptote horizontale $y = 3$.
    \item Pour cette valeur de $m$, étudier $f_m$ sur $I$ (tableau, courbe).
    \item Résoudre graphiquement $f_m(x) = 4$ sur $I$.
\end{enumerate}
\end{exercice}

\begin{exercice}[Entraînement 2 : Polynôme degré 2 et fonction associée]
Soit $f(x) = -x^2 + 4x - 1$ sur $I = [0 ; 5]$.
\begin{enumerate}
    \item Étudier $f$ (forme canonique, tableau, extremums, courbe).
    \item Construire la courbe de $g(x) = f(x-2) + 3$. Préciser les transformations.
    \item Déterminer graphiquement le nombre de solutions de $g(x) = 5$ sur l'intervalle image de $I$ par la translation.
\end{enumerate}
\end{exercice}

\begin{corrige}[Entraînement 1]
\begin{enumerate}
    \item On effectue la division : $f_m(x) = \frac{2x+m}{x-1} = 2 + \frac{m+2}{x-1}$.
    L'asymptote horizontale est $y = 2$ (quotient des coefficients directeurs), \textbf{indépendante de $m$}.
    Il n'existe \textbf{aucun} réel $m$ tel que l'asymptote horizontale soit $y=3$.
    \emph{Remark :} Si l'énoncé prévoyait un passage par un point donné (ex: $A(0;5)$), on aurait $f_m(0)=-m=5 \Rightarrow m=-5$.
    \item Comme la question 1 n'a pas de solution, on ne peut pas poursuivre "pour cette valeur de $m$".
    \textbf{Entraînement proposé :} Étudier $f_0(x) = \frac{2x}{x-1} = 2 + \frac{2}{x-1}$ sur $I=[0;3]$.
    $D_f \cap I = [0;1[ \cup ]1;3]$. $f'(x) = -\frac{2}{(x-1)^2} < 0$ sur les deux branches $\Rightarrow$ strictement décroissante.
    Asymptotes : $x=1$ (verticale), $y=2$ (horizontale).
    $f(0)=0$, $f(3)=3$.
    Branche gauche ($[0;1[$) : $f$ décroît de $f(0)=0$ vers $-\infty$ quand $x\to 1^-$ : valeurs $\in\,]-\infty ; 0]$.
    Branche droite ($]1;3]$) : $f$ décroît de $+\infty$ (quand $x\to 1^+$) vers $f(3)=3$ : valeurs $\in [3 ; +\infty[$.
    \item Pour $f_0$, équation $f_0(x)=4 \Leftrightarrow 2 + \frac{2}{x-1}=4 \Leftrightarrow \frac{2}{x-1}=2 \Leftrightarrow x-1=1 \Leftrightarrow x=2$.
    Vérification : $x=2 \in I$ et $x=2\neq 1$, donc $x=2$ est bien dans le domaine d'étude.
    Graphiquement : la droite $y=4$ coupe la branche gauche ($[0;1[$, valeurs $\in\,]-\infty;0]$) nulle part car $4\notin\,]-\infty;0]$. Elle coupe la branche droite ($]1;3]$, valeurs $\in [3;+\infty[$) car $4\in[3;+\infty[$, au point d'abscisse $x=2$.
    \textbf{Il y a exactement 1 solution sur $I$ pour $f_0$ : $x=2$.}
\end{enumerate}
\end{corrige}

\begin{corrige}[Entraînement 2]
\begin{enumerate}
    \item $f(x) = -(x^2 - 4x) - 1 = -[(x-2)^2 - 4] - 1 = -(x-2)^2 + 3$. Sommet $S(2;3)$, $a=-1<0$ (concave, max).
    $I=[0;5]$. $f(0)=-1$, $f(2)=3$, $f(5)=-25+20-1=-6$.
    Tableau : $0 \nearrow 2 \searrow 5$ ; $f: -1 \nearrow 3 \searrow -6$.
    Max absolu = 3 (en 2). Min absolu = -6 (en 5).
    \item $g(x) = f(x-2) + 3$. Translation vecteur $\vec{v}(2;3)$.
    $\mathcal{C}_g$ : Sommet $S'(4;6)$. Intervale $I$ translaté : $[2;7]$.
    $g(x) = -((x-2)-2)^2 + 3 + 3 = -(x-4)^2 + 6$.
    \item $g(x)=5 \Leftrightarrow -(x-4)^2 + 6 = 5 \Leftrightarrow (x-4)^2 = 1 \Leftrightarrow x-4 = \pm 1 \Leftrightarrow x=3$ ou $x=5$.
    Les deux sont dans $[2;7]$. \textbf{2 solutions}.
\end{enumerate}
\end{corrige}

\begin{attention}[Dernier conseil de "Prof" pour le jour J]
\begin{itemize}
    \item \textbf{Gère ton temps :} L'étude de fonction vaut souvent 6 à 8 points sur 20\%. Ne t'éternise pas sur les calculs alambiqués : si le discriminant est moche, c'est que tu as fait une erreur avant (forme canonique, simplification). Les nombres sont "propres" au Probatoire.
    \item \textbf{Soigne le tableau :} C'est la "colonne vertébrale" de ta note. Un tableau propre, aligné, avec des flèches correctes, vaut souvent 2-3 points tout seul.
    \item \textbf{Trace grand :} Une courbe sur 10cm $\times$ 10cm minimum. Petits carreaux = imprécision = points perdus en lecture graphique.
    \item \textbf{Relis l'énoncé :} "Sur l'intervalle $I$", "Donner une valeur approchée au dixième", "Encadrer". Réponds \textbf{exactement} à ce qui est demandé.
\end{itemize}
Tu as maintenant tous les outils. La rigueur paie. Bon courage pour le Probatoire !
\end{attention}

\coursec{Modélisation et problèmes concrets (Applications camerounaises)}

Dans la section précédente, nous avons mis en place une stratégie globale d'étude de fonction : ensemble de définition, dérivée, tableau de variations, courbe, résolution graphique. À quoi tout cela sert-il concrètement ? C'est la question que pose tout élève, et elle est légitime. La réponse tient en un mot : \cle{modéliser}. Derrière chaque décision économique (quel prix fixer ?), agricole (quelle surface clôturer ?), industrielle (quelle dose de produit ajouter ?), se cache une fonction à étudier. Le Probatoire le sait bien : les problèmes d'optimisation et de lecture graphique y sont des classiques. Cette dernière section te montre comment passer d'une situation camerounaise concrète à une fonction polynôme du second degré ou homographique, puis comment exploiter ton étude pour répondre à la question posée.

\courssub{La démarche de modélisation : de la situation réelle à la fonction}

Imaginons une maraîchère de Bafoussam qui veut clôturer un enclos rectangulaire avec un stock limité de grillage. Elle se demande : \og quelles dimensions donnent la plus grande surface ? \fg{} C'est une question de la vraie vie, pas une question de maths\ldots{} pour l'instant. Le travail du mathématicien consiste à la \emph{traduire} : choisir une inconnue $x$, exprimer la grandeur à optimiser en fonction de $x$, puis étudier cette fonction sur l'intervalle imposé par la réalité.

\begin{methode}[Modéliser et résoudre un problème concret]
\begin{enumerate}
\item \textbf{Identifier les variables.} Distinguer la variable de \cle{décision} $x$ (ce sur quoi on peut agir : une longueur, un prix, une quantité) et la grandeur \cle{objectif} $y$ (ce qu'on veut maximiser ou minimiser : aire, bénéfice, coût, concentration).
\item \textbf{Écrire la relation} $y = f(x)$ à l'aide des formules géométriques (périmètres, aires), physiques ou économiques (prix $\times$ quantité, coût fixe $+$ coût variable, concentration $=$ masse/volume).
\item \textbf{Déterminer l'intervalle réaliste} $I$ en traduisant les contraintes physiques : longueurs positives, budget limité, capacité maximale. C'est l'ensemble de définition \emph{du problème}.
\item \textbf{Étudier $f$ sur $I$} : dérivée, signe, tableau de variations, extremum, éventuellement résolution de $f(x) = m$.
\item \textbf{Conclure en langage courant}, avec une phrase complète et les unités : \og Le prix optimal est $450$ FCFA pour un bénéfice maximal de $125\,000$ FCFA. \fg{}
\end{enumerate}
\end{methode}

\begin{attention}[Deux pièges mortels en modélisation]
\begin{itemize}
\item \textbf{Oublier l'intervalle réaliste.} Le sommet de la parabole peut se trouver \emph{en dehors} de l'intervalle imposé par le problème. L'extremum cherché est alors atteint à une borne de $I$. Toujours vérifier que la solution mathématique a un sens physique.
\item \textbf{Mélanger les unités.} Une longueur en mètres ne s'additionne pas à une longueur en centimètres ; une concentration s'exprime en g/L, pas en g. Avant tout calcul, harmonise les unités, et vérifie à la fin que le résultat porte la bonne unité (m$^2$ pour une aire, FCFA pour un bénéfice).
\end{itemize}
\end{attention}

\courssub{Problèmes d'optimisation : maximiser ou minimiser avec un trinôme}

\textbf{Intuition.} Beaucoup de grandeurs économiques ou géométriques suivent une \og arche \fg{} : elles augmentent, atteignent un sommet, puis diminuent. Le bénéfice en fonction du prix en est l'exemple type : prix trop bas, on ne gagne rien ; prix trop haut, personne n'achète ; entre les deux, il existe un prix optimal. Mathématiquement, cette arche est une parabole $f(x) = ax^2 + bx + c$ avec $a < 0$, et l'optimum est son sommet.

\begin{exemplebox}[L'enclos de l'éleveur (maximisation d'aire)]
Un éleveur de l'Ouest dispose de $\SI{100}{m}$ de grillage pour construire un enclos rectangulaire \textbf{adossé à un mur} : seuls trois côtés sont à griller. Soit $x$ la longueur (en m) des deux côtés perpendiculaires au mur. Le côté parallèle au mur mesure alors $L = 100 - 2x$. L'aire est :
\[ A(x) = x(100 - 2x) = -2x^2 + 100x. \]
\textbf{Intervalle réaliste :} il faut $x \geq 0$ et $L = 100 - 2x \geq 0$, donc $x \in [0\,;\,50]$.
\textbf{Étude :} $A'(x) = -4x + 100$, qui s'annule en $x = 25$, positif avant, négatif après. Le tableau de variations montre un maximum en $x = 25$ :
\[ A(25) = -2 \times 625 + 100 \times 25 = -1250 + 2500 = 1250. \]
\textbf{Conclusion :} l'aire maximale est de $\SI{1250}{m^2}$, obtenue pour des côtés perpendiculaires de $\SI{25}{m}$ et un côté parallèle au mur de $\SI{50}{m}$.
\end{exemplebox}

\begin{popfigure}
\begin{center}
\begin{tikzpicture}[scale=0.55]
% mur
\draw[line width=2.5pt,popdark] (0,8.6) -- (10.8,8.6);
\foreach \i in {0,0.6,...,10.5}{\draw[popdark] (\i,8.6) -- (\i+0.35,9.1);}
\node[popdark] at (5.4,9.7) {\textbf{Mur}};
% enclos
\draw[popblue,line width=1.5pt,fill=popblueL] (0.5,0.5) rectangle (10,8.6);
\draw[popgreen,line width=2pt] (0.5,0.5) -- (0.5,8.6);
\draw[popgreen,line width=2pt] (10,0.5) -- (10,8.6);
\draw[popgreen,line width=2pt] (0.5,0.5) -- (10,0.5);
% cotes
\draw[<->,poporange,line width=1pt] (0.5,0) -- (10,0) node[midway,below] {$L = 100 - 2x$};
\draw[<->,poporange,line width=1pt] (10.5,0.5) -- (10.5,8.6) node[midway,right] {$x$};
\draw[<->,poporange,line width=1pt] (0,0.5) -- (0,8.6) node[midway,left] {$x$};
\node[popblue] at (5.25,4.5) {$A(x) = x(100-2x)$};
\end{tikzpicture}
\end{center}
\legende{L'enclos adossé au mur : trois côtés grillagés, $2x + L = 100$ m.}
\end{popfigure}

\begin{savaistu}[La forme canonique, l'optimisation sans dérivée]
Écris l'aire de l'enclos sous forme canonique :
\[ A(x) = -2x^2 + 100x = -2(x - 25)^2 + 1250. \]
Comme $-2(x-25)^2 \leq 0$ pour tout $x$, on a $A(x) \leq 1250$, avec égalité exactement quand $x = 25$. Le maximum et la valeur optimale apparaissent \emph{directement}, sans aucune dérivée ! C'est toute la puissance de l'algèbre pour l'optimisation : la forme canonique $a(x-\alpha)^2 + \beta$ exhibe l'extremum $\beta$ atteint en $\alpha$. Au Probatoire, cette lecture immédiate est un réflexe qui fait gagner un temps précieux.
\end{savaistu}

\textbf{Un deuxième grand classique : le rectangle inscrit dans un triangle.} On inscrit un rectangle dans un triangle rectangle de base $\SI{10}{cm}$ et de hauteur $\SI{8}{cm}$. Si $x$ est la largeur du rectangle, le théorème de Thalès donne sa hauteur : $h = 8 - 0{,}8x$. D'où l'aire :
\[ A(x) = x(8 - 0{,}8x) = -0{,}8x^2 + 8x, \quad x \in [0\,;\,10]. \]
Le maximum est atteint en $x = \dfrac{-8}{2 \times (-0{,}8)} = 5$, et vaut $A(5) = 5 \times 4 = \SI{20}{cm^2}$.

\begin{popfigure}
\begin{center}
\begin{tikzpicture}[scale=0.5]
% triangle rectangle
\draw[popdark,line width=1.5pt] (0,0) -- (10,0) -- (0,8) -- cycle;
% rectangle inscrit (x=5, h=4)
\draw[poporange,line width=1.5pt,fill=poporangeL] (0,0) rectangle (5,4);
\draw[<->,popblue] (0,-0.7) -- (5,-0.7) node[midway,below] {$x$};
\draw[<->,popblue] (5.5,0) -- (5.5,4) node[midway,right] {$h = 8 - 0{,}8x$};
\draw[<->,popmuted] (0,-1.6) -- (10,-1.6) node[midway,below] {base $= 10$};
\draw[<->,popmuted] (-0.7,0) -- (-0.7,8) node[midway,left] {$8$};
\node[poporange] at (2.5,2) {$A(x)$};
\end{tikzpicture}
\end{center}
\legende{Rectangle inscrit dans un triangle rectangle : par Thalès, $\frac{8-h}{8} = \frac{x}{10}$, donc $h = 8 - 0{,}8x$.}
\end{popfigure}

\begin{exoresolu}[Bénéfice maximal d'un vendeur de téléphones à Douala]
\textbf{Énoncé.} Un revendeur de smartphones au marché Mboppi estime que s'il fixe le prix de vente à $x$ centaines de FCFA (avec $x \in [200\,;\,800]$), son bénéfice hebdomadaire en FCFA est :
\[ B(x) = -2x^2 + 1800x - 200\,000. \]
\begin{enumerate}
\item Étudier les variations de $B$ sur $[200\,;\,800]$.
\item En déduire le prix de vente optimal et le bénéfice maximal.
\item Pour quels prix le bénéfice dépasse-t-il $\SI{150000}{FCFA}$ ? (Résoudre $B(x) \geq 150\,000$.)
\end{enumerate}
\tcblower
\textbf{Solution.}
\textbf{1.} $B$ est un trinôme dérivable sur $\mathbb{R}$ : $B'(x) = -4x + 1800$. On a $B'(x) = 0 \Leftrightarrow x = 450$ ; $B'(x) > 0$ pour $x < 450$ et $B'(x) < 0$ pour $x > 450$. D'où le tableau :
\begin{center}
\begin{tabular}{|c|ccccc|}\hline
$x$ & $200$ & & $450$ & & $800$ \\ \hline
$B'(x)$ & & $+$ & $0$ & $-$ & \\ \hline
$B$ & $80\,000$ & $\nearrow$ & $205\,000$ & $\searrow$ & $-40\,000$ \\ \hline
\end{tabular}
\end{center}
avec $B(200) = -80\,000 + 360\,000 - 200\,000 = 80\,000$, $B(450) = -405\,000 + 810\,000 - 200\,000 = 205\,000$ et $B(800) = -1\,280\,000 + 1\,440\,000 - 200\,000 = -40\,000$.

\textbf{2.} Le bénéfice maximal est $\SI{205000}{FCFA}$, atteint pour $x = 450$ centaines de FCFA, c'est-à-dire un prix de vente de $\SI{45000}{FCFA}$ par téléphone.

\textbf{3.} $B(x) \geq 150\,000 \Leftrightarrow -2x^2 + 1800x - 350\,000 \geq 0$. Le discriminant vaut $\Delta = 1800^2 - 8 \times 350\,000 = 3\,240\,000 - 2\,800\,000 = 440\,000$, et $\sqrt{440\,000} \approx 663{,}3$. Les racines sont $x_1 \approx \dfrac{1800 - 663{,}3}{4} \approx 284$ et $x_2 \approx \dfrac{1800 + 663{,}3}{4} \approx 616$. Le trinôme est positif entre ses racines : le bénéfice dépasse $\SI{150000}{FCFA}$ pour un prix compris environ entre $\SI{28400}{FCFA}$ et $\SI{61600}{FCFA}$.
\end{exoresolu}

\begin{popfigure}
\begin{center}
\begin{tikzpicture}
\begin{axis}[
  width=0.9\linewidth, height=6.5cm,
  axis lines=middle,
  xlabel={$x$ (centaines de FCFA)}, ylabel={$B(x)$ (FCFA)},
  xmin=150, xmax=850, ymin=-80000, ymax=240000,
  xtick={200,284,450,616,800},
  x tick label style={font=\small},
  y tick label style={font=\small},
  grid=both, grid style={popline!40},
]
\addplot[popcurve,domain=200:800,samples=150]{-2*x^2 + 1800*x - 200000};
\addplot[poporange,dashed,domain=150:850]{150000};
\addplot[popgreen,dashed,domain=150:850]{0};
\addplot[only marks,mark=*,poporange,mark size=2pt] coordinates {(450,205000)};
\node[poporange,anchor=south,font=\small] at (axis cs:450,205000) {Sommet $(450\,;\,205\,000)$};
\node[poporange,anchor=south west,font=\small] at (axis cs:620,155000) {$B(x)=150\,000$};
\end{axis}
\end{tikzpicture}
\end{center}
\legende{Courbe du bénéfice : lecture graphique de l'optimum et du seuil $B(x) = 150\,000$.}
\end{popfigure}

\courssub{Concentration et dilution : la fonction homographique en action}

\textbf{Intuition.} En chimie comme en cuisine, la concentration d'un mélange est un rapport : $C = \dfrac{\text{quantité de produit}}{\text{volume total}}$. Si on ajoute du produit pur, le numérateur augmente ; si on ajoute du solvant, le dénominateur augmente. Quand numérateur \emph{et} dénominateur dépendent de $x$ de façon affine, on obtient une fonction \cle{homographique} $C(x) = \dfrac{ax+b}{cx+d}$.

\begin{exemplebox}[Deux situations à bien distinguer]
Un réservoir contient $\SI{200}{L}$ d'eau salée à $\SI{10}{g/L}$, soit $\SI{2000}{g}$ de sel.
\begin{itemize}
\item \textbf{Cas 1 : on remplace $x$ litres de solution par de l'eau pure} (volume constant $\SI{200}{L}$). Le sel retiré est $10x$ grammes, donc $C(x) = \dfrac{2000 - 10x}{200} = 10 - 0{,}05x$ : fonction \textbf{affine} décroissante sur $[0\,;\,200]$.
\item \textbf{Cas 2 : on ajoute $x$ kg de sel pur, le volume augmentant de $x$ litres} (dissolution). Alors :
\[ C(x) = \frac{2000 + 1000x}{200 + x} \ \text{g/L}, \quad x \geq 0. \]
C'est une fonction \textbf{homographique}. Sa dérivée est :
\[ C'(x) = \frac{1000(200+x) - (2000+1000x)}{(200+x)^2} = \frac{198\,000}{(200+x)^2} > 0. \]
$C$ est strictement croissante sur $[0\,;\,+\infty[$. En réécrivant $C(x) = 1000 - \dfrac{198\,000}{x+200}$, on voit que $C(x)$ se rapproche de $\SI{1000}{g/L}$ sans jamais l'atteindre : la droite $y = 1000$ est \cle{asymptote horizontale} (concentration du sel pur, limite physique du procédé).
\end{itemize}
\end{exemplebox}

\begin{attention}[Volume constant ou volume variable ?]
C'est LA question à se poser avant d'écrire le modèle. Si le volume reste constant (on vide autant qu'on ajoute), la concentration est une fonction \emph{affine}. Si le volume varie avec $x$, la concentration devient \emph{homographique}, et il faut alors écrire soigneusement le nouveau numérateur (masse totale) et le nouveau dénominateur (volume total). Une erreur ici fausse toute l'étude.
\end{attention}

\courssub{Tarification : fonctions affines et interprétation du paramètre $m$}

\textbf{Facturation \og abonnement $+$ consommation \fg{}.} Un forfait téléphonique à $\SI{1000}{FCFA}$ d'abonnement plus $\SI{25}{FCFA}$ par minute se modélise par la fonction affine $f(x) = 1000 + 25x$, où $x$ est le nombre de minutes. Le coût fixe est l'ordonnée à l'origine, le tarif unitaire est le coefficient directeur. Comparer deux opérateurs revient à comparer deux droites et à résoudre $f_1(x) = f_2(x)$ : au-delà du point d'intersection, l'une devient plus avantageuse que l'autre.

\textbf{Le paramètre $m$ : trois lectures concrètes.} Résoudre $f(x) = m$ ou discuter le nombre de solutions selon $m$, c'est répondre à des questions très concrètes :
\begin{itemize}
\item \textbf{Seuil de rentabilité} ($m = 0$) : résoudre $B(x) = 0$ donne les prix pour lesquels le vendeur ne gagne ni ne perd. Entre les deux racines, l'activité est rentable.
\item \textbf{Objectif de production} ($m =$ valeur cible) : \og pour quels prix le bénéfice atteint-il $\SI{150000}{FCFA}$ ? \fg{} — c'est l'équation $B(x) = 150\,000$ résolue dans l'exercice précédent.
\item \textbf{Contrainte réglementaire} ($m =$ norme) : \og la concentration ne doit pas dépasser $\SI{50}{g/L}$ \fg{} se traduit par l'inéquation $C(x) \leq 50$.
\end{itemize}
Graphiquement, tout cela se lit par l'intersection de la courbe $\mathcal{C}_f$ avec la droite horizontale $y = m$ : le nombre de points d'intersection donne le nombre de solutions, et leur position relative résout les inéquations.

\courssub{Rédiger la conclusion : la phrase qui rapporte des points}

Une étude de fonction dans un problème n'est jamais finie tant que tu n'as pas répondu à la question \emph{en français}, avec les unités. Le correcteur du Probatoire cherche cette phrase.

\begin{methode}[Rédiger une conclusion de problème]
\begin{enumerate}
\item Reprendre la question posée dans l'énoncé (pas une autre !).
\item Donner la valeur de $x$ optimale \textbf{avec son unité} et sa signification concrète (prix, longueur, quantité).
\item Donner la valeur de l'extremum \textbf{avec son unité} (FCFA, m$^2$, g/L).
\item Vérifier la cohérence : la solution appartient-elle à l'intervalle réaliste ? Le résultat est-il plausible (une aire négative ou un prix de $3$ milliards de FCFA doivent te faire douter) ?
\end{enumerate}
Modèle : \og L'aire maximale de l'enclos est $\SI{1250}{m^2}$, obtenue avec des côtés de $\SI{25}{m}$ perpendiculaires au mur et $\SI{50}{m}$ parallèlement au mur. \fg{}
\end{methode}

\begin{exercice}[La boîte pliée]
D'une feuille cartonnée carrée de $\SI{30}{cm}$ de côté, on découpe quatre carrés de côté $x$ aux coins, puis on plie pour former une boîte sans couvercle.
\begin{enumerate}
\item Montrer que le volume est $V(x) = x(30 - 2x)^2$ et préciser l'intervalle réaliste.
\item On admet que $V$ admet un maximum sur cet intervalle en $x = 5$. Calculer ce volume maximal.
\end{enumerate}
\end{exercice}

\begin{corrige}[Exercice 1]
\textbf{1.} Après découpe, la base de la boîte est un carré de côté $30 - 2x$ et la hauteur est $x$, d'où $V(x) = x(30-2x)^2$. Il faut $x \geq 0$ et $30 - 2x \geq 0$, donc $x \in [0\,;\,15]$.
\textbf{2.} $V(5) = 5 \times (30 - 10)^2 = 5 \times 400 = 2\,000$ cm$^3$. Le volume maximal de la boîte est $\SI{2000}{cm^3}$ (soit $\SI{2}{L}$), obtenu en découpant des carrés de $\SI{5}{cm}$ de côté.
\end{corrige}

\begin{savaistu}[Conseil Probatoire]
Les sujets de Première A tombent très souvent sur le schéma : \og Étudier $f$ sur $[a\,;\,b]$, construire sa courbe, puis résoudre graphiquement $f(x) = m$ ou discuter selon $m$. \fg{} Deux réflexes te font gagner un quart d'heure : (1) passer par la \textbf{forme canonique} $a(x-\alpha)^2 + \beta$ pour tracer la parabole par translation de celle de $y = ax^2$ (vecteur $\alpha\vec{\imath} + \beta\vec{\jmath}$) ; (2) écrire l'homographique sous la forme $\alpha + \dfrac{k}{x - \beta}$ pour obtenir son hyperbole par translation de celle de $y = \dfrac{k}{x}$, avec ses deux asymptotes $x = \beta$ et $y = \alpha$ immédiates. La modélisation n'est alors plus qu'une affaire de lecture attentive de l'énoncé.
\end{savaistu}

Tu as désormais tous les outils : transformer une situation réelle en fonction, l'étudier rigoureusement sur un intervalle borné, exploiter les transformations de courbes, et conclure avec précision. C'est exactement la compétence que le programme attend de toi — et c'est aussi, disons-le, une manière de voir le monde : derrière chaque décision optimale, il y a une fonction qui attend d'être étudiée.

\newpage

\coursec{Activité d'intégration}

\courssub{Objectifs de l'activité}
Cette activité d'intégration a pour but de mobiliser l'ensemble des compétences travaillées dans la leçon \og Fonctions associées et étude de fonctions \fg dans un contexte unique, réaliste et motivant. Elle vise plus précisément à :
\begin{itemize}
    \item Réaliser une étude complète (ensemble de définition, forme canonique, dérivée, tableau de variations, extremums, représentation graphique) d'une fonction polynôme du second degré sur un intervalle borné.
    \item Résoudre graphiquement une équation $f(x)=m$ et discuter le nombre de solutions selon le paramètre $m$.
    \item Construire la courbe d'une fonction obtenue par translation verticale ($g(x)=f(x)+b$) et en déduire les nouveaux extremums.
    \item Construire la courbe de la fonction valeur absolue $h(x)=|f(x)|$ sur un intervalle restreint et en donner l'interprétation économique.
    \item Synthétiser les résultats pour formuler un conseil argumenté (savoir-être : communication, prise de décision).
\end{itemize}

\courssub{Situation}
Maman Ngo tient un stand de jus baptisé \og Vitalité\fg au marché Mokolo à Yaoundé. Après plusieurs semaines d'observation, elle a modélisé son bénéfice journalier net (en milliers de FCFA) en fonction du prix de vente unitaire d'un verre (en centaines de FCFA) par la fonction $f$ définie sur l'intervalle $[2\,;\,8]$ par :
\[
f(x) = -x^2 + 10x - 16.
\]
Concrètement :
\begin{itemize}
    \item $x=2$ correspond à un prix de vente de \SI{200}{FCFA} le verre ; $x=8$ correspond à \SI{800}{FCFA}.
    \item $f(x)=5$ signifie un bénéfice journalier de \SI{5000}{FCFA}.
    \item Le domaine $[2\,;\,8]$ reflète la réalité du marché : en dessous de \SI{200}{FCFA}, elle perd de l'argent ; au-dessus de \SI{800}{FCFA}, les clients ne viennent plus.
\end{itemize}
La mairie de Yaoundé IV vient d'instaurer une taxe journalière fixe de \SI{2000}{FCFA} pour tout commerce ambulant. Par ailleurs, son fournisseur principal de fruits impose désormais un prix plancher de \SI{500}{FCFA} le verre (soit $x \ge 5$) pour garantir sa marge. Maman Ngo souhaite optimiser son activité.

\courssub{Consignes}
\begin{enumerate}
    \item \textbf{Étude de la fonction initiale $f$.} Déterminer l'ensemble de définition, la forme canonique, la fonction dérivée $f'$, dresser le tableau de variations de $f$ sur $[2\,;\,8]$, préciser les extremums et tracer la courbe $\mathcal{C}_f$ dans un repère orthonormé.
    \item \textbf{Résolution graphique paramétrée.} Maman Ngo vise un bénéfice de \SI{5000}{FCFA} par jour ($m=5$). Résoudre graphiquement l'équation $f(x)=5$ sur $[2\,;\,8]$. Combien de prix de vente permettent d'atteindre cet objectif ? Justifier par le tableau de variations.
    \item \textbf{Impact de la taxe (translation).} La taxe de \SI{2000}{FCFA} modélise une translation verticale. La nouvelle fonction bénéfice est $g(x) = f(x) - 2$ sur $[2\,;\,8]$.
    \begin{itemize}
        \item Écrire $g$ sous forme canonique.
        \item Construire la courbe $\mathcal{C}_g$ par translation de $\mathcal{C}_f$ (préciser le vecteur).
        \item Quel est le nouveau bénéfice maximal journalier ? À quel prix l'atteint-elle ?
    \end{itemize}
    \item \textbf{Contrainte du fournisseur (valeur absolue).} Avec le prix plancher $x \ge 5$, on restreint l'étude à $[5\,;\,8]$. On définit $h(x) = |f(x)|$ sur $[5\,;\,8]$.
    \begin{itemize}
        \item Rappeler la définition de la valeur absolue et expliquer l'interprétation économique : \og le manque à gagner est comptabilisé comme une dette \fg.
        \item Construire la courbe $\mathcal{C}_h$ à partir de $\mathcal{C}_f$ sur $[5\,;\,8]$.
        \item $h$ admet-elle un minimum sur $[5\,;\,8]$ ? Lequel ?
    \end{itemize}
    \item \textbf{Synthèse orale (Conseil à Maman Ngo).} En vous appuyant sur les résultats précédents, rédigez un court paragraphe de conseil : quel prix fixer avant/après taxe ? Quel risque si elle descend sous \SI{500}{FCFA} ? Que représente la courbe de $h$ pour sa trésorerie ?
\end{enumerate}

\courssub{Ressources à mobiliser}
\begin{itemize}
    \item \textbf{Forme canonique :} $f(x) = a(x-\alpha)^2 + \beta$ avec $\alpha = -\frac{b}{2a}$, $\beta = f(\alpha)$. Sommet de la parabole : $S(\alpha\,;\,\beta)$.
    \item \textbf{Dérivée :} $f'(x) = 2ax + b$. Signe de $f'$ $\Leftrightarrow$ sens de variation de $f$.
    \item \textbf{Tableau de variations :} Une colonne par valeur remarquable (bornes, racine de $f'$) et une colonne \og entre \fg pour le signe de $f'$ et la flèche de variation.
    \item \textbf{Translation :} $g(x)=f(x)+b$ $\Leftrightarrow$ translation de vecteur $\vec{v}(0\,;\,b)$ (verticale). Si $b<0$, translation vers le bas.
    \item \textbf{Valeur absolue :} $|f(x)| = f(x)$ si $f(x) \ge 0$ ; $|f(x)| = -f(x)$ si $f(x) < 0$. Graphiquement : la partie de $\mathcal{C}_f$ au-dessus de l'axe des abscisses reste inchangée ; la partie en dessous est refléchie par symétrie axiale par rapport à l'axe des abscisses.
    \item \textbf{Résolution graphique $f(x)=m$ :} Abscisses des points d'intersection de $\mathcal{C}_f$ avec la droite $y=m$.
\end{itemize}

\courssub{Démarche et corrigé commenté}

\begin{exoresolu}{Question 1 : Étude complète de $f$ sur $[2\,;\,8]$}
\emph{Ensemble de définition :} $D_f = [2\,;\,8]$ (donné par le contexte).

\emph{Forme canonique :} $f(x) = -x^2 + 10x - 16$. On factorise par $-1$ :
\[
f(x) = -(x^2 - 10x) - 16 = -\big[(x-5)^2 - 25\big] - 16 = -(x-5)^2 + 25 - 16 = -(x-5)^2 + 9.
\]
Le sommet de la parabole est $S(5\,;\,9)$. Comme $a=-1<0$, la parabole est tournée vers le bas : $S$ est un maximum global sur $\mathbb{R}$, donc aussi sur $[2\,;\,8]$.

\emph{Dérivée et signe :} $f'(x) = -2x + 10 = -2(x-5)$.
Sur $[2\,;\,8]$ :
\begin{itemize}
    \item $f'(x) > 0 \Leftrightarrow x < 5$ : $f$ est strictement croissante sur $[2\,;\,5]$.
    \item $f'(x) = 0 \Leftrightarrow x = 5$ : extremum (maximum).
    \item $f'(x) < 0 \Leftrightarrow x > 5$ : $f$ est strictement décroissante sur $[5\,;\,8]$.
\end{itemize}

\emph{Tableau de variations :}
\begin{center}
\begin{tabular}{|c|ccccc|}\hline
$x$ & $2$ & & $5$ & & $8$ \\ \hline
$f'(x)$ & & $+$ & $0$ & $-$ & \\ \hline
$f$ & $0$ & $\nearrow$ & $9$ & $\searrow$ & $0$ \\ \hline
\end{tabular}
\end{center}

\emph{Extremums sur $[2\,;\,8]$ :}
\begin{itemize}
    \item Maximum absolu (global) : $9$ (soit \SI{9000}{FCFA}) atteint pour $x=5$ (prix \SI{500}{FCFA}).
    \item Minimum absolu : $0$ (bénéfice nul) atteint aux bornes $x=2$ et $x=8$ (prix \SI{200}{FCFA} et \SI{800}{FCFA}).
\end{itemize}
\end{exoresolu}

\begin{popfigure}
\begin{tikzpicture}[scale=0.7]
\begin{axis}[
    axis lines=middle,
    xlabel=$x$ (prix en 100 FCFA), ylabel=$y$ (bénéfice en 1000 FCFA),
    xmin=1, xmax=9, ymin=-1, ymax=10,
    xtick={2,3,4,5,6,7,8}, ytick={0,2,4,6,8,9},
    grid=major, grid style={popline},
    width=\linewidth, height=8cm,
    clip=false,
    axis line style={-Stealth, thick},
    tick label style={font=\small},
    label style={font=\small},
]
\addplot[popcurve, domain=2:8, samples=200, thick] {-x^2 + 10*x - 16};
\addplot[only marks, mark=*, popblue, thick] coordinates {(2,0) (5,9) (8,0)};
\node[above] at (axis cs:5,9) {$S(5;9)$};
\node[below left] at (axis cs:2,0) {$A(2;0)$};
\node[below right] at (axis cs:8,0) {$B(8;0)$};
\draw[dashed, popmuted] (axis cs:5,0) -- (axis cs:5,9);
\node[below] at (axis cs:5,0) {$5$};
\end{axis}
\end{tikzpicture}
\legende{Courbe $\mathcal{C}_f$ de la fonction bénéfice $f(x)=-x^2+10x-16$ sur $[2;8]$.}
\end{popfigure}

\begin{exoresolu}{Question 2 : Résolution graphique de $f(x)=5$ (bénéfice de 5000 FCFA)}
On trace la droite horizontale $y=5$ (en pointillés rouges sur la figure ci-dessus). Les intersections avec $\mathcal{C}_f$ donnent les solutions.
Algébriquement : $-x^2+10x-16 = 5 \Leftrightarrow -x^2+10x-21=0 \Leftrightarrow x^2-10x+21=0$.
$\Delta = 100 - 84 = 16$. $\sqrt{\Delta}=4$.
$x_1 = \frac{10-4}{2}=3$ ; $x_2 = \frac{10+4}{2}=7$.
Les deux solutions $x=3$ (\SI{300}{FCFA}) et $x=7$ (\SI{700}{FCFA}) appartiennent à $[2\,;\,8]$.
\emph{Conclusion :} Il y a \textbf{deux prix de vente} permettant d'atteindre \SI{5000}{FCFA} de bénéfice : \SI{300}{FCFA} ou \SI{700}{FCFA} le verre.
Sur le tableau de variations, la valeur $5$ est comprise entre $0$ et $9$. Comme $f$ est continue et strictement monotone par morceaux, l'équation $f(x)=5$ admet exactement une solution sur $[2\,;\,5[$ (croissante) et une sur $]5\,;\,8]$ (décroissante).
\end{exoresolu}

\begin{exoresolu}{Question 3 : Impact de la taxe : fonction $g(x)=f(x)-2$}
\emph{Forme canonique :}
$g(x) = -(x-5)^2 + 9 - 2 = -(x-5)^2 + 7$.
Le sommet de $\mathcal{C}_g$ est $S_g(5\,;\,7)$.

\emph{Translation :} $g(x) = f(x) + (-2)$. C'est la translation de vecteur $\vec{v}(0\,;\,-2)$.
Graphiquement, $\mathcal{C}_g$ s'obtient en décalant $\mathcal{C}_f$ de \textbf{2 unités vers le bas} (soit \SI{2000}{FCFA} de bénéfice en moins pour chaque prix).
\end{exoresolu}

\begin{popfigure}
\begin{tikzpicture}[scale=0.7]
\begin{axis}[
    axis lines=middle, xlabel=$x$, ylabel=$y$,
    xmin=1, xmax=9, ymin=-3, ymax=10,
    xtick={2,3,4,5,6,7,8}, ytick={-2,0,2,4,6,7,9},
    grid=major, grid style={popline}, width=\linewidth, height=8cm,
    axis line style={-Stealth, thick}, tick label style={font=\small}, label style={font=\small},
]
\addplot[popcurve, domain=2:8, samples=200, thick, popblue] {-x^2 + 10*x - 16};
\addplot[popcurve, domain=2:8, samples=200, thick, poporange, dashed] {-x^2 + 10*x - 18};
\addplot[only marks, mark=*, popblue] coordinates {(2,0) (5,9) (8,0)};
\addplot[only marks, mark=*, poporange] coordinates {(2,-2) (5,7) (8,-2)};
\node[above, popblue] at (axis cs:5,9) {$S_f(5;9)$};
\node[above, poporange] at (axis cs:5,7) {$S_g(5;7)$};
\draw[<->, popmuted, thick] (axis cs:5,9) -- (axis cs:5,7) node[midway, right] {$\vec{v}(0;-2)$};
\legend{$\mathcal{C}_f$, $\mathcal{C}_g$}
\end{axis}
\end{tikzpicture}
\legende{Translation verticale : $\mathcal{C}_g$ (orange, pointillés) image de $\mathcal{C}_f$ (bleu) par $\vec{v}(0;-2)$.}
\end{popfigure}

\begin{exoresolu}{Question 3 (suite) : Nouvel optimum}
\emph{Nouvel optimum :} Le maximum de $g$ sur $[2\,;\,8]$ est $g(5)=7$ (soit \SI{7000}{FCFA}), atteint pour le même prix optimal $x=5$ (\SI{500}{FCFA}).
\emph{Remarque :} Les zéros de $g$ ne sont plus aux bornes. $g(2)=g(8)=-2$ (perte de \SI{2000}{FCFA}). L'équation $g(x)=0$ a deux solutions dans $[2\,;\,8]$ (seuil de rentabilité déplacé vers $x=5\pm\sqrt{7} \approx 2,35$ et $7,65$).
\end{exoresolu}

\begin{exoresolu}{Question 4 : Contrainte du fournisseur : fonction $h(x)=|f(x)|$ sur $[5\,;\,8]$}
\emph{Rappel et interprétation :} Pour tout réel $y$, $|y| = y$ si $y\ge 0$ et $|y| = -y$ si $y<0$.
Ici, $h(x)=|f(x)|$. Sur l'intervalle $[5\,;\,8]$, d'après le tableau de variations, $f(x)$ décroît de $9$ à $0$. Donc $\forall x \in [5\,;\,8],\ f(x) \ge 0$.
Par conséquent, \textbf{$h(x) = f(x)$ sur tout l'intervalle $[5\,;\,8]$}.
L'interprétation \og manque à gagner comme dette \fg serait pertinente si $f(x)$ devenait négatif. Sur $[5\,;\,8]$, il n'y a pas de manque à gagner (bénéfice positif ou nul), la valeur absolue ne change rien.

\emph{Construction de $\mathcal{C}_h$ :} $\mathcal{C}_h$ coïncide exactement avec la branche droite de $\mathcal{C}_f$ (de $S(5\,;\,9)$ à $B(8\,;\,0)$).
\end{exoresolu}

\begin{popfigure}
\begin{tikzpicture}[scale=0.7]
\begin{axis}[
    axis lines=middle, xlabel=$x$, ylabel=$y$,
    xmin=4, xmax=9, ymin=-1, ymax=10,
    xtick={5,6,7,8}, ytick={0,2,4,6,8,9},
    grid=major, grid style={popline}, width=\linewidth, height=7cm,
    axis line style={-Stealth, thick}, tick label style={font=\small}, label style={font=\small},
]
\addplot[popcurve, domain=5:8, samples=200, very thick, popgreen] {-x^2 + 10*x - 16};
\addplot[only marks, mark=*, popgreen] coordinates {(5,9) (8,0)};
\node[above, popgreen] at (axis cs:5,9) {$S(5;9)$};
\node[below right, popgreen] at (axis cs:8,0) {$B(8;0)$};
\end{axis}
\end{tikzpicture}
\legende{Courbe $\mathcal{C}_h = \mathcal{C}_{|f|}$ sur $[5;8]$. Comme $f\ge 0$, $\mathcal{C}_h$ se confond avec $\mathcal{C}_f$.}
\end{popfigure}

\begin{exoresolu}{Question 4 (suite) : Minimum de $h$}
\emph{Minimum de $h$ sur $[5\,;\,8]$ :} $h$ est décroissante sur $[5\,;\,8]$. Son minimum est $h(8)=0$ (bénéfice nul au prix max). Son maximum est $h(5)=9$.
\end{exoresolu}

\begin{exoresolu}{Question 5 : Synthèse orale — Conseil à Maman Ngo}
\og Chère Maman Ngo, voici mon analyse pour votre stand \og Vitalité\fg :

\begin{itemize}
    \item \textbf{Prix optimal actuel :} Fixez votre verre à \textbf{\SI{500}{FCFA}} ($x=5$). C'est le prix qui maximise votre bénéfice journalier à \textbf{\SI{9000}{FCFA}}. Si vous vendez à \SI{300}{FCFA} ou \SI{700}{FCFA}, vous ne gagnez que \SI{5000}{FCFA}. En dessous de \SI{200}{FCFA} ou au-dessus de \SI{800}{FCFA}, vous perdez tout.
    \item \textbf{Impact de la taxe (\SI{2000}{FCFA}/jour) :} Cette taxe fixe abaisse votre bénéfice maximal à \textbf{\SI{7000}{FCFA}} (toujours à \SI{500}{FCFA} le verre). Votre seuil de rentabilité remonte : vous ne serez plus bénéficiaire si vous vendez trop bas (vers \SI{235}{FCFA}) ou trop haut (vers \SI{765}{FCFA}). \textbf{Ne changez pas votre prix optimal}, mais sachez que votre marge de manœuvre se réduit.
    \item \textbf{Contrainte fournisseur ($x \ge 5$) :} Cette contrainte vous \textbf{force à rester sur la branche décroissante} de votre courbe de bénéfice. Vous ne pouvez plus baisser le prix sous \SI{500}{FCFA} pour attirer plus de clients. Votre bénéfice ira de \SI{9000}{FCFA} (à \SI{500}{FCFA}) down to \SI{0}{FCFA} (à \SI{800}{FCFA}).
    \item \textbf{La courbe $h=|f|$ :} Sur l'intervalle imposé $[500; 800]$, votre bénéfice est toujours positif ou nul. La \og dette \fg (valeur absolue) ne s'active pas. Cela signifie que \textbf{tant que vous respectez le prix plancher, vous ne créez pas de trou de trésorerie}. En revanche, si vous aviez le droit de descendre sous \SI{500}{FCFA} (branche montante), le bénéfice resterait positif aussi (jusqu'à \SI{200}{FCFA}). Le risque de \og dette \fg (bénéfice négatif) n'apparaît qu'en dehors de $[2;8]$, zone que vous avez sagement exclue.
\end{itemize}
\textbf{Recommandation finale :} Vendez à \textbf{\SI{500}{FCFA}} le verre. C'est le point d'équilibre parfait entre volume et marge, respectant la contrainte du fournisseur et maximisant le profit net après taxe. Surveillez vos coûts variables (fruits, glace) car la taxe fixe ne pardonne pas les baisses de volume. \fg
\end{exoresolu}

\courssub{Bilan de l'activité}
Cette activité a permis de mettre en évidence la puissance de l'outil \og fonction \fg pour modéliser et piloter une activité économique réelle. Les points clés acquis/consolidés sont :
\begin{enumerate}
    \item \textbf{Maîtrise de l'étude complète sur intervalle borné :} Passage fluide entre forme développée, canonique, dérivée, tableau de variations et lecture graphique. L'identification correcte des bornes de l'intervalle comme candidats aux extremums (en plus du point critique) est cruciale.
    \item \textbf{Lien algébrique / graphique pour $f(x)=m$ :} La discussion du nombre de solutions selon $m$ (0, 1 ou 2 solutions) est visuelle sur le tableau de variations (valeur prise entre min et max $\to$ 2 solutions par monotonie ; valeur = extremum $\to$ 1 solution ; valeur hors image $\to$ 0 solution).
    \item \textbf{Translation verticale :} L'écriture $g(x)=f(x)+b$ se traduit immédiatement par un vecteur $\vec{v}(0\,;\,b)$. Les extremums se translatent de même vecteur, l'abscisse de l'optimum ne change pas, seule l'ordonnée varie. Ici, la taxe fixe ($b=-2$) baisse le bénéfice max sans changer le prix optimal.
    \item \textbf{Valeur absolue et restriction de domaine :} La construction de $|f|$ dépend \emph{entièrement} du signe de $f$ sur l'intervalle considéré. Sur $[5\,;\,8]$, $f\ge 0 \Rightarrow |f|=f$ : l'activité teste la vigilance de l'élève face à une transformation qui \og ne fait rien \fg sur l'intervalle donné, mais dont l'interprétation économique (pas de dette) est riche de sens.
    \item \textbf{Compétence transversale (Savoir-être) :} La synthèse finale impose de traduire le langage mathématique (extremum, translation, intersection) en langage de gestion (prix optimal, marge, seuil de rentabilité, trésorerie), compétence essentielle pour le Probatoire et la vie professionnelle.
\end{enumerate}
Cet exercice type \og Bac/Probatoire \fg couvre l'essentiel du programme : étude de fonction degré 2, transformations, résolution graphique, modélisation.

\newpage

\coursec{Méthodes et exercices résolus}

\coursec{Outils préliminaires : Dérivée, forme canonique et ensemble de définition}

\begin{definition}[Dérivée des fonctions usuelles]
Soit $I$ un intervalle de $\mathbb{R}$.
\begin{itemize}
\item Si $f(x)=ax^2+bx+c$ (polynôme degré 2), alors $f$ est dérivable sur $\mathbb{R}$ et $f'(x)=2ax+b$.
\item Si $f(x)=\dfrac{ax+b}{cx+d}$ (fonction homographique) avec $ad-bc\neq0$, alors $f$ est dérivable sur $D_f=\mathbb{R}\setminus\{-\frac{d}{c}\}$ et
\[f'(x)=\frac{ad-bc}{(cx+d)^2}.\]
\end{itemize}
\end{definition}

\begin{propriete}[Forme canonique d'un trinôme]
Tout trinôme $ax^2+bx+c$ ($a\neq0$) s'écrit de manière unique sous la forme canonique :
\[a\left(x+\frac{b}{2a}\right)^2+\frac{4ac-b^2}{4a}=a(x-\alpha)^2+\beta,\]
avec $\alpha=-\dfrac{b}{2a}$ (abscisse du sommet) et $\beta=f(\alpha)=-\dfrac{\Delta}{4a}$ (ordonnée du sommet, $\Delta=b^2-4ac$).
\end{propriete}

\begin{propriete}[Écriture réduite d'une fonction homographique]
Toute fonction homographique $f(x)=\dfrac{ax+b}{cx+d}$ ($c\neq0$, $ad-bc\neq0$) s'écrit de manière unique sous la forme :
\[f(x)=\alpha+\frac{k}{x-\beta},\quad\text{avec }\alpha=\frac{a}{c},\;\beta=-\frac{d}{c},\;k=\frac{bc-ad}{c^2}.\]
Cette forme met en évidence la translation de vecteur $\vec{u}\begin{pmatrix}\beta\\ \alpha\end{pmatrix}$ qui transforme la courbe de $x\mapsto\dfrac{k}{x}$ en celle de $f$.
\end{propriete}

\begin{methode}[Déterminer l'ensemble de définition d'une fonction homographique]
\begin{enumerate}
\item Identifier le dénominateur $cx+d$.
\item Résoudre $cx+d=0 \iff x=-\dfrac{d}{c}$ : c'est la \cle{valeur interdite}.
\item Écrire $D_f=\mathbb{R}\setminus\left\{-\dfrac{d}{c}\right\}=]-\infty\,;-\frac{d}{c}[\cup]-\frac{d}{c}\,;+\infty[$.
\item \textbf{Vérifier} que l'intervalle d'étude $I$ donné par l'énoncé ne contient pas cette valeur interdite.
\end{enumerate}
\end{methode}

\begin{savaistu}[Histoire : Al-Khwarizmi et le « complétion du carré »]
Au IX\textsuperscript{e} siècle, le mathématicien perse Al-Khwarizmi (dont le nom a donné « algorithme ») résolvait les équations du second degré par une méthode géométrique : « compléter le carré ». Il transformait $x^2+bx=c$ en ajoutant $(b/2)^2$ des deux côtés pour former un carré parfait. C'est l'ancêtre direct de notre forme canonique $a(x-\alpha)^2+\beta$ !
\end{savaistu}

\coursec{Étude complète d'une fonction polynôme du second degré sur un intervalle borné}

\begin{methode}[{Étudier les variations de $f(x)=ax^2+bx+c$ sur un intervalle borné $I=[u\,;v]$}]
\begin{enumerate}
\item \textbf{Ensemble d'étude :} Préciser $I$ donné. Justifier que $f$ est dérivable sur $I$ (polynôme $\Rightarrow$ dérivable sur $\mathbb{R}$).
\item \textbf{Dérivée :} Calculer $f'(x)=2ax+b$.
\item \textbf{Valeurs critiques :} Résoudre $f'(x)=0 \iff x_0=-\dfrac{b}{2a}$. Vérifier si $x_0\in I$.
\item \textbf{Signe de $f'$ :} $f'$ est affine, de coefficient directeur $2a$.
\begin{itemize}
\item Si $a>0$ : $f'(x)<0$ pour $x<x_0$, $f'(x)>0$ pour $x>x_0$.
\item Si $a<0$ : $f'(x)>0$ pour $x<x_0$, $f'(x)<0$ pour $x>x_0$.
\end{itemize}
\item \textbf{Conclusion rédactionnelle :} « $f$ est strictement décroissante sur $[u\,;x_0]$ et strictement croissante sur $[x_0\,;v]$ » (adapter selon $a$ et la position de $x_0$ par rapport à $I$).
\item \textbf{Tableau de variations :} Calculer $f(u)$, $f(x_0)$ (si $x_0\in I$), $f(v)$. Remplir le tableau à 3 lignes ($x$, $f'$, $f$) et 5 colonnes (si $x_0\in I$) ou 3 colonnes (si $x_0\notin I$).
\item \textbf{Extremum :} Si $x_0\in I$, $f$ admet un extremum en $x_0$ : \cle{minimum} si $a>0$, \cle{maximum} si $a<0$. Sinon, les extremums sont aux bornes.
\end{enumerate}
\end{methode}

\begin{exoresolu}[Variations d'un trinôme sur un intervalle borné]
On considère la fonction $f$ définie sur $I=[-1\,;4]$ par $f(x)=x^2-2x-3$.
\begin{enumerate}
\item Étudier le sens de variation de $f$ sur $I$.
\item Dresser le tableau de variations de $f$ sur $I$.
\item En déduire le minimum de $f$ sur $I$ et la valeur en laquelle il est atteint.
\end{enumerate}
\tcblower
\textbf{1. Sens de variation.}

$f$ est une fonction polynôme, donc dérivable sur $\mathbb{R}$, en particulier sur $I=[-1\,;4]$. Pour tout $x\in I$ :
\[f'(x)=2x-2.\]
Résolvons $f'(x)=0$ : $2x-2=0 \iff x=1$. Or $1\in[-1\,;4]$, cette valeur appartient à l'ensemble d'étude.

$f'$ est une fonction affine de coefficient directeur $2>0$, donc :
\begin{itemize}
\item pour $x\in[-1\,;1[$, $f'(x)<0$ ;
\item pour $x\in]1\,;4]$, $f'(x)>0$.
\end{itemize}
\textbf{Conclusion :} $f$ est strictement décroissante sur $[-1\,;1]$ et strictement croissante sur $[1\,;4]$.

\textbf{2. Tableau de variations.}

Calculons les images aux bornes et en $1$ :
\[f(-1)=(-1)^2-2(-1)-3=1+2-3=0\,;\quad f(1)=1-2-3=-4\,;\quad f(4)=16-8-3=5.\]
\begin{center}
\begin{tabular}{|c|ccccc|}\hline
$x$ & $-1$ & & $1$ & & $4$ \\ \hline
$f'(x)$ & & $-$ & $0$ & $+$ & \\ \hline
$f$ & $0$ & $\searrow$ & $-4$ & $\nearrow$ & $5$ \\ \hline
\end{tabular}
\end{center}

\textbf{3. Minimum.}

D'après le tableau de variations, $f$ admet sur $I$ un \cle{minimum} égal à $-4$, atteint en $x=1$. On remarque que $1=-\dfrac{b}{2a}=-\dfrac{-2}{2}$, ce qui confirme le résultat.
\end{exoresolu}

\begin{experience}[Découverte guidée : Lien forme canonique / variations]
Soit $f(x)=2x^2-8x+5$.
\begin{enumerate}
\item Mettre $f$ sous la forme $a(x-\alpha)^2+\beta$.
\item En déduire directement le sens de variation de $f$ sur $\mathbb{R}$ sans calculer $f'(x)$.
\item Vérifier en calculant $f'(x)$.
\end{enumerate}
\tcblower
\textbf{1.} $f(x)=2(x^2-4x)+5=2[(x-2)^2-4]+5=2(x-2)^2-3$. Donc $\alpha=2$, $\beta=-3$, $a=2>0$.
\textbf{2.} Comme $a>0$, le carré $(x-2)^2$ est minimal (nul) en $x=2$. $f$ décroît sur $]-\infty\,;2]$ et croît sur $[2\,;+\infty[$. Minimum $-3$ en $2$.
\textbf{3.} $f'(x)=4x-8=4(x-2)$. Signe de $f'$ : négatif si $x<2$, nul en $2$, positif si $x>2$. Même conclusion.
\end{experience}

\coursec{Étude complète d'une fonction homographique sur un intervalle borné}

\begin{methode}[Étudier $f(x)=\dfrac{ax+b}{cx+d}$ sur un intervalle borné $I$]
\begin{enumerate}
\item \textbf{Ensemble de définition :} Résoudre $cx+d=0 \Rightarrow x=-\frac{d}{c}$ (valeur interdite). $D_f=\mathbb{R}\setminus\{-\frac{d}{c}\}$. Vérifier $I\subset D_f$.
\item \textbf{Asymptote verticale :} La droite $\mathcal{V}$ d'équation $x=-\frac{d}{c}$ est asymptote à la courbe, parallèle à l'axe des ordonnées.
\item \textbf{Dérivée :} $f'(x)=\dfrac{ad-bc}{(cx+d)^2}$. Le dénominateur $>0$ sur $D_f$ $\Rightarrow$ signe de $f'$ = signe constant de $ad-bc$.
\item \textbf{Variations :}
\begin{itemize}
\item Si $ad-bc>0$ : $f$ strictement croissante sur \textbf{chaque} intervalle de $D_f$.
\item Si $ad-bc<0$ : $f$ strictement décroissante sur \textbf{chaque} intervalle de $D_f$.
\end{itemize}
\textbf{Interdiction :} Ne jamais écrire de réunion d'intervalles pour la monotonie.
\item \textbf{Tableau de variations sur $I$ :} Calculer les images des bornes de $I$. Tracer la branche d'hyperbole limitée à $I$.
\end{enumerate}
\end{methode}

\begin{exoresolu}[Étude d'une fonction homographique]
Soit $f$ la fonction définie par $f(x)=\dfrac{2x+1}{x-1}$.
\begin{enumerate}
\item Déterminer l'ensemble de définition de $f$ et une équation de l'asymptote à la courbe de $f$ parallèle à l'axe des ordonnées.
\item Étudier le sens de variation de $f$ sur $I=[2\,;5]$ et dresser son tableau de variations.
\end{enumerate}
\tcblower
\textbf{1. Ensemble de définition et asymptote.}

$f(x)$ existe si et seulement si $x-1\neq0$, c'est-à-dire $x\neq1$. Donc
\[D_f=\mathbb{R}\setminus\{1\}=]-\infty\,;1[\cup]1\,;+\infty[.\]
La droite d'équation $x=1$ est l'\cle{asymptote verticale} (parallèle à l'axe des ordonnées) à la courbe de $f$.

\textbf{2. Variations sur $[2\,;5]$.}

L'intervalle $[2\,;5]$ ne contient pas la valeur interdite $1$, donc $f$ est bien définie et dérivable sur $I$.

Avec $u(x)=2x+1$ et $v(x)=x-1$, on a $u'(x)=2$ et $v'(x)=1$, d'où :
\[f'(x)=\frac{2(x-1)-(2x+1)\times1}{(x-1)^2}=\frac{2x-2-2x-1}{(x-1)^2}=\frac{-3}{(x-1)^2}.\]
Pour tout $x\in[2\,;5]$, $(x-1)^2>0$, donc $f'(x)<0$ : $f$ est \textbf{strictement décroissante} sur $[2\,;5]$.

Images des bornes :
\[f(2)=\frac{4+1}{2-1}=5\qquad;\qquad f(5)=\frac{10+1}{5-1}=\frac{11}{4}=2{,}75.\]
\begin{center}
\begin{tabular}{|c|ccc|}\hline
$x$ & $2$ & & $5$ \\ \hline
$f'(x)$ & & $-$ & \\ \hline
$f$ & $5$ & $\searrow$ & $\dfrac{11}{4}$ \\ \hline
\end{tabular}
\end{center}
\textbf{Conclusion :} $f$ est strictement décroissante sur $[2\,;5]$, de $f(2)=5$ vers $f(5)=\dfrac{11}{4}$.
\end{exoresolu}

\begin{popfigure}
\begin{tikzpicture}[scale=0.8, >=Stealth]
\begin{axis}[
    axis lines=middle,
    xlabel=$x$, ylabel=$y$,
    xmin=-1, xmax=6,
    ymin=-1, ymax=6,
    xtick={1,2,3,4,5}, ytick={1,2,3,4,5},
    domain=2:5, samples=100,
    restrict y to domain=-1:6,
    width=\linewidth, height=0.6\linewidth,
    clip=false
]
\addplot[popcurve, thick, domain=2:5] { (2*x+1)/(x-1) };
\addplot[dashed, popmuted, thick] coordinates {(1,-1) (1,6)} node[above right] {$x=1$};
\addplot[only marks, mark=*, popblue] coordinates {(2,5) (5,2.75)};
\node[popblue, anchor=south west] at (axis cs:2,5) {$(2;5)$};
\node[popblue, anchor=north west] at (axis cs:5,2.75) {$(5;2,75)$};
\end{axis}
\end{tikzpicture}
\legende{Courbe de $f(x)=\frac{2x+1}{x-1}$ sur $[2;5]$ et asymptote verticale $x=1$.}
\end{popfigure}

\coursec{Fonctions associées : Transformations géométriques des courbes}

\courssub{Forme canonique et translation de la parabole}

\begin{methode}[Construire la courbe de $f(x)=a(x-\alpha)^2+\beta$ à partir de $g(x)=ax^2$]
\begin{enumerate}
\item \textbf{Forme canonique :} Si $f(x)=ax^2+bx+c$, compléter le carré pour obtenir $f(x)=a(x-\alpha)^2+\beta$ avec $\alpha=-\frac{b}{2a}$, $\beta=f(\alpha)$.
\item \textbf{Vecteur translation :} La courbe $\mathcal{C}_f$ est l'image de $\mathcal{C}_g$ (parabole $y=ax^2$) par la translation de vecteur $\vec{u}=\alpha\,\vec{\imath}+\beta\,\vec{\jmath}$.
\item \textbf{Éléments remarquables :} Le sommet $O(0;0)$ de $\mathcal{C}_g$ devient le sommet $S(\alpha;\beta)$ de $\mathcal{C}_f$. L'axe de symétrie est la droite $x=\alpha$.
\item \textbf{Tracé :} Placer $S$, calculer $f(\alpha\pm1)$, $f(\alpha\pm2)$ pour la forme, tracer la parabole.
\end{enumerate}
\end{methode}

\begin{exoresolu}[Forme canonique et construction d'une parabole]
Soit $f$ la fonction définie sur $\mathbb{R}$ par $f(x)=2x^2-8x+5$.
\begin{enumerate}
\item Écrire $f(x)$ sous forme canonique.
\item En déduire le vecteur de la translation qui transforme la courbe de $g:x\mapsto 2x^2$ en celle de $f$, puis préciser le sommet et l'axe de symétrie.
\end{enumerate}
\tcblower
\textbf{1. Forme canonique.}
\begin{align*}
f(x)&=2\left(x^2-4x\right)+5\\
&=2\left[(x-2)^2-4\right]+5\\
&=2(x-2)^2-8+5\\
&=2(x-2)^2-3.
\end{align*}
Donc $\alpha=2$, $\beta=-3$. \textbf{Vérification :} $\alpha=-\frac{b}{2a}=2$, $f(2)=-3=\beta$.

\textbf{2. Translation et éléments.}
Vecteur translation : $\vec{u}=2\,\vec{\imath}-3\,\vec{\jmath}$.
Sommet $S(2\,;-3)$. Axe de symétrie : $x=2$. Comme $a=2>0$, parabole tournée vers le haut, minimum $-3$ en $2$.
\end{exoresolu}

\begin{popfigure}
\begin{tikzpicture}[scale=0.7, >=Stealth]
\begin{axis}[
    axis lines=middle, xlabel=$x$, ylabel=$y$,
    xmin=-1, xmax=5, ymin=-5, ymax=5,
    xtick={0,1,2,3,4}, ytick={-3,-2,-1,1,2,3,4},
    width=\linewidth, height=0.6\linewidth,
    clip=false
]
% g(x) = 2x^2
\addplot[popmuted, dashed, thick, domain=-1.5:1.5, samples=100] {2*x^2};
\node[popmuted, anchor=south] at (axis cs:1.5, 4.5) {$g(x)=2x^2$};
% f(x) = 2(x-2)^2 - 3
\addplot[popcurve, thick, domain=0:4, samples=100] {2*(x-2)^2 - 3};
\node[popcurve, anchor=south west] at (axis cs:4, 5) {$f(x)=2(x-2)^2-3$};
% Translation vector arrow
\draw[poporange, thick, ->] (axis cs:0,0) -- (axis cs:2,-3) node[midway, above right, poporange] {$\vec{u}(2;-3)$};
% Points
\addplot[only marks, mark=*, poporange] coordinates {(0,0) (2,-3)};
\node[poporange, anchor=north east] at (axis cs:0,0) {$O$};
\node[poporange, anchor=north] at (axis cs:2,-3) {$S$};
\end{axis}
\end{tikzpicture}
\legende{Translation de la parabole $y=2x^2$ par $\vec{u}(2;-3)$ pour obtenir $f$.}
\end{popfigure}

\courssub{Écriture réduite d'une fonction homographique et translation}

\begin{methode}[Écrire $f(x)=\dfrac{ax+b}{cx+d}$ sous la forme $\alpha+\dfrac{k}{x-\beta}$ et construire sa courbe]
\begin{enumerate}
\item \textbf{Identifier $\beta$} : valeur interdite $\beta=-\dfrac{d}{c}$.
\item \textbf{Division (méthode recommandée) :} Écrire le numérateur en fonction du dénominateur :
\[ax+b = \frac{a}{c}(cx+d) + \left(b-\frac{ad}{c}\right).\]
Diviser par $cx+d$ : $f(x)=\frac{a}{c} + \frac{bc-ad}{c^2}\cdot\frac{1}{x+\frac{d}{c}}$.
Donc $\alpha=\frac{a}{c}$, $k=\frac{bc-ad}{c^2}$.
\item \textbf{Transformation :} La courbe de $f$ est l'image de l'hyperbole de référence $h(x)=\dfrac{k}{x}$ par la translation de vecteur $\vec{u}=\beta\,\vec{\imath}+\alpha\,\vec{\jmath}$.
\item \textbf{Conséquences :}
\begin{itemize}
\item Centre de symétrie : $\Omega(\beta\,;\alpha)$.
\item Asymptotes : verticale $x=\beta$, horizontale $y=\alpha$.
\end{itemize}
\end{enumerate}
\end{methode}

\begin{exoresolu}[Écriture réduite et construction d'une hyperbole]
Soit $f$ la fonction définie par $f(x)=\dfrac{3x-1}{x-2}$.
\begin{enumerate}
\item Montrer que pour tout $x\neq2$, $f(x)=3+\dfrac{5}{x-2}$.
\item En déduire le vecteur de translation, les asymptotes et le centre de symétrie.
\end{enumerate}
\tcblower
\textbf{1. Écriture réduite.}
Pour $x\neq2$ :
\[3+\frac{5}{x-2}=\frac{3(x-2)+5}{x-2}=\frac{3x-6+5}{x-2}=\frac{3x-1}{x-2}=f(x).\]
On a bien $\alpha=3$, $k=5$, $\beta=2$.

\textbf{2. Translation, asymptotes, centre.}
La courbe de $f$ est l'image de $h(x)=\frac{5}{x}$ par la translation de vecteur $\vec{u}=2\,\vec{\imath}+3\,\vec{\jmath}$.
\begin{itemize}
\item Asymptote verticale $x=0 \to x=2$.
\item Asymptote horizontale $y=0 \to y=3$.
\item Centre de symétrie $O(0;0) \to \Omega(2;3)$.
\end{itemize}
\end{exoresolu}

\begin{popfigure}
\begin{tikzpicture}[scale=0.7, >=Stealth]
\begin{axis}[
    axis lines=middle, xlabel=$x$, ylabel=$y$,
    xmin=-1, xmax=5, ymin=-1, ymax=6,
    xtick={0,1,2,3,4}, ytick={0,1,2,3,4,5},
    width=\linewidth, height=0.6\linewidth,
    clip=false,
    restrict y to domain=-2:6
]
% h(x) = 5/x
\addplot[popmuted, dashed, thick, domain=0.3:4, samples=100] {5/x};
\addplot[popmuted, dashed, thick, domain=-1:-0.3, samples=100] {5/x};
\node[popmuted, anchor=south west] at (axis cs:3, 1.8) {$h(x)=5/x$};
% Asymptotes of h
\draw[popmuted, thin, dashed] (axis cs:0,-1) -- (axis cs:0,6);
\draw[popmuted, thin, dashed] (axis cs:-1,0) -- (axis cs:5,0);
% f(x) = 3 + 5/(x-2)
\addplot[popcurve, thick, domain=2.3:5, samples=100] {3 + 5/(x-2)};
\addplot[popcurve, thick, domain=-1:1.7, samples=100] {3 + 5/(x-2)};
\node[popcurve, anchor=south west] at (axis cs:4, 4.5) {$f(x)=3+\frac{5}{x-2}$};
% Asymptotes of f
\draw[poporange, thick, dashed] (axis cs:2,-1) -- (axis cs:2,6) node[above] {$x=2$};
\draw[poporange, thick, dashed] (axis cs:-1,3) -- (axis cs:5,3) node[right] {$y=3$};
% Center
\addplot[only marks, mark=*, poporange] coordinates {(2,3)};
\node[poporange, anchor=south west] at (axis cs:2,3) {$\Omega(2;3)$};
% Translation arrow
\draw[poporange, thick, ->] (axis cs:0,0) -- (axis cs:2,3) node[midway, above left, poporange] {$\vec{u}(2;3)$};
\end{axis}
\end{tikzpicture}
\legende{Construction de l'hyperbole $f$ par translation de $h(x)=5/x$.}
\end{popfigure}

\courssub{Construction des courbes des fonctions associées : $-f$, $|f|$, $f(x-a)$, $f(x)+b$}

\begin{methode}[Construire la courbe d'une fonction associée à partir de $\mathcal{C}_f$]
Soit $\mathcal{C}_f$ la courbe de $f$ dans un repère $(O\,;\vec{\imath},\vec{\jmath})$.
\begin{enumerate}
\item \textbf{$x\mapsto -f(x)$ :} Symétrie de $\mathcal{C}_f$ par rapport à l'\cle{axe des abscisses}. $(x;y)\to(x;-y)$.
\item \textbf{$x\mapsto |f(x)|$ :} Conserver la partie de $\mathcal{C}_f$ \textbf{au-dessus} de l'axe des abscisses ($f(x)\ge0$). Réfléchir (symétrie axiale) la partie \textbf{en-dessous} ($f(x)<0$). La courbe obtenue est entièrement au-dessus de l'axe.
\item \textbf{$x\mapsto f(x-a)$ :} Translation de $\mathcal{C}_f$ de vecteur $a\,\vec{\imath}$ (vers la droite si $a>0$). \emph{Astuce : « le signe à l'intérieur est trompeur »}.
\item \textbf{$x\mapsto f(x)+b$ :} Translation de $\mathcal{C}_f$ de vecteur $b\,\vec{\jmath}$ (vers le haut si $b>0$).
\item \textbf{$x\mapsto f(x-a)+b$ :} Translation de vecteur $a\,\vec{\imath}+b\,\vec{\jmath}$.
\end{enumerate}
\end{methode}

\begin{exoresolu}[Tracés de fonctions associées]
On donne $f$ définie sur $[-3\,;3]$ par $f(x)=x^2-2$. $\mathcal{C}_f$ sa courbe.
\begin{enumerate}
\item Calculer $f(-3)$, $f(0)$, $f(3)$ et résoudre $f(x)=0$ dans $[-3\,;3]$.
\item Expliquer comment construire, à partir de $\mathcal{C}_f$, les courbes de $g(x)=-f(x)$, $h(x)=|f(x)|$ et $p(x)=f(x-1)+2$.
\end{enumerate}
\tcblower
\textbf{1. Calculs préliminaires.}
$f(-3)=7$, $f(0)=-2$, $f(3)=7$.
$f(x)=0 \iff x^2=2 \iff x=\pm\sqrt{2}\in[-3\,;3]$.
$\mathcal{C}_f$ coupe l'axe des abscisses en $A(-\sqrt{2};0)$, $B(\sqrt{2};0)$, sommet $S(0;-2)$, points $(-3;7)$, $(3;7)$.
$f(x)<0$ sur $]-\sqrt{2}\,;\sqrt{2}[$, $f(x)>0$ sur $[-3\,;-\sqrt{2}[\cup]\sqrt{2}\,;3]$.

\textbf{2. Constructions.}

\textbf{Courbe de $g(x)=-f(x)=-x^2+2$ :} Symétrie axiale de $\mathcal{C}_f$ par rapport à l'axe des abscisses.
$S(0;-2)\to S'(0;2)$, $(-3;7)\to(-3;-7)$, $(3;7)\to(3;-7)$, $A,B$ fixes. Parabole tournée vers le bas.

\textbf{Courbe de $h(x)=|f(x)|$ :}
\begin{itemize}
\item Sur $[-3\,;-\sqrt{2}]\cup[\sqrt{2}\,;3]$ : $f\ge0 \Rightarrow h=f$ (on conserve).
\item Sur $[-\sqrt{2}\,;\sqrt{2}]$ : $f\le0 \Rightarrow h=-f$ (on symétrise l'arc passant par $S$). Nouvel arc passant par $(0;2)$.
\end{itemize}
Courbe entièrement au-dessus de l'axe.

\textbf{Courbe de $p(x)=f(x-1)+2$ :} Translation de $\mathcal{C}_f$ par $\vec{u}=\vec{\imath}+2\vec{\jmath}$ ($a=1,b=2$).
$S(0;-2)\to S''(1;0)$, $A\to(1-\sqrt{2};2)$, $B\to(1+\sqrt{2};2)$.
\end{exoresolu}

\begin{popfigure}
\begin{tikzpicture}[scale=0.65, >=Stealth]
\begin{axis}[
    axis lines=middle, xlabel=$x$, ylabel=$y$,
    xmin=-4, xmax=4, ymin=-3, ymax=4,
    xtick={-3,-2,-1,1,2,3}, ytick={-2,-1,1,2,3},
    width=0.48\linewidth, height=0.45\linewidth,
    clip=false, title style={font=\small}, title={$\mathcal{C}_f : y=x^2-2$}
]
\addplot[popcurve, thick, domain=-3:3] {x^2-2};
\addplot[only marks, mark=*, popblue] coordinates {(-1.414,0) (1.414,0) (0,-2) (-3,7) (3,7)};
\end{axis}
\begin{axis}[
    axis lines=middle, xlabel=$x$, ylabel=$y$,
    xmin=-4, xmax=4, ymin=-8, ymax=3,
    xtick={-3,-2,-1,1,2,3}, ytick={-7,-2,2},
    width=0.48\linewidth, height=0.45\linewidth,
    clip=false, title style={font=\small}, title={$\mathcal{C}_g : y=-x^2+2$}, at={(current axis.right of origin)}, anchor=left of origin, xshift=1cm
]
\addplot[poporange, thick, domain=-3:3] {-x^2+2};
\addplot[only marks, mark=*, poporange] coordinates {(-1.414,0) (1.414,0) (0,2) (-3,-7) (3,-7)};
\end{axis}
\begin{axis}[
    axis lines=middle, xlabel=$x$, ylabel=$y$,
    xmin=-4, xmax=4, ymin=-1, ymax=4,
    xtick={-3,-2,-1,1,2,3}, ytick={1,2,3},
    width=0.48\linewidth, height=0.45\linewidth,
    clip=false, title style={font=\small}, title={$\mathcal{C}_h : y=|x^2-2|$}, at={(current axis.below origin)}, anchor=above origin, yshift=-1cm
]
\addplot[popgreen, thick, domain=-3:-1.414] {x^2-2};
\addplot[popgreen, thick, domain=-1.414:1.414] {-x^2+2};
\addplot[popgreen, thick, domain=1.414:3] {x^2-2};
\addplot[only marks, mark=*, popgreen] coordinates {(-1.414,0) (1.414,0) (0,2) (-3,7) (3,7)};
\end{axis}
\begin{axis}[
    axis lines=middle, xlabel=$x$, ylabel=$y$,
    xmin=-2, xmax=6, ymin=-1, ymax=10,
    xtick={-1,0,1,2,3,4,5}, ytick={2,4,6,8},
    width=0.48\linewidth, height=0.45\linewidth,
    clip=false, title style={font=\small}, title={$\mathcal{C}_p : y=(x-1)^2$}, at={(current axis.right of origin)}, anchor=left of origin, xshift=1cm
]
\addplot[poppurple, thick, domain=-2:4] {(x-1)^2};
\addplot[only marks, mark=*, poppurple] coordinates {(1-1.414,2) (1+1.414,2) (1,0) (-2,9) (4,9)};
\draw[poppurple, thick, ->] (axis cs:0,-2) -- (axis cs:1,0) node[midway, above right, poppurple] {$\vec{u}(1;2)$};
\end{axis}
\end{tikzpicture}
\legende{Transformations successives : $-f$ (symétrie), $|f|$ (pliage), $f(x-1)+2$ (translation).}
\end{popfigure}

\coursec{Résolution graphique d'équations et d'inéquations ; Paramètre m}

\begin{methode}[Résoudre graphiquement $f(x)=m$, $f(x)\geq m$ sur un intervalle borné $I$]
\begin{enumerate}
\item \textbf{Tracer soigneusement} $\mathcal{C}_f$ sur $I$ (tableau de variations, points remarquables) et la droite horizontale $\Delta : y=m$.
\item \textbf{Équation $f(x)=m$ :} Les solutions sont les \cle{abscisses} des points d'intersection de $\mathcal{C}_f$ et $\Delta$ appartenant à $I$. Lire les valeurs sur l'axe des abscisses.
\item \textbf{Nombre de solutions selon $m$ :} Faire « glisser » mentalement $\Delta$ verticalement. Comparer $m$ aux extremums de $f$ sur $I$ et aux images des bornes $f(u), f(v)$.
\item \textbf{Inéquation $f(x)\geq m$ :} Solutions = abscisses des points de $\mathcal{C}_f$ situés \textbf{au-dessus} de $\Delta$ (inclus si $\ge$). Pour $f(x)\leq m$ : en dessous. Pour $f(x)=0$ ou $\ge0$ : $\Delta$ est l'axe des abscisses.
\item \textbf{Réponse :} Ensemble de solutions $\mathcal{S}$ sous forme d'intervalle(s), en précisant « lecture graphique » et en restant dans $I$.
\end{enumerate}
\end{methode}

\begin{exoresolu}[Lecture graphique avec paramètre]
Soit $f$ définie sur $I=[-2\,;4]$ par $f(x)=x^2-2x-1$. On admet : $f$ décroissante sur $[-2\,;1]$, croissante sur $[1\,;4]$, $f(-2)=7$, $f(1)=-2$, $f(4)=7$.
\begin{enumerate}
\item Dresser le tableau de variations de $f$ sur $I$.
\item Déterminer graphiquement, suivant les valeurs de $m\in\mathbb{R}$, le nombre de solutions de $f(x)=m$ dans $I$.
\item Résoudre graphiquement $f(x)\geq0$ dans $I$ (on donne $1-\sqrt{2}\approx-0{,}4$, $1+\sqrt{2}\approx2{,}4$).
\end{enumerate}
\tcblower
\textbf{1. Tableau de variations.}
\begin{center}
\begin{tabular}{|c|ccccc|}\hline
$x$ & $-2$ & & $1$ & & $4$ \\ \hline
$f$ & $7$ & $\searrow$ & $-2$ & $\nearrow$ & $7$ \\ \hline
\end{tabular}
\end{center}

\textbf{2. Nombre de solutions de $f(x)=m$.}
Intersections de $\mathcal{C}_f$ (parabole sommet $S(1;-2)$, bornes $(-2;7)$, $(4;7)$) avec $y=m$ :
\begin{itemize}
\item $m < -2$ : $\Delta$ sous le minimum $\Rightarrow$ \textbf{0 solution}.
\item $m = -2$ : $\Delta$ tangente au sommet $\Rightarrow$ \textbf{1 solution} ($x=1$).
\item $-2 < m \le 7$ : $\Delta$ coupe les deux branches $\Rightarrow$ \textbf{2 solutions} (une dans $[-2;1[$, une dans $]1;4]$).
\item $m > 7$ : $\Delta$ au-dessus des bornes $\Rightarrow$ \textbf{0 solution} dans $I$.
\end{itemize}

\textbf{3. Inéquation $f(x)\geq0$.}
Zéros : $x^2-2x-1=0 \iff x=1\pm\sqrt{2}$. Dans $I$ : $\alpha=1-\sqrt{2}\approx-0,4$, $\beta=1+\sqrt{2}\approx2,4$.
Parabole tournée vers le haut $\Rightarrow$ courbe au-dessus de l'axe \textbf{en dehors} des racines.
\[\mathcal{S}=[-2\,;\,1-\sqrt{2}]\cup[1+\sqrt{2}\,;\,4].\]
\end{exoresolu}

\begin{popfigure}
\begin{tikzpicture}[scale=0.7, >=Stealth]
\begin{axis}[
    axis lines=middle, xlabel=$x$, ylabel=$y$,
    xmin=-2.5, xmax=4.5, ymin=-3, ymax=8,
    xtick={-2,-1,1,2,3,4}, ytick={-2,0,2,4,6,7},
    width=\linewidth, height=0.6\linewidth,
    clip=false
]
% f(x)
\addplot[popcurve, thick, domain=-2:4] {x^2-2*x-1};
% Horizontal lines for m
\addplot[popmuted, dashed, domain=-2:4] {-2} node[right, popmuted] {$m=-2$};
\addplot[popblue, dashed, domain=-2:4] {3} node[right, popblue] {$m=3$};
\addplot[poporange, dashed, domain=-2:4] {7} node[right, poporange] {$m=7$};
% Points
\addplot[only marks, mark=*, popblue] coordinates {(-2,7) (1,-2) (4,7)};
\node[popblue, anchor=south west] at (axis cs:-2,7) {$(-2;7)$};
\node[popblue, anchor=north] at (axis cs:1,-2) {$S(1;-2)$};
\node[popblue, anchor=south east] at (axis cs:4,7) {$(4;7)$};
% Roots
\addplot[only marks, mark=*, popgreen] coordinates {({1-sqrt(max(0,2))},0) ({1+sqrt(max(0,2))},0)};
\node[popgreen, anchor=north] at (axis cs:{1-sqrt(2)},0) {$\alpha$};
\node[popgreen, anchor=north] at (axis cs:{1+sqrt(2)},0) {$\beta$};
% Shaded region for f(x) >= 0
\addplot[popgreen!20, domain=-2:{1-sqrt(max(0,2))}] {x^2-2*x-1} \closedcycle;
\addplot[popgreen!20, domain={1+sqrt(max(0,2))}:4] {x^2-2*x-1} \closedcycle;
\end{axis}
\end{tikzpicture}
\legende{Analyse graphique du nombre de solutions de $f(x)=m$ et résolution de $f(x)\ge0$.}
\end{popfigure}

\coursec{Synthèse et stratégie globale d'étude de fonction}

\begin{aretenir}[Check-list « Étude complète d'une fonction sur un intervalle borné $I$ »]
\textbf{Étape 1 : Ensemble de définition $D_f$ et intervalle d'étude $I$.}
\begin{itemize}
\item Polynôme : $D_f=\mathbb{R}$, $I$ donné.
\item Homographique : $D_f=\mathbb{R}\setminus\{x_0\}$ (valeur interdite). Vérifier $x_0\notin I$.
\end{itemize}
\textbf{Étape 2 : Asymptotes (homographique seulement).}
Asymptote verticale $x=x_0$. Asymptote horizontale $y=\alpha$ (via forme réduite).
\textbf{Étape 3 : Dérivée et signe.}
$f'(x)$ calculée, signe étudié sur $I$ (tableau de signes de $f'$).
\textbf{Étape 4 : Tableau de variations.}
3 lignes ($x$, $f'$, $f$). Colonnes : bornes de $I$ + valeurs où $f'=0$ dans $I$. Calculer images exactes.
\textbf{Étape 5 : Extremums.}
Min/Max sur $I$ (valeurs et abscisses). Distinguer extremum relatif (intérieur) et absolu (bornes).
\textbf{Étape 6 : Représentation graphique.}
Repère orthonormé, échelle adaptée. Points caractéristiques, asymptotes (pointillés), branche(s) de courbe.
\textbf{Étape 7 : Résolution graphique (si demandée).}
Équation $f(x)=m$ (abscisses intersections), Inéquation $f(x)\ge m$ (abscisses au-dessus).
\end{aretenir}

\begin{exemplebox}[Stratégie au Probatoire : Gestion du temps]
\begin{enumerate}
\item \textbf{Lisez tout le sujet} : repérez l'intervalle $I$, la fonction ($f$ ou $g$), les questions (tableau, graphique, équation).
\item \textbf{Brouillon :} Calculez $D_f$, $f'$, signe de $f'$, valeurs aux bornes et points critiques. Vérifiez la cohérence (ex: $f(1)=-4$ compatible avec décroissant puis croissant ?).
\item \textbf{Rédaction :} Phrases complètes pour les conclusions de variations. Tableau de variations \textbf{au stylo}, à la règle. Courbe \textbf{au crayon}, point par point, traits continus.
\item \textbf{Lecture graphique :} Tracez la droite $y=m$ en pointillés. Entourez les intersections. Lisez les \textbf{abscisses} sur l'axe $Ox$. Écrivez $\mathcal{S}=[\dots]$.
\end{enumerate}
\end{exemplebox}

\coursec{Modélisation et problèmes concrets (Applications camerounaises)}

\begin{exoresolu}[Optimisation d'un enclos agricole (Fonction polynôme)]
Un agriculteur de la région de l'Ouest dispose de \SI{100}{m} de grillage pour clôturer un enclos rectangulaire contre un mur existant (un seul côté non grillagé). On note $x$ (en mètres) la longueur du côté perpendiculaire au mur.
\begin{enumerate}
\item Exprimer la surface $S(x)$ de l'enclos en fonction de $x$. Préciser l'intervalle de définition physique $I$.
\item Étudier les variations de $S$ sur $I$ et déterminer les dimensions de l'enclos de surface maximale.
\item Quelle est cette surface maximale ?
\end{enumerate}
\tcblower
\textbf{1. Modélisation.}
Le grillage sert pour les deux largeurs ($x$ chacune) et une longueur $L$. $2x+L=100 \Rightarrow L=100-2x$.
Surface $S(x)=x\times L = x(100-2x) = -2x^2+100x$.
Contraintes physiques : $x>0$ et $L>0 \Rightarrow 100-2x>0 \Rightarrow x<50$.
Donc $I=[0\,;50]$ (bornes incluses pour l'étude, bien que $x=0$ ou $50$ donnent surface nulle).

\textbf{2. Étude de $S(x)=-2x^2+100x$ sur $[0\,;50]$.}
$S'(x)=-4x+100$.
$S'(x)=0 \iff x=25 \in [0\,;50]$.
$a=-2<0 \Rightarrow S'$ positif puis négatif.
$S$ croissante sur $[0\,;25]$, décroissante sur $[25\,;50]$.
Tableau de variations :
\begin{center}
\begin{tabular}{|c|ccccc|}\hline
$x$ & $0$ & & $25$ & & $50$ \\ \hline
$S'(x)$ & & $+$ & $0$ & $-$ & \\ \hline
$S$ & $0$ & $\nearrow$ & $1250$ & $\searrow$ & $0$ \\ \hline
\end{tabular}
\end{center}

\textbf{3. Conclusion.}
Surface maximale $1250\,\text{m}^2$ atteinte pour $x=25\,\text{m}$.
Longueur $L=100-2\times25=50\,\text{m}$.
Dimensions optimales : \textbf{25 m (largeur) $\times$ 50 m (longueur)}.
\end{exoresolu}

\begin{exoresolu}[Coût de raccordement électrique (Fonction homographique)]
CAMTEL doit raccorder un village. Le coût total $C(x)$ (en millions de FCFA) pour installer $x$ km de câble est modélisé par $C(x)=\dfrac{50x+200}{x+2}$ pour $x\in[0\,;20]$.
\begin{enumerate}
\item Déterminer l'ensemble de définition de $C$ et vérifier que $[0\,;20]\subset D_C$.
\item Écrire $C(x)$ sous la forme $\alpha+\dfrac{k}{x-\beta}$. En déduire les asymptotes.
\item Étudier les variations de $C$ sur $[0\,;20]$. Interpréter $\alpha$ économiquement.
\item Calculer $C(0)$ et $C(20)$. Que représente $C(0)$ ?
\end{enumerate}
\tcblower
\textbf{1. Ensemble de définition.}
Dénominateur $x+2=0 \iff x=-2$. $D_C=\mathbb{R}\setminus\{-2\}$. Comme $-2\notin[0\,;20]$, $C$ est bien définie sur $I$.

\textbf{2. Forme réduite et asymptotes.}
$C(x)=\dfrac{50x+200}{x+2}=\dfrac{50(x+2)+100}{x+2}=50+\dfrac{100}{x+2}$.
Donc $\alpha=50$, $k=100$, $\beta=-2$.
Translation de $h(x)=\frac{100}{x}$ par $\vec{u}(-2;50)$.
Asymptotes : verticale $x=-2$ (hors intervalle), horizontale $y=50$.

\textbf{3. Variations.}
$C'(x)=\frac{ad-bc}{(x+2)^2}=\frac{50\times2 - 200\times1}{(x+2)^2}=\frac{-100}{(x+2)^2}<0$ sur $[0\,;20]$.
$C$ est \textbf{strictement décroissante} sur $[0\,;20]$.
Tableau :
\begin{center}
\begin{tabular}{|c|ccc|}\hline
$x$ & $0$ & & $20$ \\ \hline
$C'(x)$ & & $-$ & \\ \hline
$C$ & $100$ & $\searrow$ & $50+\frac{100}{22}\approx54,5$ \\ \hline
\end{tabular}
\end{center}
\textbf{Interprétation :} $\alpha=50$ millions FCFA est le \cle{coût marginal asymptotique} (coût par km pour de très grandes distances, les frais fixes $200$ M sont « dilués »).

\textbf{4. Valeurs aux bornes.}
$C(0)=\frac{200}{2}=100$ : Coût fixe d'installation (borne, transformateur) sans câble.
$C(20)=50+\frac{100}{22}\approx54,55$ : Coût moyen au km pour 20 km $\approx 2,73$ M FCFA/km.
\end{exoresolu}

\begin{exercice}[Exercice d'entraînement : Modélisation télécoms]
Une antenne relais MTN émet un signal de puissance $P_0$. La puissance reçue $P(d)$ (en $\mu$W) à une distance $d$ (en km) est modélisée par $P(d)=\frac{k}{d^2}$ pour $d\ge1$. On mesure $P(2)=25$.
\begin{enumerate}
\item Déterminer $k$ et écrire l'expression de $P(d)$.
\item Étudier les variations de $P$ sur $[1\,;10]$. Que se passe-t-il quand on s'éloigne ?
\item Résoudre graphiquement l'inéquation $P(d)\geq 4$ sur $[1\,;10]$. Interpréter.
\end{enumerate}
\end{exercice}

\begin{corrige}[Exercice d'entraînement]
\textbf{1.} $P(2)=\frac{k}{4}=25 \Rightarrow k=100$. $P(d)=\frac{100}{d^2}$ sur $[1\,;10]$.
\textbf{2.} $P'(d)=-\frac{200}{d^3}<0$ sur $[1\,;10]$. $P$ strictement décroissante.
Tableau : $P(1)=100$, $P(10)=1$. La puissance chute très vite au début, puis se stabilise faiblement.
\textbf{3.} $P(d)\ge4 \iff \frac{100}{d^2}\ge4 \iff d^2\le25 \iff d\le5$ (car $d>0$).
Sur $[1\,;10]$, $\mathcal{S}=[1\,;5]$.
\textbf{Interprétation :} Pour capter un signal correct ($\ge4\mu$W), il faut être à moins de 5 km de l'antenne.
\end{corrige}

\begin{attention}[Les 5 erreurs qui coûtent des points au Probatoire]
\begin{enumerate}
\item \textbf{Oublier la valeur interdite.} Avant toute étude d'une fonction homographique, on doit écrire l'ensemble de définition. Une dérivée calculée sur un intervalle contenant la valeur interdite invalide toute l'étude. Vérifie toujours que l'intervalle borné proposé ne contient pas $-\dfrac{d}{c}$.
\item \textbf{Se tromper de signe dans la dérivée homographique.} $f'(x)=\dfrac{ad-bc}{(cx+d)^2}$ : une erreur de signe sur $ad-bc$ inverse tout le sens de variation. Détaille le calcul $u'v-uv'$ au brouillon et vérifie avec une valeur numérique (par exemple le sens entre deux images calculées).
\item \textbf{Écrire une réunion d'intervalles pour une fonction homographique.} On dit « $f$ est décroissante sur $]-\infty\,;1[$ \textbf{et} sur $]1\,;+\infty[$ », jamais « sur $]-\infty\,;1[\cup]1\,;+\infty[$ » : la monotonie ne se formule que sur un intervalle.
\item \textbf{Confondre les transformations.} La courbe de $f(x-a)$ se déplace de $+a$ horizontalement (sens opposé au signe écrit), celle de $f(x)+b$ de $+b$ verticalement (même signe). Pour $|f|$, seule la partie \emph{sous} l'axe des abscisses est réfléchie ; pour $-f$, c'est toute la courbe qui est symétrisée.
\item \textbf{Donner des ordonnées au lieu d'abscisses en lecture graphique.} Les solutions de $f(x)=m$ ou de $f(x)\geq0$ sont des valeurs de $x$, lues sur l'axe des abscisses, et doivent être présentées sous forme d'ensemble $\mathcal{S}=\dots$ en restant dans l'intervalle borné imposé par l'énoncé.
\end{enumerate}
\end{attention}

\newpage

\coursec{Exercices}

\courssub{Je vérifie mes connaissances}

\begin{exercice}[QCM : dérivée et variations]
Pour chaque question, une seule réponse est exacte. Recopie la bonne lettre.
\begin{enumerate}
\item La fonction dérivée de $f(x)=2x^{2}-8x+5$ est :
\begin{enumerate}
\item[(a)] $f'(x)=2x-8$ \quad (b) $f'(x)=4x-8$ \quad (c) $f'(x)=4x+5$
\end{enumerate}
\item La fonction $g(x)=\dfrac{3x-1}{x-2}$ est définie sur :
\begin{enumerate}
\item[(a)] $\mathbb{R}$ \quad (b) $\mathbb{R}\setminus\{2\}$ \quad (c) $\mathbb{R}\setminus\left\{\dfrac{1}{3}\right\}$
\end{enumerate}
\item Si $f'(x)<0$ pour tout $x\in]1\,;\,4[$, alors $f$ est :
\begin{enumerate}
\item[(a)] croissante sur $[1\,;\,4]$ \quad (b) décroissante sur $[1\,;\,4]$ \quad (c) constante sur $[1\,;\,4]$
\end{enumerate}
\item La courbe de $h(x)=\dfrac{5}{x-3}$ admet pour asymptote parallèle à l'axe des ordonnées la droite d'équation :
\begin{enumerate}
\item[(a)] $x=3$ \quad (b) $y=3$ \quad (c) $x=5$
\end{enumerate}
\end{enumerate}
\end{exercice}

\begin{exercice}[Vrai ou faux, justifié]
Réponds par Vrai ou Faux en justifiant chaque réponse.
\begin{enumerate}
\item La forme canonique de $f(x)=x^{2}-4x+3$ est $f(x)=(x-2)^{2}-1$.
\item La courbe de la fonction $x\mapsto f(x)+2$ s'obtient à partir de celle de $f$ par une translation de vecteur $2\vec{i}$.
\item La courbe de la fonction $x\mapsto -f(x)$ est symétrique de celle de $f$ par rapport à l'axe des abscisses.
\item Si $f$ admet un minimum en $x=2$ sur $[0\,;\,5]$, alors $f'(2)=0$.
\end{enumerate}
\end{exercice}

\begin{exercice}[Ensembles de définition]
Détermine l'ensemble de définition de chacune des fonctions suivantes :
\begin{enumerate}
\item $f(x)=3x^{2}-7x+2$ ;
\item $g(x)=\dfrac{2x+5}{x-4}$ ;
\item $h(x)=\dfrac{x-1}{2x+6}$ ;
\item $k(x)=\dfrac{4}{x^{2}-9}$.
\end{enumerate}
\end{exercice}

\begin{exercice}[Calculs de dérivées]
Calcule la fonction dérivée de chacune des fonctions suivantes, puis étudie son signe sur l'intervalle indiqué.
\begin{enumerate}
\item $f(x)=x^{2}-6x+1$ sur $[0\,;\,6]$ ;
\item $g(x)=-2x^{2}+4x+3$ sur $[-1\,;\,3]$ ;
\item $h(x)=\dfrac{2x+1}{x-1}$ sur $[2\,;\,5]$.
\end{enumerate}
\end{exercice}

\courssub{Je m'entraîne}

\begin{exercice}[Forme canonique et extremum]
Soit $f$ la fonction définie sur $[0\,;\,5]$ par $f(x)=x^{2}-6x+10$.
\begin{enumerate}
\item Écris $f(x)$ sous forme canonique.
\item Déduis-en le sens de variation de $f$ sur $[0\,;\,5]$ sans calculer la dérivée.
\item Détermine l'extremum de $f$ sur $[0\,;\,5]$ : précise sa nature, sa valeur et le point où il est atteint.
\item Dresse le tableau des variations de $f$ sur $[0\,;\,5]$.
\end{enumerate}
\end{exercice}

\begin{exercice}[Fonction homographique : écriture réduite]
Soit $f$ la fonction définie par $f(x)=\dfrac{2x+3}{x-1}$.
\begin{enumerate}
\item Détermine l'ensemble de définition de $f$.
\item Montre que pour tout $x\neq 1$, $f(x)=2+\dfrac{5}{x-1}$.
\item Donne une équation de chacune des deux asymptotes à la courbe de $f$.
\item Calcule $f'(x)$ et étudie son signe.
\item Dresse le tableau des variations de $f$ sur l'intervalle $[2\,;\,6]$.
\end{enumerate}
\end{exercice}

\begin{exercice}[Transformations de courbes]
Soit $f$ la fonction définie par $f(x)=x^{2}$, de courbe représentative $\mathcal{C}_{f}$.
\begin{enumerate}
\item Pour chacune des fonctions suivantes, indique la transformation géométrique permettant de construire sa courbe à partir de $\mathcal{C}_{f}$ :
\[ g(x)=x^{2}-3 \;;\quad h(x)=(x-2)^{2} \;;\quad k(x)=-x^{2} \;;\quad l(x)=|x^{2}-4| \;;\quad m(x)=(x+1)^{2}+2. \]
\item Pour la fonction $m$, précise le vecteur de la translation et les coordonnées du sommet de la parabole obtenue.
\item Soit $u(x)=\dfrac{1}{x}$. Indique la transformation permettant d'obtenir la courbe de $v(x)=\dfrac{1}{x-3}+2$ à partir de celle de $u$, et donne les équations des asymptotes de la courbe de $v$.
\end{enumerate}
\end{exercice}

\begin{exercice}[Lecture graphique]
La courbe ci-dessous représente une fonction $f$ définie sur $[-1\,;\,5]$.
\begin{center}
\begin{tikzpicture}[scale=0.8]
\draw[->,thick] (-1.5,0) -- (5.5,0) node[below]{$x$};
\draw[->,thick] (0,-2.5) -- (0,6) node[left]{$y$};
\foreach \x in {-1,1,2,3,4,5} \draw (\x,0.08) -- (\x,-0.08) node[below]{\small $\x$};
\foreach \y in {-2,-1,1,2,3,4,5} \draw (0.06,\y) -- (-0.06,\y) node[left]{\small $\y$};
\draw[very thick,popblue,domain=-1:5,samples=100] plot (\x,{0.5*(\x-2)*(\x-2)-1});
\draw[dashed,popmuted] (2,0) -- (2,-1) -- (0,-1);
\node[popblue] at (4.6,4.6) {$\mathcal{C}_{f}$};
\end{tikzpicture}
\end{center}
À l'aide du graphique :
\begin{enumerate}
\item Lis les images de $-1$, $2$ et $5$ par $f$.
\item Résous graphiquement l'équation $f(x)=0$ sur $[-1\,;\,5]$.
\item Résous graphiquement l'inéquation $f(x)\leq 3$ sur $[-1\,;\,5]$.
\item Détermine graphiquement le nombre de solutions de l'équation $f(x)=m$ dans chacun des cas : $m=-2$ ; $m=-1$ ; $m=0$ ; $m=6$.
\end{enumerate}
\end{exercice}

\courssub{Je me prépare au Probatoire}

\begin{exercice}[Type Probatoire : polynôme du second degré et paramètre]
On considère la fonction $f$ définie sur l'intervalle $[0\,;\,4]$ par
\[ f(x)=2x^{2}-8x+5 \]
et on note $\mathcal{C}_{f}$ sa courbe représentative dans un repère orthonormé $(O\,;\vec{i},\vec{j})$ (unité : 1 cm sur chaque axe).
\begin{enumerate}
\item \begin{enumerate}
\item Montre que pour tout réel $x$, $f(x)=2(x-2)^{2}-3$.
\item En déduire le sens de variation de $f$ sur $[0\,;\,2]$ et sur $[2\,;\,4]$.
\end{enumerate}
\item \begin{enumerate}
\item Calcule $f'(x)$ et vérifie que le signe de $f'(x)$ confirme les variations trouvées à la question 1.
\item Dresse le tableau des variations complet de $f$ sur $[0\,;\,4]$ (on y fera figurer $f(0)$, $f(2)$ et $f(4)$).
\item Quel est le minimum de $f$ sur $[0\,;\,4]$ ? En quel point est-il atteint ?
\end{enumerate}
\item Recopie et complète le tableau de valeurs suivant, puis construis soigneusement $\mathcal{C}_{f}$ :
\begin{center}
\begin{tabular}{|c|ccccc|}\hline
$x$ & $0$ & $1$ & $2$ & $3$ & $4$ \\ \hline
$f(x)$ & & & & & \\ \hline
\end{tabular}
\end{center}
\item \begin{enumerate}
\item Résous graphiquement l'équation $f(x)=1$ sur $[0\,;\,4]$.
\item Retrouve les valeurs exactes des solutions par le calcul.
\end{enumerate}
\item Détermine, suivant les valeurs du paramètre réel $m$, le nombre de solutions de l'équation $f(x)=m$ sur $[0\,;\,4]$. On présentera les résultats dans un tableau.
\item Résous graphiquement l'inéquation $f(x)\geq 1$ sur $[0\,;\,4]$.
\end{enumerate}
\end{exercice}

\begin{exercice}[Type Probatoire : fonction homographique]
On considère la fonction $f$ définie par
\[ f(x)=\dfrac{3x-1}{x-2} \]
et on note $\mathcal{C}_{f}$ sa courbe représentative dans un repère orthonormé $(O\,;\vec{i},\vec{j})$ (unité : 1 cm).
\begin{enumerate}
\item Détermine l'ensemble de définition $D_{f}$ de $f$.
\item \begin{enumerate}
\item Montre que pour tout $x\in D_{f}$, $f(x)=3+\dfrac{5}{x-2}$.
\item En déduire que $\mathcal{C}_{f}$ est l'image de la courbe de la fonction $x\mapsto\dfrac{5}{x}$ par une translation dont on précisera le vecteur.
\item Donne les coordonnées du centre de symétrie de $\mathcal{C}_{f}$.
\end{enumerate}
\item Donne une équation de l'asymptote à $\mathcal{C}_{f}$ parallèle à l'axe des ordonnées, puis une équation de l'asymptote parallèle à l'axe des abscisses.
\item \begin{enumerate}
\item Calcule $f'(x)$ pour tout $x\in D_{f}$.
\item Étudie le signe de $f'(x)$ et dresse le tableau des variations de $f$ sur l'intervalle $[3\,;\,6]$.
\end{enumerate}
\item Calcule $f(3)$, $f(4)$, $f(5)$ et $f(6)$, puis construis $\mathcal{C}_{f}$ sur $[3\,;\,6]$ en traçant également les deux asymptotes.
\item \begin{enumerate}
\item Résous graphiquement sur $[3\,;\,6]$ l'inéquation $f(x)\geq 4$.
\item Retrouve le résultat par le calcul.
\end{enumerate}
\item Soit $g$ la fonction définie par $g(x)=f(x)-2$. Sans refaire l'étude, indique comment construire la courbe de $g$ à partir de $\mathcal{C}_{f}$ et donne une équation de son asymptote horizontale.
\end{enumerate}
\end{exercice}

\begin{exercice}[Type Probatoire : fonctions associées et modélisation agricole]
\textbf{Partie A : transformations successives.}

Soit $f$ la fonction définie sur $[-1\,;\,5]$ par $f(x)=x^{2}-4x+3$, de courbe représentative $\mathcal{C}_{f}$.
\begin{enumerate}
\item Écris $f(x)$ sous forme canonique et précise les coordonnées du sommet $S$ de $\mathcal{C}_{f}$.
\item Dresse le tableau des variations de $f$ sur $[-1\,;\,5]$ et construis $\mathcal{C}_{f}$ dans un repère orthonormé (unité : 1 cm).
\item On considère les fonctions définies par :
\[ g(x)=f(x-2) \;;\quad h(x)=-f(x-2) \;;\quad k(x)=\left|-f(x-2)\right| \;;\quad l(x)=\left|-f(x-2)\right|+1. \]
Pour chacune des courbes $\mathcal{C}_{g}$, $\mathcal{C}_{h}$, $\mathcal{C}_{k}$ et $\mathcal{C}_{l}$ :
\begin{enumerate}
\item indique la transformation géométrique qui permet de l'obtenir à partir de la courbe précédente ;
\item donne les coordonnées de l'image du sommet $S$ après cette transformation.
\end{enumerate}
\item Sur quel intervalle la fonction $g$ est-elle définie si l'on veut que $x-2\in[-1\,;\,5]$ ?
\end{enumerate}

\textbf{Partie B : le pré de M. Etoundi.}

M. Etoundi, agriculteur de la région de l'Ouest, veut clôturer un pré rectangulaire de \SI{400}{m^{2}} adossé à une rivière. Le côté longeant la rivière n'a pas besoin de grillage. On note $x$ (en mètres) la longueur de chaque côté perpendiculaire à la rivière, avec $x\in[5\,;\,80]$. Le grillage coûte \SI{2500}{FCFA} le mètre.
\begin{enumerate}
\item Exprime en fonction de $x$ la longueur du côté parallèle à la rivière.
\item Montre que le coût total en FCFA est donné par $C(x)=5000x+\dfrac{1000000}{x}$.
\item On admet que $C'(x)=\dfrac{5000(x^{2}-200)}{x^{2}}$. Étudie le signe de $C'(x)$ sur $[5\,;\,80]$ et dresse le tableau des variations de $C$.
\item Détermine les dimensions du pré qui rendent le coût minimal, et calcule ce coût minimal en FCFA (on donnera une valeur approchée à 100 FCFA près).
\end{enumerate}
\end{exercice}

\newpage

\coursec{Corrigés des exercices}

\coursec{Exercices corrigés : Étude de fonctions et transformations}

\begin{corrige}[Exercice 1 — QCM : dérivée et variations]
\begin{enumerate}
\item \textbf{Réponse (b).} $f(x)=2x^{2}-8x+5$ est une fonction polynôme : on dérive terme à terme, $(x^{2})'=2x$ et $(x)'=1$, donc $f'(x)=2\times 2x-8=4x-8$.
\item \textbf{Réponse (b).} Une fonction homographique est définie partout sauf là où son dénominateur s'annule : $x-2=0\Leftrightarrow x=2$. Donc $D_{g}=\mathbb{R}\setminus\{2\}$. (La valeur $\frac{1}{3}$ annule le \emph{numérateur} : elle donne $g\left(\frac{1}{3}\right)=0$, ce qui est parfaitement licite.)
\item \textbf{Réponse (b).} Théorème fondamental : si $f'(x)<0$ sur un intervalle, alors $f$ est strictement décroissante sur cet intervalle.
\item \textbf{Réponse (a).} La courbe de $x\mapsto\dfrac{k}{x-\beta}$ admet pour asymptote verticale la droite d'équation $x=\beta$. Ici $\beta=3$, donc l'asymptote parallèle à l'axe des ordonnées a pour équation $x=3$.
\end{enumerate}
\end{corrige}

\begin{corrige}[Exercice 2 — Vrai ou faux, justifié]
\begin{enumerate}
\item \textbf{Vrai.} $x^{2}-4x$ est le début du carré $(x-2)^{2}=x^{2}-4x+4$. Donc
\[ f(x)=x^{2}-4x+3=(x-2)^{2}-4+3=(x-2)^{2}-1. \]
\item \textbf{Faux.} La courbe de $x\mapsto f(x)+2$ s'obtient par la translation de vecteur $2\vec{j}$ (translation \emph{verticale}, vers le haut). Le vecteur $2\vec{i}$ correspondrait à la courbe de $x\mapsto f(x-2)$.
\item \textbf{Vrai.} Le point $M(x\,;f(x))$ de $\mathcal{C}_{f}$ a pour correspondant $M'(x\,;-f(x))$ sur la courbe de $-f$ : même abscisse, ordonnées opposées. C'est exactement la symétrie par rapport à l'axe des abscisses.
\item \textbf{Faux en général.} Le résultat n'est valable que si $f$ est \emph{dérivable} en $2$. Contre-exemple : $f(x)=|x-2|$ admet sur $[0\,;\,5]$ un minimum en $x=2$ (de valeur $0$), mais $f$ n'est pas dérivable en $2$, donc $f'(2)$ n'existe pas. \textbf{À retenir :} si $f$ est dérivable sur un intervalle ouvert contenant $2$ et admet un extremum en $2$, alors $f'(2)=0$ ; la réciproque est fausse.
\end{enumerate}
\end{corrige}

\begin{corrige}[Exercice 3 — Ensembles de définition]
\begin{enumerate}
\item $f$ est une fonction polynôme : elle est définie pour tout réel. $D_{f}=\mathbb{R}$.
\item $g$ est homographique : $x-4=0\Leftrightarrow x=4$. Donc $D_{g}=\mathbb{R}\setminus\{4\}$.
\item $2x+6=0\Leftrightarrow x=-3$. Donc $D_{h}=\mathbb{R}\setminus\{-3\}$.
\item $x^{2}-9=0\Leftrightarrow x=3$ ou $x=-3$. Donc $D_{k}=\mathbb{R}\setminus\{-3\,;\,3\}$.
\end{enumerate}
\textbf{Méthode :} Seuls deux obstacles limitent l'ensemble de définition à ce niveau : le dénominateur nul et (plus tard) la quantité sous une racine carrée négative. Un polynôme ne pose jamais de problème.
\end{corrige}

\begin{corrige}[Exercice 4 — Calculs de dérivées]
\begin{enumerate}
\item $f'(x)=2x-6=2(x-3)$. C'est une fonction affine de coefficient directeur positif, nulle en $x=3$. Donc $f'(x)<0$ sur $[0\,;\,3[$ et $f'(x)>0$ sur $]3\,;\,6]$ : $f$ décroît sur $[0\,;\,3]$ puis croît sur $[3\,;\,6]$.
\item $g'(x)=-4x+4=-4(x-1)$, nulle en $x=1$, de coefficient directeur $-4<0$. Donc $g'(x)>0$ sur $[-1\,;\,1[$ et $g'(x)<0$ sur $]1\,;\,3]$ : $g$ croît sur $[-1\,;\,1]$ puis décroît sur $[1\,;\,3]$.
\item $h=\dfrac{u}{v}$ avec $u(x)=2x+1$, $u'(x)=2$, $v(x)=x-1$, $v'(x)=1$. Par la formule $\left(\dfrac{u}{v}\right)'=\dfrac{u'v-uv'}{v^{2}}$ :
\[ h'(x)=\frac{2(x-1)-(2x+1)\times 1}{(x-1)^{2}}=\frac{2x-2-2x-1}{(x-1)^{2}}=\frac{-3}{(x-1)^{2}}. \]
Pour tout $x\in[2\,;\,5]$, $(x-1)^{2}>0$, donc $h'(x)<0$ : $h$ est strictement décroissante sur $[2\,;\,5]$.
\end{enumerate}
\end{corrige}

\begin{corrige}[Exercice 5 — Forme canonique et extremum]
\begin{enumerate}
\item $x^{2}-6x$ est le début de $(x-3)^{2}=x^{2}-6x+9$. Donc
\[ f(x)=(x-3)^{2}-9+10=(x-3)^{2}+1. \]
\item Pour tout réel $x$, $(x-3)^{2}\geq 0$. La quantité $(x-3)^{2}$ diminue quand $x$ se rapproche de $3$ et augmente quand $x$ s'en éloigne. Donc $f$ est \textbf{décroissante sur $[0\,;\,3]$} et \textbf{croissante sur $[3\,;\,5]$}.
\item $(x-3)^{2}\geq 0$ donc $f(x)\geq 1$, avec égalité si et seulement si $x=3$. Comme $3\in[0\,;\,5]$, $f$ admet sur $[0\,;\,5]$ un \textbf{minimum égal à $1$, atteint en $x=3$}.
\item $f(0)=(0-3)^{2}+1=10$ ; $f(3)=1$ ; $f(5)=(5-3)^{2}+1=5$.
\begin{center}
\begin{tabular}{|c|ccccc|}\hline $x$ & $0$ & & $3$ & & $5$ \\ \hline $f'(x)$ & & $-$ & $0$ & $+$ & \\ \hline $f$ & $10$ & $\searrow$ & $1$ & $\nearrow$ & $5$ \\ \hline\end{tabular}
\end{center}
\end{enumerate}
\end{corrige}

\begin{corrige}[Exercice 6 — Fonction homographique : écriture réduite]
\begin{enumerate}
\item $x-1=0\Leftrightarrow x=1$, donc $D_{f}=\mathbb{R}\setminus\{1\}$.
\item Pour tout $x\neq 1$ :
\[ 2+\frac{5}{x-1}=\frac{2(x-1)+5}{x-1}=\frac{2x-2+5}{x-1}=\frac{2x+3}{x-1}=f(x). \]
\item La courbe de $f$ est l'image de l'hyperbole $y=\dfrac{5}{x}$ par la translation de vecteur $\vec{i}+2\vec{j}$. Elle admet donc deux asymptotes : la droite d'équation $x=1$ (parallèle à l'axe des ordonnées) et la droite d'équation $y=2$ (parallèle à l'axe des abscisses).
\item $f(x)=2+5(x-1)^{-1}$ donc $f'(x)=-\dfrac{5}{(x-1)^{2}}$. Pour tout $x\neq 1$, $(x-1)^{2}>0$, donc $f'(x)<0$ : $f$ est strictement décroissante sur chaque intervalle de son ensemble de définition.
\item Sur $[2\,;\,6]$ (qui ne contient pas $1$) : $f(2)=\dfrac{7}{1}=7$ et $f(6)=\dfrac{15}{5}=3$.
\begin{center}
\begin{tabular}{|c|ccc|}\hline $x$ & $2$ & & $6$ \\ \hline $f'(x)$ & & $-$ & \\ \hline $f$ & $7$ & $\searrow$ & $3$ \\ \hline\end{tabular}
\end{center}
\end{enumerate}
\end{corrige}

\begin{corrige}[Exercice 7 — Transformations de courbes]
\begin{enumerate}
\item \begin{itemize}
\item $g(x)=f(x)-3$ : translation de vecteur $-3\vec{j}$ (3 unités vers le bas).
\item $h(x)=f(x-2)$ : translation de vecteur $2\vec{i}$ (2 unités vers la droite).
\item $k(x)=-f(x)$ : symétrie par rapport à l'axe des abscisses.
\item $l(x)=|x^{2}-4|$ : on conserve la partie de la courbe de $x\mapsto x^{2}-4$ située au-dessus de l'axe des abscisses, et on applique la symétrie d'axe $(Ox)$ à la partie située en dessous (pour $|x|<2$).
\item $m(x)=f(x+1)+2$ : translation de vecteur $-\vec{i}+2\vec{j}$.
\end{itemize}
\item Pour $m(x)=(x+1)^{2}+2=(x-(-1))^{2}+2$ : translation de vecteur $\vec{t}=-\vec{i}+2\vec{j}$. Le sommet $O(0\,;\,0)$ de $\mathcal{C}_{f}$ a pour image le point $S(-1\,;\,2)$, sommet de la nouvelle parabole.
\item $v(x)=u(x-3)+2$ : la courbe de $v$ est l'image de celle de $u$ par la translation de vecteur $3\vec{i}+2\vec{j}$. Les asymptotes $x=0$ et $y=0$ de l'hyperbole de $u$ deviennent les droites d'équations $x=3$ et $y=2$.
\end{enumerate}
\end{corrige}

\begin{corrige}[Exercice 8 — Lecture graphique]
La courbe est la parabole de sommet $S(2\,;\,-1)$ passant par $(0\,;\,1)$.
\begin{enumerate}
\item On lit : $f(-1)\approx 3,5$ ; $f(2)=-1$ (sommet) ; $f(5)\approx 3,5$.
\item Les solutions de $f(x)=0$ sont les abscisses des points d'intersection de $\mathcal{C}_{f}$ avec l'axe des abscisses : $x_{1}\approx 0,6$ et $x_{2}\approx 3,4$.
\item $f(x)\leq 3$ : on cherche les abscisses des points de $\mathcal{C}_{f}$ situés sous la droite $y=3$. Celle-ci coupe la parabole en $x\approx -0,8$ et $x\approx 4,8$. L'ensemble des solutions est approximativement $[-0,8\,;\,4,8]$.
\item Le minimum de $f$ sur $[-1\,;\,5]$ est $-1$ et les valeurs aux bornes sont $f(-1)=f(5)=3,5$ (maximal sur l'intervalle). Une droite horizontale $y=m$ coupe donc la courbe :
\begin{itemize}
\item $m=-2$ : $-2<-1$, aucune solution ;
\item $m=-1$ : la droite passe par le sommet, \textbf{une} solution ($x=2$) ;
\item $m=0$ : $-1<0\leq 3,5$, \textbf{deux} solutions ;
\item $m=6$ : $6>3,5$, aucune solution.
\end{itemize}
\end{enumerate}
\textbf{Point de vigilance :} Une lecture graphique ne donne que des valeurs approchées ; on les signale avec le symbole $\approx$ et on ne les présente jamais comme des valeurs exactes.
\end{corrige}

\begin{corrige}[Exercice 9 — Type Probatoire : polynôme du second degré et paramètre]
\textbf{1. a)} Pour tout réel $x$ :
\[ 2(x-2)^{2}-3=2(x^{2}-4x+4)-3=2x^{2}-8x+8-3=2x^{2}-8x+5=f(x). \]
\textbf{b)} $(x-2)^{2}\geq 0$, minimal (nul) en $x=2$. Donc $f$ est \textbf{décroissante sur $[0\,;\,2]$} et \textbf{croissante sur $[2\,;\,4]$}.

\textbf{2. a)} $f'(x)=4x-8=4(x-2)$ : $f'(x)<0$ sur $[0\,;\,2[$, $f'(2)=0$, $f'(x)>0$ sur $]2\,;\,4]$. Cela confirme les variations de la question 1.

\textbf{b)} $f(0)=5$ ; $f(2)=-3$ ; $f(4)=2(4-2)^{2}-3=5$.
\begin{center}
\begin{tabular}{|c|ccccc|}\hline $x$ & $0$ & & $2$ & & $4$ \\ \hline $f'(x)$ & & $-$ & $0$ & $+$ & \\ \hline $f$ & $5$ & $\searrow$ & $-3$ & $\nearrow$ & $5$ \\ \hline\end{tabular}
\end{center}

\textbf{c)} Le minimum de $f$ sur $[0\,;\,4]$ est $\mathbf{-3}$, atteint en $x=2$.

\textbf{3.} $f(1)=2-8+5=-1$ ; $f(3)=18-24+5=-1$.
\begin{center}
\begin{tabular}{|c|ccccc|}\hline
$x$ & $0$ & $1$ & $2$ & $3$ & $4$ \\ \hline
$f(x)$ & $5$ & $-1$ & $-3$ & $-1$ & $5$ \\ \hline
\end{tabular}
\end{center}
On place les cinq points, on marque le sommet $S(2\,;\,-3)$, puis on trace une parabole régulière (symétrique par rapport à la droite $x=2$).

\textbf{4. a)} Graphiquement, la droite $y=1$ coupe $\mathcal{C}_{f}$ en deux points d'abscisses $x_{1}\approx 0,6$ et $x_{2}\approx 3,4$.

\textbf{b)} Par le calcul : $f(x)=1\Leftrightarrow 2(x-2)^{2}-3=1\Leftrightarrow (x-2)^{2}=2\Leftrightarrow x-2=\sqrt{2}$ ou $x-2=-\sqrt{2}$. Les solutions exactes sont $x_{1}=2-\sqrt{2}\approx 0,59$ et $x_{2}=2+\sqrt{2}\approx 3,41$, toutes deux dans $[0\,;\,4]$.

\textbf{5.} D'après le tableau des variations, $f$ prend sur $[0\,;\,4]$ toutes les valeurs de $[-3\,;\,5]$, deux fois chaque valeur de $]-3\,;\,5]$ (une sur chaque branche) :
\begin{center}
\begin{tabularx}{\linewidth}{|l|X|}\hline
Valeurs de $m$ & Nombre de solutions de $f(x)=m$ \\ \hline
$m<-3$ & $0$ \\ \hline
$m=-3$ & $1$ ($x=2$) \\ \hline
$-3<m\leq 5$ & $2$ \\ \hline
$m>5$ & $0$ \\ \hline
\end{tabularx}
\end{center}

\textbf{6.} $f(x)\geq 1$ : la courbe est au-dessus de la droite $y=1$ en dehors des solutions de la question 4 :
\[ \mathcal{S}=[0\,;\,2-\sqrt{2}]\cup[2+\sqrt{2}\,;\,4]. \]

\textbf{Barème indicatif (sur 10) :} 1.a) 1 pt ; 1.b) 1 pt ; 2.a) 1 pt ; 2.b) 1 pt ; 2.c) 0,5 pt ; 3. tableau + courbe 1,5 pt ; 4. 1,5 pt ; 5. 1,5 pt ; 6. 1 pt.

\textbf{Points de vigilance :} Citer la forme canonique pour justifier les variations sans dérivée ; vérifier que les solutions trouvées appartiennent bien à $[0\,;\,4]$ ; dans la discussion en $m$, ne pas oublier les cas $m=-3$ et $m=5$ (bornes incluses).
\end{corrige}

\begin{corrige}[Exercice 10 — Type Probatoire : fonction homographique]
\textbf{1.} $x-2=0\Leftrightarrow x=2$, donc $D_{f}=\mathbb{R}\setminus\{2\}$.

\textbf{2. a)} Pour tout $x\neq 2$ :
\[ 3+\frac{5}{x-2}=\frac{3(x-2)+5}{x-2}=\frac{3x-6+5}{x-2}=\frac{3x-1}{x-2}=f(x). \]
\textbf{b)} $f(x)=u(x-2)+3$ avec $u(x)=\dfrac{5}{x}$ : $\mathcal{C}_{f}$ est l'image de la courbe de $u$ par la translation de vecteur $\vec{t}=2\vec{i}+3\vec{j}$.
\textbf{c)} Le centre de symétrie $O(0\,;\,0)$ de l'hyperbole $y=\frac{5}{x}$ a pour image $\Omega(2\,;\,3)$, centre de symétrie de $\mathcal{C}_{f}$.

\textbf{3.} Asymptote parallèle à l'axe des ordonnées : $x=2$ ; asymptote parallèle à l'axe des abscisses : $y=3$.

\textbf{4. a)} $f(x)=3+5(x-2)^{-1}$, donc $f'(x)=-\dfrac{5}{(x-2)^{2}}$.
\textbf{b)} Pour tout $x\neq 2$, $(x-2)^{2}>0$, donc $f'(x)<0$ : $f$ est strictement décroissante sur $[3\,;\,6]$. $f(3)=\dfrac{8}{1}=8$ ; $f(6)=\dfrac{17}{4}=4,25$.
\begin{center}
\begin{tabular}{|c|ccc|}\hline $x$ & $3$ & & $6$ \\ \hline $f'(x)$ & & $-$ & \\ \hline $f$ & $8$ & $\searrow$ & $4,25$ \\ \hline\end{tabular}
\end{center}

\textbf{5.} $f(4)=\dfrac{11}{2}=5,5$ ; $f(5)=\dfrac{14}{3}\approx 4,67$. On trace les asymptotes $x=2$ et $y=3$ en pointillés, on place les points $(3\,;\,8)$, $(4\,;\,5,5)$, $(5\,;\,4,67)$, $(6\,;\,4,25)$, puis on trace la branche d'hyperbole qui s'approche des asymptotes sans jamais les toucher.

\textbf{6. a)} Graphiquement, sur $[3\,;\,6]$ la courbe est entièrement au-dessus de la droite $y=4$ (même le point le plus bas, $(6\,;\,4,25)$, est au-dessus) : $\mathcal{S}=[3\,;\,6]$.
\textbf{b)} Par le calcul, pour $x\in[3\,;\,6]$ on a $x-2>0$, donc :
\[ f(x)\geq 4\Leftrightarrow \frac{5}{x-2}\geq 1\Leftrightarrow x-2\leq 5\Leftrightarrow x\leq 7. \]
Or $[3\,;\,6]\subset]-\infty\,;\,7]$, donc $\mathcal{S}=[3\,;\,6]$. Les deux méthodes concordent.

\textbf{7.} $g(x)=f(x)-2$ : la courbe de $g$ s'obtient à partir de $\mathcal{C}_{f}$ par la translation de vecteur $-2\vec{j}$. L'asymptote horizontale $y=3$ devient la droite d'équation $y=1$.

\textbf{Barème indicatif (sur 10) :} 1. 0,5 pt ; 2.a) 1 pt ; 2.b) 1 pt ; 2.c) 0,5 pt ; 3. 1 pt ; 4.a) 1 pt ; 4.b) 1 pt ; 5. 1,5 pt ; 6. 1,5 pt ; 7. 1 pt.

\textbf{Points de vigilance :} Exclure la valeur interdite $2$ avant tout calcul ; pour l'inéquation, justifier que $x-2>0$ avant de multiplier les deux membres par $x-2$ (sinon le sens de l'inégalité n'est pas garanti) ; la courbe ne doit jamais couper ses asymptotes.
\end{corrige}

\begin{corrige}[Exercice 11 — Type Probatoire : fonctions associées et modélisation agricole]
\textbf{Partie A.}

\textbf{1.} $f(x)=x^{2}-4x+3=(x-2)^{2}-4+3=(x-2)^{2}-1$. Le sommet est $S(2\,;\,-1)$.

\textbf{2.} $f$ décroît sur $[-1\,;\,2]$ puis croît sur $[2\,;\,5]$ ; $f(-1)=1+4+3=8$ ; $f(2)=-1$ ; $f(5)=25-20+3=8$.
\begin{center}
\begin{tabular}{|c|ccccc|}\hline $x$ & $-1$ & & $2$ & & $5$ \\ \hline $f'(x)$ & & $-$ & $0$ & $+$ & \\ \hline $f$ & $8$ & $\searrow$ & $-1$ & $\nearrow$ & $8$ \\ \hline\end{tabular}
\end{center}
On place le sommet $S(2\,;\,-1)$ et les points $(-1\,;\,8)$, $(0\,;\,3)$, $(1\,;\,0)$, $(3\,;\,0)$, $(4\,;\,3)$, $(5\,;\,8)$, puis on trace la parabole.

\textbf{3.} \begin{itemize}
\item $\mathcal{C}_{g}$ : $g(x)=f(x-2)$, translation de vecteur $2\vec{i}$ appliquée à $\mathcal{C}_{f}$. Image de $S$ : $S_{g}(2+2\,;\,-1)=(4\,;\,-1)$.
\item $\mathcal{C}_{h}$ : $h(x)=-g(x)$, symétrie par rapport à l'axe des abscisses appliquée à $\mathcal{C}_{g}$. Image de $S_{g}$ : $S_{h}(4\,;\,1)$.
\item $\mathcal{C}_{k}$ : $k(x)=|h(x)|$ : on conserve la partie de $\mathcal{C}_{h}$ au-dessus de l'axe des abscisses et on symétrise la partie en dessous. Comme $h(4)=1>0$, le sommet est conservé : $S_{k}(4\,;\,1)$.
\item $\mathcal{C}_{l}$ : $l(x)=k(x)+1$, translation de vecteur $\vec{j}$ appliquée à $\mathcal{C}_{k}$. Image : $S_{l}(4\,;\,2)$.
\end{itemize}

\textbf{4.} Il faut $x-2\in[-1\,;\,5]$, c'est-à-dire $-1\leq x-2\leq 5$, soit $1\leq x\leq 7$. Donc $g$ est définie sur $[1\,;\,7]$.

\textbf{Partie B.}

\textbf{1.} L'aire vaut $400\,m^{2}$ : si $x$ est la longueur de chaque côté perpendiculaire, le côté parallèle à la rivière mesure $\dfrac{400}{x}$ mètres.

\textbf{2.} La longueur de grillage nécessaire est $L(x)=2x+\dfrac{400}{x}$ (deux côtés perpendiculaires et un seul côté parallèle). À $2500$ FCFA le mètre :
\[ C(x)=2500\left(2x+\frac{400}{x}\right)=5000x+\frac{1\,000\,000}{x}. \]

\textbf{3.} $C'(x)=\dfrac{5000(x^{2}-200)}{x^{2}}$. Sur $[5\,;\,80]$, $x^{2}>0$ et $5000>0$, donc $C'(x)$ a le signe de $x^{2}-200$. Or $x^{2}-200=0\Leftrightarrow x=\sqrt{200}=10\sqrt{2}\approx 14,14$ (valeur positive, car $x>0$). Donc $C'(x)<0$ sur $[5\,;\,10\sqrt{2}[$ et $C'(x)>0$ sur $]10\sqrt{2}\,;\,80]$ : $C$ décroît puis croît.
\[ C(5)=25\,000+200\,000=225\,000 \;;\quad C(10\sqrt{2})=50\,000\sqrt{2}+\frac{1\,000\,000}{10\sqrt{2}}=50\,000\sqrt{2}+50\,000\sqrt{2}=100\,000\sqrt{2}\approx 141\,421 \;;\quad C(80)=400\,000+12\,500=412\,500. \]
\begin{center}
\begin{tabular}{|c|ccccc|}\hline $x$ & $5$ & & $10\sqrt{2}$ & & $80$ \\ \hline $C'(x)$ & & $-$ & $0$ & $+$ & \\ \hline $C$ & $225\,000$ & $\searrow$ & $100\,000\sqrt{2}$ & $\nearrow$ & $412\,500$ \\ \hline\end{tabular}
\end{center}

\textbf{4.} Le coût est minimal pour $x=10\sqrt{2}\approx 14,14$ m. Le côté parallèle mesure alors $\dfrac{400}{10\sqrt{2}}=\dfrac{40}{\sqrt{2}}=20\sqrt{2}\approx 28,28$ m. Le coût minimal est $100\,000\sqrt{2}\approx 141\,421$ FCFA, soit environ \textbf{141\,400 FCFA} au cent FCFA près.

\textbf{Barème indicatif (sur 10) :} Partie A : 1. 1 pt ; 2. 1,5 pt ; 3. 2 pts ; 4. 0,5 pt. Partie B : 1. 1 pt ; 2. 1 pt ; 3. 1,5 pt ; 4. 1,5 pt.

\textbf{Points de vigilance :} Pour les transformations successives, chaque transformation s'applique à la courbe \emph{précédente}, pas à $\mathcal{C}_{f}$ ; dans la Partie B, ne pas oublier que la rivière remplace un côté (trois côtés grillagés seulement) ; justifier le signe de $C'$ en précisant que $x^{2}>0$ sur $[5\,;\,80]$ ; conclure par une phrase répondant concrètement à la question posée (dimensions \emph{et} coût).
\end{corrige}

\newpage

\coursec{Fiche bilan}

\begin{popfigure}
\begin{tikzpicture}[mindmap, grow cyclic, every node/.style={concept, font=\small},
  root concept/.append style={concept color=popblue, font=\bfseries},
  level 1/.append style={level distance=3.4cm, sibling angle=60},
  level 2/.append style={level distance=2.2cm, sibling angle=45, font=\scriptsize}]
\node[root concept]{Fonctions associées et étude de fonctions}
  child[concept color=poporange]{ node {Outils préliminaires}
    child { node {Dérivée $f'$} }
    child { node {Forme canonique} }
    child { node {Ensemble de définition} }
  }
  child[concept color=popgreen]{ node {Polynôme de degré 2}
    child { node {Extremum en $x=-\frac{b}{2a}$} }
    child { node {Signe de $a$} }
  }
  child[concept color=poppurple]{ node {Homographique}
    child { node {Asymptote $x=-\frac{d}{c}$} }
    child { node {Forme $\alpha+\frac{k}{x-\beta}$} }
  }
  child[concept color=poppink]{ node {Fonctions associées}
    child { node {$-f$,\ $|f|$} }
    child { node {$f(x-a)+b$} }
  }
  child[concept color=popgold]{ node {Lecture graphique}
    child { node {$f(x)=0$} }
    child { node {$f(x)=m$} }
    child { node {Inéquations} }
  }
  child[concept color=popmuted]{ node {Modélisation}
    child { node {Problèmes concrets} }
  };
\end{tikzpicture}
\legende{Carte mentale de la leçon : les six compétences à maîtriser.}
\end{popfigure}

\begin{bilan}
\textbf{Définitions clés}
\begin{itemize}
  \item \cle{Ensemble de définition} : ensemble des réels $x$ pour lesquels $f(x)$ existe. Pour $f(x)=\dfrac{ax+b}{cx+d}$, $D_f=\mathbb{R}\setminus\left\{-\dfrac{d}{c}\right\}$.
  \item \cle{Asymptote verticale} : droite d'équation $x=-\dfrac{d}{c}$ pour la courbe d'une fonction homographique (parallèle à l'axe des ordonnées).
  \item \cle{Extremum relatif} : valeur $f(\alpha)$ telle que $f$ change de sens de variation en $\alpha$ ; minimum si $a>0$, maximum si $a<0$.
  \item \cle{Fonction associée} : fonction obtenue à partir de $f$ par une transformation géométrique de sa courbe.
\end{itemize}

\textbf{Formules à connaître par cœur}
\begin{center}
\begin{tabularx}{\linewidth}{|l|X|X|}
\hline
\rowcolor{popblueL} \textbf{Objet} & \textbf{Polynôme $ax^2+bx+c$} & \textbf{Homographique $\dfrac{ax+b}{cx+d}$} \\
\hline
Dérivée & $f'(x)=2ax+b$ & $f'(x)=\dfrac{ad-bc}{(cx+d)^2}$ \\
\hline
Forme de référence & $f(x)=a(x-\alpha)^2+\beta$, $\alpha=-\dfrac{b}{2a}$, $\beta=f(\alpha)$ & $f(x)=\alpha+\dfrac{k}{x-\beta}$, $\beta=-\dfrac{d}{c}$, $k=\dfrac{bc-ad}{c^2}$ \\
\hline
Sens de variation & Signe de $f'$ ; extremum en $x=\alpha$ & Signe de $ad-bc$ (constant sur chaque intervalle) \\
\hline
Élément remarquable & Sommet $S(\alpha\,;\,\beta)$ & Asymptotes $x=\beta$ et $y=\alpha$ \\
\hline
Translation & Vecteur $\alpha\vec{i}+\beta\vec{j}$ depuis $y=ax^2$ & Vecteur $\beta\vec{i}+\alpha\vec{j}$ depuis $y=\dfrac{k}{x}$ \\
\hline
\end{tabularx}
\end{center}

\textbf{Transformations des courbes}
\begin{itemize}
  \item $y=-f(x)$ : symétrie de $\mathcal{C}_f$ par rapport à l'axe des abscisses.
  \item $y=|f(x)|$ : on conserve la partie au-dessus de l'axe $(Ox)$, on symétrise la partie en dessous.
  \item $y=f(x-a)$ : translation de vecteur $a\vec{i}$ (vers la droite si $a>0$).
  \item $y=f(x)+b$ : translation de vecteur $b\vec{j}$ (vers le haut si $b>0$).
  \item $y=f(x-a)+b$ : translation de vecteur $a\vec{i}+b\vec{j}$.
\end{itemize}

\textbf{Méthodes express}
\begin{enumerate}
  \item \emph{Étudier $f$} : ensemble de définition $\to$ dérivée $\to$ signe de $f'$ $\to$ tableau de variations $\to$ extremum $\to$ courbe (points + asymptotes).
  \item \emph{Résoudre graphiquement $f(x)=0$} : lire les abscisses des intersections de $\mathcal{C}_f$ avec l'axe $(Ox)$.
  \item \emph{Équation $f(x)=m$} : tracer la droite $y=m$ et compter les intersections ; discuter selon la position de $m$ par rapport aux extrema.
  \item \emph{Inéquation $f(x)\geq m$} : repérer les abscisses des points de $\mathcal{C}_f$ situés au-dessus de la droite $y=m$.
  \item \emph{Mettre sous forme $\alpha+\dfrac{k}{x-\beta}$} : diviser le numérateur par le dénominateur ou identifier.
\end{enumerate}
\end{bilan}

\begin{aretenir}[Checklist avant le Probatoire]
\begin{itemize}
  \item[$\square$] Je sais déterminer l'ensemble de définition d'une fonction polynôme ou homographique.
  \item[$\square$] Je sais calculer la dérivée de $ax^2+bx+c$ et de $\dfrac{ax+b}{cx+d}$ sans erreur de signe.
  \item[$\square$] Je sais écrire la forme canonique $a(x-\alpha)^2+\beta$ et en déduire le sommet.
  \item[$\square$] Je sais donner une équation de l'asymptote verticale d'une fonction homographique.
  \item[$\square$] Je sais dresser un tableau de variations complet (signe de $f'$, flèches, valeurs aux bornes).
  \item[$\square$] Je sais déterminer l'extremum d'un trinôme sur un intervalle borné et vérifier si $\alpha$ appartient à l'intervalle.
  \item[$\square$] Je sais construire la courbe d'un trinôme ou d'une homographique avec ses éléments remarquables.
  \item[$\square$] Je sais construire les courbes de $-f$, $|f|$, $f(x-a)$, $f(x)+b$ par transformation.
  \item[$\square$] Je sais écrire une homographique sous la forme $\alpha+\dfrac{k}{x-\beta}$ et donner le vecteur de translation.
  \item[$\square$] Je sais résoudre graphiquement $f(x)=0$, $f(x)=m$ et des inéquations sur un ensemble borné.
  \item[$\square$] Je sais discuter le nombre de solutions de $f(x)=m$ selon les valeurs du paramètre $m$.
  \item[$\square$] Je rédige proprement : justifications, unités, conclusion en français.
\end{itemize}
\end{aretenir}

\findecours

\end{document}
